\documentclass[11pt,a4paper,fleqn,oneside,a4paper,usenames,dvipsnames,linedheaders]
{book}
\renewcommand{\familydefault}{cmss}

\usepackage[nobottomtitles,pagestyles]{titlesec}%outermarks option?
\renewpagestyle{plain}[\bfseries]{\footrule\setfoot{}{\thepage}{}}
\newpagestyle{mystyle}[\bfseries]{
\headrule
\sethead[\thepage][][F. Della Pietra]{\sectiontitle}{}{\thepage}
}

\usepackage[italian]{babel}                     
\usepackage[T1]{fontenc}                   
\usepackage[utf8]{inputenc}
\usepackage{graphicx}      
\usepackage{yhmath}
\usepackage[font=small]{quoting}            %  
\usepackage{mparhack,relsize}      %
\usepackage{comment}      
\usepackage{caption,multicol}
\usepackage{array}
\usepackage{pgf,tikz}
\usepackage{pgfplots}
\usepackage{polynom}
\usepackage{pifont}
\usepackage[framemethod=TikZ,nobreak=true]{mdframed}
\mdfsetup{skipabove=\topskip,skipbelow=\topskip}

\usetikzlibrary{lindenmayersystems}
\usetikzlibrary[shadings]
\usetikzlibrary{arrows,decorations}
\usetikzlibrary{snakes,patterns}
\usetikzlibrary{decorations.markings}
\makeatletter
\tikzset{% customization of pattern
         % based on <m.wibrow@gm...> - 2013-03-24 07:20: 
        hatch distance/.store in=\hatchdistance,
        hatch distance=5pt,
        hatch thickness/.store in=\hatchthickness,
        hatch thickness=5pt
        }
\pgfdeclarepatternformonly[\hatchdistance,\hatchthickness]{north east hatch}% name
    {\pgfqpoint{-1pt}{-1pt}}% below left
    {\pgfqpoint{\hatchdistance}{\hatchdistance}}% above right
    {\pgfpoint{\hatchdistance-1pt}{\hatchdistance-1pt}}%
    {
        \pgfsetcolor{\tikz@pattern@color}
        \pgfsetlinewidth{\hatchthickness}
        \pgfpathmoveto{\pgfqpoint{0pt}{0pt}}
        \pgfpathlineto{\pgfqpoint{\hatchdistance}{\hatchdistance}}
        \pgfusepath{stroke}
    }
\makeatother		

\usepackage{tcolorbox}
\tcbuselibrary{skins}
\definecolor{oldlace}{rgb}{0.99, 0.96, 0.9}	                    %	Old lace 	                        #FDF5E6 
\tcbset{enhanced,colback=gray!10!white
%,drop shadow
}
\newcommand*\rfrac[2]{{}^{#1}\!/_{#2}}
\usepackage{enumitem,wrapfig}
\setenumerate[1]{label=(\arabic*)}
\setitemize[1]{label=$\textcolor{gray}{\bullet}$}
\definecolor{MyColor1}{rgb}{0.2,0.4,0.6} %mix personal color
\definecolor{lightcoral}{rgb}{0.94, 0.5, 0.5}
\definecolor{lightcornflowerblue}{rgb}{0.6, 0.81, 0.93}
\definecolor{lapislazuli}{rgb}{0.15, 0.38, 0.61}
\definecolor{indigo}{rgb}{0.0, 0.25, 0.42}
\definecolor{internationalkleinblue}{rgb}{0.0, 0.18, 0.65}
\definecolor{ao(english)}{rgb}{0.0, 0.5, 0.0}	               
\definecolor{chiaro}{RGB}{111,180,180}
\definecolor{arancia}{RGB}{253,134,40}
\definecolor{cadmiumgreen}{rgb}{0.0, 0.42, 0.24}	 
\newcommand{\eqre}[1]{\textcolor{internationalkleinblue}{{(\ref{#1})}}}   
\newcommand{\boxx}[1]{
\begin{tcolorbox}[opacityframe=0.75,
colframe=lightcornflowerblue%,drop fuzzy shadow=lightcornflowerblue!70!white
] 
#1
\end{tcolorbox}
}
\newcommand{\boxb}[1]{
\begin{tcolorbox}[opacityframe=0.75,
colframe=green%,drop fuzzy shadow=lightcornflowerblue!70!white
] 
#1
\end{tcolorbox}
}
\newcommand{\boxr}[1]{
\begin{tcolorbox}[opacityframe=0.75,
colframe=lightcoral%,drop fuzzy shadow=lightcoral!50!white
] 
#1
\end{tcolorbox}
}
\newcommand{\boxa}[1]{
\begin{tcolorbox}[opacityframe=0.75,
colframe=ao(english)%,drop fuzzy shadow=lightcornflowerblue!70!white
] 
#1
\end{tcolorbox}
}
\pgfplotsset{compat=1.16}
\newcommand{\lp}{\left(}
\newcommand{\rp}{\right)}

%\newcommand{\boxx}[1]{
%\begin{tcolorbox}[opacityframe=0.75,colback=white,
%colframe=blue] 
%#1
%\end{tcolorbox}
%}
%\newcommand{\boxr}[1]{
%\begin{tcolorbox}[opacityframe=0.75,colback=white,
%colframe=red] 
%#1
%\end{tcolorbox}
%}
\usepackage{xcolor}
       
\usepackage[colorlinks=true,linkcolor=internationalkleinblue]{hyperref}
 %\usepackage[top=2.5cm, bottom=2.5cm, left=2.5cm, right=2.5cm]{geometry}

         % 
%\usepackage[nochapters,beramono,eulermath,%                 
%           pdfspacing,%                    %
%           listings,%                      % 
%            ]{classicthesis}                % 
%\usepackage{arsclassica}     

%\usepackage{framed}
\usepackage[hyperpageref]{backref}
               
%\usepackage{geometry}
%\input impostazioni-articolo.tex 
%\usepackage{showkeys}
%\usepackage{cite}
%\usepackage{refcheck}
\usepackage{picture}
\usepackage{amsmath,amsthm,amsfonts,amssymb}
\usepackage{comment}
\usepackage{titling} 



			%%%%%%%%%%%%%%%%%%%
			%%			 	 %%
			%%    Comandi    %%
			%%			 	 %%
			%%%%%%%%%%%%%%%%%%%

\newcommand{\diesis}{^\#}

\theoremstyle{definition}
\newcounter{definiz}[chapter]
\setcounter{definiz}{0}
\renewcommand{\thedefiniz}{\arabic{chapter}.\arabic{definiz}}
\newenvironment{definiz}[2][]{%
\refstepcounter{definiz}%
\ifstrempty{#1}%
{\mdfsetup{%
backgroundcolor=cadmiumgreen!5,%
roundcorner=10pt,
innertopmargin=5pt,
innerbottommargin=10pt,linecolor=cadmiumgreen!50,%
linewidth=1.5pt,topline=true,
frametitle={%
\tikz[baseline=(current bounding box.east),outer sep=0pt]
\node[anchor=east,rectangle,line width=1pt,draw=cadmiumgreen,fill=cadmiumgreen!10,rounded corners]
{\strut Definizione~\thedefiniz};}}
}%
{\mdfsetup{%
backgroundcolor=cadmiumgreen!5,%
roundcorner=10pt,
innertopmargin=5pt,
innerbottommargin=10pt,linecolor=cadmiumgreen!50,%
linewidth=1.5pt,topline=true,
frametitle={%
\tikz[baseline=(current bounding box.east),outer sep=0pt]
\node[anchor=east,rectangle,line width=1pt,draw=cadmiumgreen,fill=cadmiumgreen!10,rounded corners] 
{\strut Definizione~\thedefiniz:~#1};}}%
}%
\mdfsetup{
backgroundcolor=cadmiumgreen!5,%
roundcorner=10pt,
innertopmargin=5pt,
innerbottommargin=10pt,linecolor=cadmiumgreen!50,%
linewidth=1.5pt,topline=true,
frametitleaboveskip=\dimexpr-\ht\strutbox\relax,}
\begin{mdframed}[]\relax%
}{\end{mdframed}}

%%%%Assioma
\newcounter{assioma}[chapter]
\setcounter{assioma}{0}
\renewcommand{\theassioma}{\arabic{chapter}.\arabic{assioma}}
\newenvironment{assioma}[2][]{%
\refstepcounter{assioma}%
\ifstrempty{#1}%
{\mdfsetup{%
backgroundcolor=cadmiumgreen!5,%
roundcorner=10pt,
innertopmargin=5pt,
innerbottommargin=10pt,linecolor=cadmiumgreen!50,%
linewidth=1.5pt,topline=true,
frametitle={%
\tikz[baseline=(current bounding box.east),outer sep=0pt]
\node[anchor=east,rectangle,line width=1pt,draw=cadmiumgreen,fill=cadmiumgreen!10,rounded corners]
{\strut Assioma~\theassioma};}}
}%
{\mdfsetup{%
backgroundcolor=cadmiumgreen!5,%
roundcorner=10pt,
innertopmargin=5pt,
innerbottommargin=10pt,linecolor=cadmiumgreen!50,%
linewidth=1.5pt,topline=true,
frametitle={%
\tikz[baseline=(current bounding box.east),outer sep=0pt]
\node[anchor=east,rectangle,line width=1pt,draw=cadmiumgreen,fill=cadmiumgreen!10,rounded corners] 
{\strut Assioma~\theassioma:~#1};}}%
}%
\mdfsetup{
backgroundcolor=cadmiumgreen!5,%
roundcorner=10pt,
innertopmargin=5pt,
innerbottommargin=10pt,linecolor=cadmiumgreen!50,%
linewidth=1.5pt,topline=true,
frametitleaboveskip=\dimexpr-\ht\strutbox\relax,}
\begin{mdframed}[]\relax%
}{\end{mdframed}}

\newcounter{vuoto}[chapter]
\setcounter{vuoto}{0}
\renewcommand{\thevuoto}{\arabic{chapter}.\arabic{vuoto}}
\newenvironment{vuoto}[2][]{%
\refstepcounter{vuoto}%
\ifstrempty{#1}%
{\mdfsetup{%
backgroundcolor=cadmiumgreen!5,%
roundcorner=10pt,
innertopmargin=5pt,
innerbottommargin=10pt,linecolor=cadmiumgreen!50,%
linewidth=1.5pt,topline=true,
frametitle={%
\tikz[baseline=(current bounding box.east),outer sep=0pt]
\node[anchor=east,rectangle,line width=1pt,draw=cadmiumgreen,fill=cadmiumgreen!10,rounded corners]
{\strut };}}
}%
{\mdfsetup{%
backgroundcolor=cadmiumgreen!5,%
roundcorner=10pt,
innertopmargin=5pt,
innerbottommargin=10pt,linecolor=cadmiumgreen!50,%
linewidth=1.5pt,topline=true,
frametitle={%
\tikz[baseline=(current bounding box.east),outer sep=0pt]
\node[anchor=east,rectangle,line width=1pt,draw=cadmiumgreen,fill=cadmiumgreen!10,rounded corners] 
{\strut ~#1};}}%
}%
\mdfsetup{
backgroundcolor=cadmiumgreen!5,%
roundcorner=10pt,
innertopmargin=5pt,
innerbottommargin=10pt,linecolor=cadmiumgreen!50,%
linewidth=1.5pt,topline=true,
frametitleaboveskip=\dimexpr-\ht\strutbox\relax,}
\begin{mdframed}[nobreak=false]\relax%
}{\end{mdframed}}


\newtheorem*{claim}{Claim}
%\newtheorem{rem}[definiz]{Osservazione}

\theoremstyle{definition}
\newtheorem{ex}{Esempio}
%\newcounter{ex}[chapter]
\setcounter{ex}{0}
\renewcommand{\theex}{\thechapter.\arabic{ex}}
\newtheorem{esercizio}{Esercizio}
%\newcounter{ex}[chapter]
\setcounter{esercizio}{0}
\renewcommand{\theesercizio}{\thechapter.\arabic{esercizio}}

\newtheorem*{notazione}{Notazione}
\newtheorem*{sol}{Soluzione}
\theoremstyle{plain}
%\newtheorem{theorem}[definiz]{Teorema}
\newtheorem*{princ}{Principio di induzione matematica}
\newtheorem{lemma}[definiz]{Lemma}
\newtheorem{cor}[definiz]{Corollario}

\newcounter{prop}[chapter]%\setcounter{prop}{0}
\renewcommand{\theprop}{\arabic{chapter}.\arabic{prop}}
\newenvironment{prop}[2][]{%
\refstepcounter{prop}%
\ifstrempty{#1}%
{\mdfsetup{%
backgroundcolor=arancia!7,%
roundcorner=10pt,
innertopmargin=5pt,
innerbottommargin=10pt,
linecolor=arancia,%
linewidth=1.5pt,topline=true,
frametitle={%
\tikz[baseline=(current bounding box.east),outer sep=0pt]
\node[anchor=east,rectangle,line width=1pt,draw=arancia,fill=arancia!15,rounded corners]
{\strut Proposizione~\theprop};}}
}%
{\mdfsetup{%
backgroundcolor=arancia!7,%
roundcorner=10pt,
innertopmargin=5pt,
innerbottommargin=10pt,
linecolor=arancia,%
linewidth=1.5pt,topline=true,
frametitle={%
\tikz[baseline=(current bounding box.east),outer sep=0pt]
\node[draw=arancia,line width=1pt,anchor=east,line width=1pt,rectangle,draw=arancia,fill=arancia!15,rounded corners] 
{\strut Proposizione~\theprop:~#1};}}%
}%
\mdfsetup{
backgroundcolor=arancia!7,%
roundcorner=10pt,
innertopmargin=5pt,
innerbottommargin=10pt,
linecolor=arancia,%
linewidth=1.5pt,topline=true,
frametitleaboveskip=\dimexpr-\ht\strutbox\relax,}
\begin{mdframed}[]\relax%
}{\end{mdframed}}

\newcounter{propvuota}[chapter]%\setcounter{prop}{0}
\renewcommand{\thepropvuota}{\arabic{chapter}.\arabic{propvuota}}
\newenvironment{propvuota}[2][]{%
\refstepcounter{propvuota}%
\ifstrempty{#1}%
{\mdfsetup{%
backgroundcolor=arancia!7,%
roundcorner=10pt,
innertopmargin=5pt,
innerbottommargin=10pt,
linecolor=arancia,%
linewidth=1.5pt,topline=true,
frametitle={%
\tikz[baseline=(current bounding box.east),outer sep=0pt]
\node[anchor=east,rectangle,line width=1pt,draw=arancia,fill=arancia!15,rounded corners]
{\strut ~#1};}}
}%
{\mdfsetup{%
backgroundcolor=arancia!7,%
roundcorner=10pt,
innertopmargin=5pt,
innerbottommargin=10pt,
linecolor=arancia,%
linewidth=1.5pt,topline=true,
frametitle={%
\tikz[baseline=(current bounding box.east),outer sep=0pt]
\node[draw=arancia,line width=1pt,anchor=east,line width=1pt,rectangle,draw=arancia,fill=arancia!15,rounded corners] 
{\strut ~#1};}}%
}%
\mdfsetup{
backgroundcolor=arancia!7,%
roundcorner=10pt,
innertopmargin=5pt,
innerbottommargin=10pt,
linecolor=arancia,%
linewidth=1.5pt,topline=true,
frametitleaboveskip=\dimexpr-\ht\strutbox\relax,}
\begin{mdframed}[]\relax%
}{\end{mdframed}}


\newcounter{theorem}[chapter]%\setcounter{prop}{0}
\renewcommand{\thetheorem}{\arabic{chapter}.\arabic{theorem}}
\newenvironment{theorem}[2][]{%
\refstepcounter{theorem}%
\ifstrempty{#1}%
{\mdfsetup{%
backgroundcolor=arancia!5,%
roundcorner=10pt,
innertopmargin=5pt,
innerbottommargin=10pt,linecolor=chiaro,%
linewidth=1.5pt,topline=true,
frametitle={%
\tikz[baseline=(current bounding box.east),outer sep=0pt]
\node[anchor=east,rectangle,draw=arancia,fill=arancia!15,rounded corners]
{\strut Teorema~\thetheorem};}}
}%
{\mdfsetup{%
backgroundcolor=arancia!5,%
roundcorner=10pt,
innertopmargin=5pt,
innerbottommargin=10pt,linecolor=chiaro,%
linewidth=1.5pt,topline=true,
frametitle={%
\tikz[baseline=(current bounding box.east),outer sep=0pt]
\node[anchor=east,rectangle,draw=arancia,fill=arancia!20,rounded corners] 
{\strut Teorema~\thetheorem:~#1};}}%
}%
\mdfsetup{
backgroundcolor=arancia!5,%
roundcorner=10pt,
innertopmargin=5pt,
innerbottommargin=10pt,linecolor=arancia!20,%
linewidth=1.5pt,topline=true,
frametitleaboveskip=\dimexpr-\ht\strutbox\relax,}
\begin{mdframed}[]\relax%
}{\end{mdframed}}





%%% OSSERVAZIONE

\newcounter{rem}[chapter]
\setcounter{rem}{0}
\renewcommand{\therem}{\arabic{chapter}.\arabic{rem}}
\newenvironment{rem}[2][]{%
\refstepcounter{rem}%
\ifstrempty{#1}%
{\mdfsetup{%
backgroundcolor=chiaro!5,%
roundcorner=10pt,
innertopmargin=5pt,
innerbottommargin=10pt,linecolor=chiaro,%
linewidth=1.5pt,topline=true,
frametitle={%
\tikz[baseline=(current bounding box.east),outer sep=0pt]
\node[line width=1pt,draw=chiaro,anchor=east,rectangle,fill=chiaro!20,rounded corners]
{\strut Osservazione~\therem};}}
}%
{\mdfsetup{%
backgroundcolor=chiaro!5,%
roundcorner=10pt,
innertopmargin=5pt,
innerbottommargin=10pt,linecolor=chiaro,%
linewidth=1.5pt,topline=true,
frametitle={%
\tikz[baseline=(current bounding box.east),outer sep=0pt]
\node[line width=1pt,draw=chiaro,anchor=east,rectangle,fill=chiaro!20,rounded corners]
{\strut Osservazione~\therem:~#1};}}%
}%
\mdfsetup{
backgroundcolor=chiaro!5,%
roundcorner=10pt,
innertopmargin=5pt,
innerbottommargin=10pt,linecolor=chiaro,%
linewidth=1.5pt,topline=true,
frametitleaboveskip=\dimexpr-\ht\strutbox\relax,}
\begin{mdframed}[nobreak=false]\relax%
\label{#2}
}{\end{mdframed}}
\numberwithin{equation}{chapter}

\newcommand{\norm}[1]{\left\Vert#1\right\Vert}
\newcommand{\ds}{\displaystyle}
\newcommand{\ol}[1]{\overline{#1}}
\newcommand{\ul}[1]{\underline{#1}}
\newcommand{\abs}[1]{\left\vert#1\right\vert}
\newcommand{\set}[1]{\left\{#1\right\}}
\newcommand{\R}{\mathbb R}
\newcommand{\Q}{\mathbb Q}
\newcommand{\N}{\mathbb N}
\newcommand{\Z}{\mathbb Z}
\newcommand{\Rn}{\mathbb R^n}
\newcommand{\W}{\mathcal W}
\newcommand{\de}{\partial}
\newcommand{\eps}{\varepsilon}
\newcommand{\To}{\longrightarrow}
\newcommand{\wS}{\overline w_{\eps,\delta}}
\newcommand{\ws}{\underline w_{\eps,\delta}}
\newcommand{\lto}{\left(}
\newcommand{\rto}{\right)}
\newcommand{\lqu}{\left[}
\newcommand{\rqu}{\right]}
\newcommand{\lgr}{\left\{}
\newcommand{\rgr}{\right\}}
\newcommand\restr[2]{{% we make the whole thing an ordinary symbol
  \left.\kern-\nulldelimiterspace % automatically resize the bar with \right
  #1 % the function
  \vphantom{\big|} % pretend it's a little taller at normal size
  \right|_{#2} % this is the delimiter
  }}
\newcommand{\bul}{$\textcolor{gray}{\bullet}$}
\newcommand{\enclose}[1]{{\left( #1 \right)}}
\newcommand{\Enclose}[1]{{\left[ #1 \right]}}
\newcommand{\ENCLOSE}[1]{{\left\{ #1 \right\}}}
\DeclareMathOperator{\divergenza}{div}
\DeclareMathOperator{\epi}{epi}
\DeclareMathOperator{\ess}{ess}
\DeclareMathOperator{\esssup}{ess\,sup}
\DeclareMathOperator{\dist}{dist}
\DeclareMathOperator{\D}{D}
\DeclareMathOperator{\sign}{sign}
\DeclareMathOperator{\Sk}{S}
\DeclareMathOperator{\Imm}{Im}
\DeclareMathOperator{\osc}{osc}

\newcounter{es}
\setcounter{es}{1}
\newcommand{\esn}{\thees\stepcounter{es}}

\usepackage{textgreek}
\makeatletter
\newcommand{\romangreek}[1]{\c@romangreek{#1}}
\newcommand{\c@romangreek}[1]{%
  {%
    \ifcase\number\value{#1} %
    \or a\or b\or c\or d\or e\or f\or g\or h\or i\or j\or k\or l\or m\or n\or o\or p\or q\or r\or s\or t\or u\or v\or w\or x\or y\or z\or \textalpha \or \textbeta \or \textgamma \or \textdelta \or \textepsilon \or \textzeta \or \texteta \or \texttheta  %\or \straighttheta
    \or    \textiota \or \textkappa \or \textlambda \or \textmu \or \textnu    \or \textxi \or \textomikron \or \textrho \or \textpi \or \textsigma \or \texttau \or \textupsilon \or \textphi
 \or   \textchi \or  \textpsi \or \textomega
 \or A\or B\or C\or D\or E\or F\or G\or H\or I\or J\or K\or M\or N\or O\or P\or Q\or R\or S\or T\or \or U \or V \or W \or X \or Y \or Z
  \textAlpha \or \textBeta \or \textGamma \or \textDelta \or \textEpsilon \or \textZeta
\or \textEta \or \textTheta \or \textIota \or \textKappa \or \textLambda\or\textMu\or\textNu\or\textXi\or\textOmikron\or\textRho\or
\textPi\or\textSigma\or\textTau\or\textUpsilon\or\textPhi\or\textChi\or\textPsi\or\textOmega\fi
  }%
}


\AddEnumerateCounter*{\romangreek}{\c@romangreek}{5}


%\newtheoremstyle{note}% name
% {3pt}% measure of space to leave above the theorem. E.g.: 3pt
%   {3pt}% measure of space to leave below the theorem. E.g.: 3pt
%   {\ttfamily}% name of font to use in the body of the theorem
%   {1pt}% measure of space to indent
%   {\bfseries}% name of head font
%   {\textbf{.}}% punctuation between head and body
%   {10mm}% space after theorem head
%   {}% Manually specify head
%
%\theoremstyle{note}
%\newtheorem{note}{Osservazione}[chapter]


%\includeonly{alfabetogreco}
%\includeonly{reali}
%\includeonly{successioni}
%\includeonly{serie}
%\includeonly{continue}
%\includeonly{calcolodiff}
%\includeonly{integrali}
%\includeonly{biblio}




%%%%%%%%%%%%%%%%%%%%%%%%%%%%%%%%%%%%%%%%%%%%%%%%%%%%%%%%%%%%%%%%%%%%%
%% Impostazioni specifiche della raccolta di esercizi
\makeatletter\@addtoreset{esercizio}{chapter}\makeatother   % numerazione per scheda
\addto\captionsitalian{\renewcommand{\chaptername}{Scheda}}
\newif\ifsoluzioni
\soluzionitrue      % \soluzionifalse per la versione senza soluzioni
% --- pagina ---
\usepackage{geometry}
\geometry{a4paper,top=2.3cm,bottom=2.3cm,left=2.5cm,right=2.5cm}
% --- soluzioni raccolte in fondo a ogni scheda ---
% \risp{...} dopo ogni esercizio mette da parte la soluzione (con il numero
% dell'esercizio); \stampasoluzioni, alla fine di ogni scheda, le stampa
% tutte a partire da una pagina nuova. Con \soluzionifalse non stampa nulla.
\makeatletter
\newcommand{\solbuf}{}
\newcommand{\risp}[1]{\ifsoluzioni\edef\sol@num{\theesercizio}%
  \expandafter\sol@append\expandafter{\sol@num}{#1}\fi}
\newcommand{\sol@append}[2]{\g@addto@macro\solbuf{\solitem{#1}{#2}}}
\newcommand{\solitem}[2]{\par\noindent\textbf{\textcolor{internationalkleinblue}{#1.}}\ #2\par\medskip}
\newpagestyle{solstyle}[\bfseries]{\headrule\sethead{Soluzioni --- Scheda \thechapter}{}{\thepage}}
\newcommand{\stampasoluzioni}{\ifsoluzioni\clearpage\pagestyle{solstyle}
  \section*{Soluzioni degli esercizi della Scheda~\thechapter}
  \addcontentsline{toc}{section}{Soluzioni}
  {\small\solbuf}\gdef\solbuf{}\clearpage\pagestyle{mystyle}\fi}
\makeatother
\newcommand{\liv}[1]{\ifcase#1\or\ding{72}\or\ding{72}\ding{72}\or\ding{72}\ding{72}\ding{72}\fi}
%%%%%%%%%%%%%%%%%%%%%%%%%%%%%%%%%%%%%%%%%%%%%%%%%%%%%%%%%%%%%%%%%%%%%
\begin{document}
\pagestyle{mystyle}
\title{\large\textsc{Istituzioni di Matematica --- Esercizi} \\[.2cm] \normalsize Corso di Laurea in Scienze per la Natura e per l'Ambiente \\[.3cm] \textit{Schede settimanali, aggiornate progressivamente}}
\author{\normalsize \textsc{F. Della Pietra}}
\date{A.A. 2026--2027}

\maketitle

\tableofcontents

% !TEX root =esercizi.tex
\chapter{Lezioni 1--3 (21--28 settembre 2026)}

\begin{vuoto}[Argomenti svolti]{}
\begin{itemize}
\item[\bul] \textbf{Lezione 1 --- Insiemi}. Insiemi, appartenenza, inclusione, insieme vuoto, insieme delle parti; unione, intersezione, differenza, complementare. Quantificatori $\forall$ ed $\exists$. (Le leggi di De Morgan e la negazione dei quantificatori sono nell'Approfondimento~1.1 degli appunti: gli esercizi 1.10 e 1.12 qui sotto vanno con quello.) Prodotto cartesiano, partizioni. Cardinalit\`a: principio di addizione, principio di inclusione-esclusione (richiamato all'inizio della sezione 1.3), $|A\times B|=|A|\,|B|$.
\item[\bul] \textbf{Lezione 2 --- Induzione e calcolo combinatorio}. Principio di induzione (forma insiemistica e dimostrazioni per induzione). Principio di moltiplicazione. Disposizioni con ripetizione, $|\mathcal P(A)|=2^{n}$. Fattoriale, disposizioni semplici, permutazioni. Combinazioni e coefficienti binomiali, triangolo di Tartaglia, binomio di Newton.
\item[\bul] \textbf{Lezione 3 --- I numeri reali}. Naturali, interi, razionali; sviluppo decimale dei razionali; irrazionalit\`a di $\sqrt2$. I numeri reali come allineamenti decimali; troncamenti; ordine e operazioni. Retta reale, intervalli, valore assoluto e distanza, disuguaglianza triangolare. (Estremo superiore e completezza, \S3.8, sono nella Scheda~2.)
\end{itemize}
\end{vuoto}

Il livello di difficolt\`a \`e indicato con \liv{1} (esercizio di base), \liv{2} (esercizio standard), \liv{3} (esercizio pi\`u impegnativo).

%-------------------------------------------------------------------
\section{Insiemi e operazioni}

\begin{esercizio}[\liv{1}]
Siano $U=\{1,2,\ldots,12\}$, $A$ l'insieme dei multipli di $2$, $B$ quello dei multipli di $3$, $C$ quello dei divisori di $12$ (tutti in $U$). Scrivere per elencazione
\begin{gather*}
A\cap B\cap C,\qquad (A\cup B)\setminus C,\\
C^{c},\qquad A\cap(B\cup C).
\end{gather*}
\end{esercizio}
\risp{$A=\{2,4,6,8,10,12\}$, $B=\{3,6,9,12\}$, $C=\{1,2,3,4,6,12\}$. Si ha $A\cap B\cap C=\{6,12\}$, $(A\cup B)\setminus C=\{8,9,10\}$, $C^{c}=\{5,7,8,9,10,11\}$, $A\cap(B\cup C)=\{2,4,6,12\}$.}

\begin{esercizio}[\liv{1}]
Descrivere mediante una propriet\`a caratteristica gli insiemi
\[
\{1,4,9,16,25\},\qquad\{\ldots,-6,-3,0,3,6,\ldots\},\qquad\Bigl\{\frac12,\frac14,\frac18,\frac1{16}\Bigr\}.
\]
\end{esercizio}
\risp{$\{n^{2}\colon n\in\N,\ n\le5\}$; \ $\{3k\colon k\in\Z\}$; \ $\{2^{-n}\colon n\in\N,\ n\le4\}$.}

\begin{esercizio}[\liv{1}]
Scrivere per elencazione gli insiemi
\[
\{x\in\N\colon x^{2}-5x+6=0\},\quad\{x\in\Z\colon |x|<3\},\quad\{x\in\N\colon 3x+1<11\}.
\]
\end{esercizio}
\risp{$\{2,3\}$; \ $\{-2,-1,0,1,2\}$; \ $\{1,2,3\}$.}

\begin{esercizio}[\liv{1}]
Stabilire quali delle seguenti affermazioni sono vere:
\begin{gather*}
\emptyset\in\emptyset,\qquad\emptyset\subseteq\emptyset,\qquad\emptyset\in\{\emptyset\},\\
\{1\}\subseteq\{\{1\},2\},\qquad\{1\}\in\{\{1\},2\},\qquad\{1,2\}\subseteq\{1,2,\{1,2\}\}.
\end{gather*}
\end{esercizio}
\risp{Sono vere la seconda, la terza, la quinta e la sesta.}

\begin{esercizio}[\liv{2}]
Scrivere l'insieme delle parti di $\{1,\{2\}\}$.
\end{esercizio}
\risp{$\bigl\{\emptyset,\ \{1\},\ \{\{2\}\},\ \{1,\{2\}\}\bigr\}$.}

\begin{esercizio}[\liv{2}]
Provare le \textit{leggi di assorbimento}: $A\cup(A\cap B)=A$ e $A\cap(A\cup B)=A$.
\end{esercizio}
\risp{Poiché $A\cap B\subseteq A$, unire $A\cap B$ ad $A$ non aggiunge elementi: $A\cup(A\cap B)=A$. Poiché $A\subseteq A\cup B$, gli elementi comuni ad $A$ e ad $A\cup B$ sono esattamente quelli di $A$: $A\cap(A\cup B)=A$.}

\begin{esercizio}[\liv{2}]
Dimostrare o confutare con un controesempio:
\begin{enumerate}
\item se $A\setminus B=B\setminus A$, allora $A=B$;
\item se $A\cup B=A\cup C$, allora $B=C$;
\item $A\cap B=A$ se e solo se $A\subseteq B$.
\end{enumerate}
\end{esercizio}
\risp{(1) Vera: se esistesse $x\in A\setminus B$, sarebbe $x\in B\setminus A$, e quindi $x\in B$, assurdo; dunque $A\setminus B=\emptyset$, cio\`e $A\subseteq B$, e analogamente $B\subseteq A$. (2) Falsa: $A=\{1,2\}$, $B=\{1\}$, $C=\{2\}$. (3) Vera.}

\begin{esercizio}[\liv{1}]
Sia $U$ l'insieme delle specie presenti in una riserva, $M$ l'insieme dei mammiferi, $N$ quello delle specie notturne, $P$ quello delle specie protette. Esprimere mediante operazioni tra insiemi: (a) i mammiferi notturni non protetti; (b) le specie che non sono n\'e mammiferi n\'e protette; (c) le specie protette che non sono mammiferi notturni.
\end{esercizio}
\risp{(a) $M\cap N\cap P^{c}$; (b) $M^{c}\cap P^{c}$; (c) $P\setminus(M\cap N)$.}

\begin{esercizio}[\liv{2}]
Si definisce \textit{differenza simmetrica} di $A$ e $B$ l'insieme $A\,\Delta\,B=(A\setminus B)\cup(B\setminus A)$. Provare che
\[
A\,\Delta\,B=(A\cup B)\setminus(A\cap B).
\]
Calcolare $N\,\Delta\,S$ per i due stagni dell'Esempio 1.1 degli appunti, e interpretare il risultato.
\end{esercizio}
\risp{$N\,\Delta\,S=\{g,c,p,a,s\}$: le specie osservate in uno solo dei due stagni.}

%-------------------------------------------------------------------
\section{Quantificatori}

\begin{esercizio}[\liv{1}]
Scrivere con i quantificatori, e poi negare senza usare ``non è vero che'':
(a) ``ogni campione ha pH minore di $8$''; (b) ``esiste una specie presente in tutti i siti''; (c) ``in ogni sito c'è almeno una specie protetta''.
\end{esercizio}
\risp{(a) $\forall$ campione $c$: $\mathrm{pH}(c)<8$; negazione: esiste un campione con pH maggiore o uguale a $8$. (b) $\exists$ specie $s$ $\forall$ sito $x$: $s$ è presente in $x$; negazione: per ogni specie esiste un sito in cui non è presente. (c) $\forall$ sito $x$ $\exists$ specie protetta $s$ presente in $x$; negazione: esiste un sito in cui non c'è alcuna specie protetta.}

\begin{esercizio}[\liv{1}]
Stabilire se sono vere o false:
(a) $\forall n\in\N:\ n^2\ge n$; \ (b) $\exists n\in\N:\ n^2=2n$; \ (c) $\forall n\in\N\ \exists m\in\N:\ m>n$; \ (d) $\exists m\in\N\ \forall n\in\N:\ m>n$.
\end{esercizio}
\risp{(a) vera; (b) vera ($n=2$); (c) vera (basta $m=n+1$); (d) falsa: nessun naturale è maggiore di tutti i naturali (non lo è di sé stesso). La (c) e la (d) differiscono solo per l'ordine dei quantificatori.}

\begin{esercizio}[\liv{2}]
Sia $A\subseteq\R$. Scrivere con i quantificatori la frase ``esiste un numero reale maggiore o uguale a tutti gli elementi di $A$'' (cioè: $A$ è limitato superiormente), e scriverne la negazione.
\end{esercizio}
\risp{$\exists\,M\in\R\ \forall x\in A:\ x\le M$. Negazione: $\forall M\in\R\ \exists\,x\in A:\ x>M$ (per quanto grande si scelga $M$, c'è un elemento di $A$ che lo supera).}

%-------------------------------------------------------------------
\section{Prodotto cartesiano e cardinalit\`a}

\begin{vuoto}[Il principio di inclusione-esclusione]{}
Se $A$ e $B$ sono insiemi finiti \emph{disgiunti}, $|A\cup B|=|A|+|B|$: è il principio di addizione. Se invece hanno elementi in comune, sommando $|A|$ e $|B|$ gli elementi di $A\cap B$ vengono contati \emph{due volte}, una in $A$ e una in $B$; per correggere si sottraggono una volta:
\[
|A\cup B|=|A|+|B|-|A\cap B| .
\]
Esempio: se in $10$ stagni c'è la rana verde, in $7$ il tritone e in $4$ entrambi, gli stagni con almeno una delle due specie sono $10+7-4=13$. Con tre insiemi si ragiona allo stesso modo: si sommano le cardinalità, si tolgono quelle delle intersezioni a due a due (contate due volte), e si rimette quella dell'intersezione dei tre, che a quel punto era stata aggiunta tre volte e tolta tre volte:
\[
|A\cup B\cup C|=|A|+|B|+|C|-|A\cap B|-|A\cap C|-|B\cap C|+|A\cap B\cap C| .
\]
\end{vuoto}

\begin{esercizio}[\liv{1}]
Siano $|A|=3$ e $|B|=4$. Quanti elementi hanno $A\times B$, $\mathcal P(A)$, $\mathcal P(A\times B)$?
\end{esercizio}
\risp{$12$, $8$, $2^{12}=4096$.}

\begin{esercizio}[\liv{1}]
Elencare gli elementi di $\{0,1\}\times\{a,b,c\}$ e di $\{a,b,c\}\times\{0,1\}$. I due insiemi sono uguali?
\end{esercizio}
\risp{$\{(0,a),(0,b),(0,c),(1,a),(1,b),(1,c)\}$ e $\{(a,0),(a,1),(b,0),(b,1),(c,0),(c,1)\}$. Non sono uguali: le coppie sono ordinate, e $(0,a)\ne(a,0)$. Hanno però lo stesso numero di elementi, $6$.}

\begin{esercizio}[\liv{1}]
Un piano di monitoraggio prevede $4$ siti ($s_1,\dots,s_4$) e $3$ campagne di prelievo (primavera, estate, autunno). Quanti prelievi si fanno in tutto? Elencare quelli relativi al sito $s_1$.
\end{esercizio}
\risp{$|S\times C|=4\cdot3=12$. Per il sito $s_1$: $(s_1,\text{primavera})$, $(s_1,\text{estate})$, $(s_1,\text{autunno})$.}

\begin{esercizio}[\liv{1}]
Siano $A=\{\text{rosso},\text{verde}\}$ e $B=\{1,2,3\}$. Quante sono le coppie $(a,b)$ con $a\in A$ e $b\in B$? Quante di esse hanno $b$ dispari?
\end{esercizio}
\risp{$|A\times B|=6$; con $b$ dispari: $2\cdot2=4$.}

\begin{esercizio}[\liv{2}]
In un'indagine su $80$ siti, la specie $X$ \`e presente in $50$ siti, la specie $Y$ in $35$, e $10$ siti non ospitano nessuna delle due. In quanti siti sono presenti entrambe? In quanti soltanto $X$?
\end{esercizio}
\risp{$|X\cup Y|=70$, quindi $|X\cap Y|=50+35-70=15$; soltanto $X$: $35$; soltanto $Y$: $20$.}

\begin{esercizio}[\liv{2}]
Su $60$ campioni di suolo, $25$ superano la soglia per i nitrati ($N$), $20$ quella per i fosfati ($F$), $15$ quella per i metalli pesanti ($M$). Inoltre $|N\cap F|=8$, $|N\cap M|=5$, $|F\cap M|=4$, e $2$ campioni superano tutte e tre le soglie. Quanti campioni superano almeno una soglia? Quanti nessuna? Quanti esattamente una?
\end{esercizio}
\risp{Almeno una: $25+20+15-8-5-4+2=45$; nessuna: $15$. Soltanto $N$: $25-8-5+2=14$; soltanto $F$: $10$; soltanto $M$: $8$; esattamente una: $32$.}

\begin{esercizio}[\liv{2}]
Stabilire quali delle seguenti famiglie sono partizioni dell'intervallo $[0,14]$ dei valori di pH:
\begin{gather*}
\text{(a) } [0,7[,\ \{7\},\ ]7,14];\qquad \text{(b) } [0,6],\ [6,8],\ [8,14];\\
\text{(c) } [0,6[,\ [6,8[,\ ]8,14].
\end{gather*}
\end{esercizio}
\risp{Solo (a). In (b) i valori $6$ e $8$ appartengono a due classi; in (c) il valore $8$ non appartiene ad alcuna classe.}

\begin{esercizio}[\liv{2}]
Tre siti ospitano le specie $A=\{1,2,3,4,5,6\}$, $B=\{4,5,6,7,8\}$, $C=\{1,2,9\}$. Calcolare l'indice di Jaccard (ossia il numero $J(X,Y):=|X\cap Y|/|X\cup Y|$) per ciascuna coppia di siti. Quali sono i due siti pi\`u simili?
\end{esercizio}
\risp{$J(A,B)=\frac38$, $J(A,C)=\frac27$, $J(B,C)=0$. I pi\`u simili sono $A$ e $B$.}

\begin{esercizio}[\liv{2}]
Trenta campioni d'acqua sono classificati come ``acidi'', ``neutri'' o ``basici'', e le tre classi formano una partizione. Se $12$ campioni sono acidi e $7$ neutri, quanti sono basici? Perché, per rispondere, è importante che le classi formino una partizione?
\end{esercizio}
\risp{$30-12-7=11$. Poiché le classi sono disgiunte e la loro unione è tutto l'insieme dei campioni, le numerosità si sommano (principio di addizione): ogni campione è contato una e una sola volta.}

\begin{esercizio}[\liv{3}]
Siano $A$ e $B$ insiemi finiti. Provare che $|A\setminus B|=|A|-|A\cap B|$, e che se $A\subseteq B$ allora $|A|\le|B|$.
\end{esercizio}
\risp{$A$ \`e unione disgiunta di $A\setminus B$ e $A\cap B$. Se $A\subseteq B$, $B$ \`e unione disgiunta di $A$ e $B\setminus A$, dunque $|B|=|A|+|B\setminus A|\ge|A|$.}

%-------------------------------------------------------------------
\section{Principio di induzione}

\begin{esercizio}[\liv{1}]
Sia $\mathbb I=\{n\in\N\colon 2^{n}\ge n+1\}$. Verificare che $1\in\mathbb I$ e che, se $n\in\mathbb I$, allora $n+1\in\mathbb I$. Che cosa si conclude, per il principio di induzione?
\end{esercizio}
\risp{$2^1=2\ge2$, quindi $1\in\mathbb I$. Se $2^n\ge n+1$, allora $2^{n+1}=2\cdot2^n\ge2(n+1)=2n+2\ge n+2$, quindi $n+1\in\mathbb I$. Per il principio di induzione $\mathbb I=\N$: $2^n\ge n+1$ per ogni $n\in\N$.}

\begin{esercizio}[\liv{1}]
Provare per induzione che, per ogni $n\in\N$,
\[
\sum_{k=1}^{n}k(k+1)=\frac{n(n+1)(n+2)}{3}.
\]
\end{esercizio}
\risp{Per $n=1$: $2=2$. Passo: $\frac{n(n+1)(n+2)}{3}+(n+1)(n+2)=\frac{(n+1)(n+2)(n+3)}{3}$.}

\begin{esercizio}[\liv{1}]
Provare per induzione che, per ogni $n\in\N$,
\[
\sum_{k=1}^{n}\frac{1}{k(k+1)}=\frac{n}{n+1}.
\]
\end{esercizio}
\risp{Passo: $\frac{n}{n+1}+\frac{1}{(n+1)(n+2)}=\frac{n^{2}+2n+1}{(n+1)(n+2)}=\frac{n+1}{n+2}$.}

\begin{esercizio}[\liv{2}]
Sia $q\ne1$. Provare per induzione che, per ogni $n\in\N_{0}$,
\[
\sum_{k=0}^{n}q^{k}=\frac{1-q^{n+1}}{1-q}.
\]
\end{esercizio}
\risp{Per $n=0$: $1=\frac{1-q}{1-q}$. Passo: $\frac{1-q^{n+1}}{1-q}+q^{n+1}=\frac{1-q^{n+1}+q^{n+1}-q^{n+2}}{1-q}=\frac{1-q^{n+2}}{1-q}$.}

\begin{esercizio}[\liv{2}]
Provare che $n^{3}-n$ \`e divisibile per $3$ per ogni $n\in\N$.
\end{esercizio}
\risp{Per $n=1$ si ha $0$. Passo: $(n+1)^{3}-(n+1)=(n^{3}-n)+3(n^{2}+n)$.}

\begin{esercizio}[\liv{2}]
Provare che $2^{n}>n^{2}$ per ogni $n\ge5$. Per quali $n\le4$ la disuguaglianza \`e falsa?
\end{esercizio}
\risp{Base: $32>25$. Passo: $2^{n+1}=2\cdot2^{n}>2n^{2}\ge(n+1)^{2}$, perch\'e $n^{2}-2n-1=(n-1)^{2}-2\ge0$ per $n\ge3$. La disuguaglianza \`e falsa per $n=2,3,4$ (vera per $n=1$).}

\begin{esercizio}[\liv{2}]
Una popolazione conta $P_{1}=100$ individui e aumenta di $20$ individui ogni anno: $P_{n+1}=P_{n}+20$. Un'altra conta $Q_{1}=100$ individui e aumenta del $10\,\%$ ogni anno: $Q_{n+1}=1{,}1\,Q_{n}$. Congetturare formule esplicite per $P_{n}$ e $Q_{n}$ e dimostrarle per induzione.
\end{esercizio}
\risp{$P_{n}=100+20(n-1)$, \ $Q_{n}=100\cdot1{,}1^{\,n-1}$.}

\begin{esercizio}[\liv{3}]
Trovare l'errore nel seguente ragionamento, che ``dimostra'' che tutti i cavalli hanno lo stesso colore. \textit{Sia $P(n)$: ``in ogni gruppo di $n$ cavalli, tutti hanno lo stesso colore''. $P(1)$ \`e ovvia. Supponiamo $P(n)$ vera e consideriamo $n+1$ cavalli $c_{1},\ldots,c_{n+1}$. Per l'ipotesi induttiva, $c_{1},\ldots,c_{n}$ hanno lo stesso colore, e anche $c_{2},\ldots,c_{n+1}$. Poich\'e i due gruppi hanno cavalli in comune, tutti gli $n+1$ cavalli hanno lo stesso colore.}
\end{esercizio}
\risp{Il passo da $n=1$ a $n=2$ non funziona: i gruppi $\{c_{1}\}$ e $\{c_{2}\}$ non hanno cavalli in comune.}

%-------------------------------------------------------------------
\section{Calcolo combinatorio}

\begin{esercizio}[\liv{1}]
Quanti codici di $4$ cifre (da $0$ a $9$) esistono? Quanti hanno le cifre tutte diverse?
\end{esercizio}
\risp{$10^{4}$; \ $10\cdot9\cdot8\cdot7=5040$.}

\begin{esercizio}[\liv{1}]
In quanti modi si possono disporre $6$ provette in un portaprovette con $6$ posti?
\end{esercizio}
\risp{$6!=720$.}

\begin{esercizio}[\liv{1}]
Calcolare $\binom83$, $\binom{10}{5}$, $\binom{20}{2}$, $\binom{20}{18}$.
\end{esercizio}
\risp{$56$, $252$, $190$, $190$.}

\begin{esercizio}[\liv{1}]
Tra $12$ studenti si devono scegliere un responsabile, un vice e un segretario di un gruppo di lavoro. In quanti modi? E se si devono scegliere semplicemente $3$ rappresentanti?
\end{esercizio}
\risp{$D_{12,3}=1320$; \ $\binom{12}{3}=220$.}

\begin{esercizio}[\liv{2}]
Una riserva ha $9$ sentieri. Un'escursione ne percorre $3$ distinti, uno dopo l'altro. Quante escursioni diverse sono possibili? E se si tiene conto solo di quali sentieri si percorrono?
\end{esercizio}
\risp{$9\cdot8\cdot7=504$; \ $\binom93=84$.}

\begin{esercizio}[\liv{2}]
In un campione di $15$ insetti, $5$ sono coleotteri. Se ne scelgono $4$ da esaminare. Quante scelte contengono esattamente $2$ coleotteri? Quante almeno un coleottero?
\end{esercizio}
\risp{$\binom52\binom{10}{2}=450$; \ $\binom{15}{4}-\binom{10}{4}=1365-210=1155$.}

\begin{esercizio}[\liv{2}]
Quante sequenze di DNA di lunghezza $5$ contengono almeno una volta la base A?
\end{esercizio}
\risp{Dal totale si tolgono quelle senza A: $4^{5}-3^{5}=1024-243=781$.}

\begin{esercizio}[\liv{2}]
Quanti sottoinsiemi di $\{1,2,\ldots,10\}$ contengono il numero $1$? Quanti hanno esattamente $3$ elementi e contengono il numero $1$?
\end{esercizio}
\risp{$2^{9}=512$; \ $\binom92=36$.}

\begin{esercizio}[\liv{2}]
In una riunione ciascun partecipante stringe la mano a tutti gli altri, e le strette di mano sono $45$. Quanti sono i partecipanti?
\end{esercizio}
\risp{$\binom n2=\frac{n(n-1)}{2}=45$, da cui $n=10$.}

\begin{esercizio}[\liv{2}]
Quante stringhe di $8$ cifre binarie contengono esattamente tre cifre $1$? Quante contengono almeno una cifra $1$?
\end{esercizio}
\risp{$\binom83=56$; \ $2^{8}-1=255$.}

\begin{esercizio}[\liv{2}]
Un esperimento combina $4$ temperature, $3$ salinit\`a e $2$ specie di alghe, con $3$ repliche per trattamento. Quante unit\`a sperimentali servono? Se ogni vasca contiene $6$ unit\`a, quante vasche servono?
\end{esercizio}
\risp{$4\cdot3\cdot2\cdot3=72$ unit\`a; $12$ vasche.}

\begin{esercizio}[\liv{2}]
Scrivere la riga $n=7$ del triangolo di Tartaglia e verificare che la somma dei suoi elementi \`e $2^{7}$.
\end{esercizio}
\risp{$1,7,21,35,35,21,7,1$; somma $128$.}

\begin{esercizio}[\liv{2}]
Determinare il coefficiente di $x^{3}$ nello sviluppo di $(1+2x)^{6}$ e il termine che non contiene $x$ nello sviluppo di $\left(x+\frac1x\right)^{6}$.
\end{esercizio}
\risp{$\binom63\,2^{3}=160$; \ $\binom63=20$.}

\begin{esercizio}[\liv{3}]
Quanti sono gli anagrammi (anche privi di significato) della parola NATURA? E della parola AMBIENTE? (Suggerimento: contare prima gli anagrammi supponendo le lettere uguali distinguibili, poi chiedersi quante volte \`e stato contato ciascun anagramma.)
\end{esercizio}
\risp{$\frac{6!}{2!}=360$; \ $\frac{8!}{2!}=20\,160$.}

\begin{esercizio}[\liv{1}]
Un menù offre $3$ primi, $4$ secondi e $2$ dolci. Quanti pasti diversi (un primo, un secondo e un dolce) si possono comporre?
\end{esercizio}
\risp{$3\cdot4\cdot2=24$ (principio di moltiplicazione).}

\begin{esercizio}[\liv{1}]
Calcolare $5!$, \ $\dfrac{7!}{5!}$, \ $\dfrac{10!}{8!\,2!}$ senza calcolare i fattoriali per intero.
\end{esercizio}
\risp{$120$; \ $7\cdot6=42$; \ $\frac{10\cdot9}{2}=45$.}

\begin{esercizio}[\liv{3}]
Provare che $k\binom nk=n\binom{n-1}{k-1}$ per $1\le k\le n$, sia con la formula per $\binom nk$ sia con un argomento combinatorio. (Suggerimento: contare in due modi le commissioni di $k$ persone scelte tra $n$, con un presidente scelto tra i membri della commissione.)
\end{esercizio}
\risp{Con la formula: $k\binom nk=k\frac{n!}{k!(n-k)!}=\frac{n!}{(k-1)!(n-k)!}=n\frac{(n-1)!}{(k-1)!(n-k)!}=n\binom{n-1}{k-1}$. Contando in due modi: si sceglie la commissione ($\binom nk$ modi) e poi il presidente al suo interno ($k$ modi); oppure si sceglie prima il presidente ($n$ modi) e poi gli altri $k-1$ membri fra le $n-1$ persone rimaste ($\binom{n-1}{k-1}$ modi).}

\begin{esercizio}[\liv{3}]
Provare che, per $n\in\N$,
\[
\sum_{k=0}^{n}(-1)^{k}\binom nk=0.
\]
Dedurre che un insieme con $n$ elementi ha tanti sottoinsiemi con un numero pari di elementi quanti con un numero dispari, cio\`e $2^{n-1}$.
\end{esercizio}
\risp{Con $a=1$, $b=-1$ nel binomio di Newton: $0=(1-1)^n=\sum_{k=0}^n\binom nk(-1)^k$. La somma dei $\binom nk$ con $k$ pari è quindi uguale a quella con $k$ dispari; poiché in tutto i sottoinsiemi sono $2^n$, ciascuna delle due somme vale $2^{n-1}$ (per $n\ge1$).}

%-------------------------------------------------------------------
\section{Numeri razionali e irrazionali}

\begin{esercizio}[\liv{1}]
Scrivere come frazioni ridotte ai minimi termini: $0{,}\overline{7}$, \ $0{,}1\overline{2}$, \ $0{,}0\overline{45}$, \ $3{,}\overline{142}$.
\end{esercizio}
\risp{$\frac79$; \ $\frac{11}{90}$; \ $\frac{1}{22}$; \ $\frac{3139}{999}$.}

\begin{esercizio}[\liv{1}]
Calcolare la rappresentazione decimale di $\frac38$, $\frac7{12}$, $\frac{13}{40}$, $\frac{5}{6}$. Quali sono finite? Formulare una congettura sui denominatori delle frazioni (ridotte) che hanno rappresentazione finita.
\end{esercizio}
\risp{$0{,}375$; $0{,}58\overline{3}$; $0{,}325$; $0{,}8\overline{3}$. Una frazione ridotta ai minimi termini ha rappresentazione finita se e solo se il denominatore non ha fattori primi diversi da $2$ e $5$.}

\begin{esercizio}[\liv{2}]
Calcolare le rappresentazioni decimali di $\frac17$ e di $\frac{1}{13}$. Confrontare la lunghezza del periodo con la stima che si ricava dalla dimostrazione del Teorema 3.2 degli appunti.
\end{esercizio}
\risp{$\frac17=0{,}\overline{142857}$, $\frac1{13}=0{,}\overline{076923}$: entrambi hanno periodo di $6$ cifre (la stima d\`a al pi\`u $6$ e $12$).}

\begin{esercizio}[\liv{2}]
Siano $q\in\Q$ e $x\notin\Q$. Provare che $q+x\notin\Q$ e che, se $q\ne0$, $qx\notin\Q$.
\end{esercizio}
\risp{Se fosse $q+x=r\in\Q$, si avrebbe $x=r-q\in\Q$, assurdo. Analogamente $x=\frac{qx}{q}$.}

\begin{esercizio}[\liv{2}]
La somma di due numeri irrazionali \`e sempre irrazionale? E il prodotto?
\end{esercizio}
\risp{No in entrambi i casi: $\sqrt2+(-\sqrt2)=0$, $\sqrt2\cdot\sqrt2=2$.}

\begin{esercizio}[\liv{2}]
Usando l'esercizio precedente e l'irrazionalità di $\sqrt2$, provare che $3\sqrt2$, \ $1+\sqrt2$ \ e \ $\frac{\sqrt2}{5}$ sono irrazionali.
\end{esercizio}
\risp{$3\sqrt2$ e $\frac{\sqrt2}5$ sono prodotti di un razionale non nullo ($3$ e $\frac15$) per l'irrazionale $\sqrt2$; $1+\sqrt2$ è somma di un razionale e di un irrazionale.}

%-------------------------------------------------------------------
\section{Numeri reali e ordine}

\begin{esercizio}[\liv{1}]
Stabilire quali dei seguenti numeri sono razionali, giustificando la risposta (nel secondo, a ogni blocco si aggiunge un $1$):
\begin{gather*}
0{,}\overline{12},\qquad 0{,}12112111211112\ldots,\qquad \sqrt{\tfrac{16}{9}},\\
3{,}14,\qquad \sqrt2\cdot\sqrt8 .
\end{gather*}
\end{esercizio}
\risp{$0{,}\overline{12}=\frac{12}{99}=\frac4{33}$ \`e razionale (periodico); il secondo non \`e periodico, quindi \`e irrazionale; $\sqrt{16/9}=\frac43$; $3{,}14=\frac{157}{50}$; $\sqrt2\cdot\sqrt8=\sqrt{16}=4$. Sono tutti razionali tranne il secondo.}

\begin{esercizio}[\liv{1}]
Disporre in ordine crescente:
\begin{gather*}
1{,}41,\qquad \sqrt2=1{,}41421\ldots,\qquad 1{,}4142,\\
\tfrac{99}{70},\qquad 1{,}\overline{41}.
\end{gather*}
\end{esercizio}
\risp{$\frac{99}{70}=1{,}4\overline{142857}=1{,}41428\ldots$ e $1{,}\overline{41}=1{,}414141\ldots$. L'ordine \`e $1{,}41<1{,}\overline{41}<1{,}4142<\sqrt2<\frac{99}{70}$.}

\begin{esercizio}[\liv{2}]
Trovare un numero razionale e un numero irrazionale compresi fra $0{,}4999$ e $0{,}5$.
\end{esercizio}
\risp{Ad esempio $0{,}49995$ (razionale) e $0{,}4999+\sqrt2\cdot10^{-5}=0{,}49991414\ldots$ (irrazionale: somma di un razionale e di un irrazionale, esercizio della sezione precedente).}

\begin{esercizio}[\liv{2}]
Mostrare con un controesempio che da $x<y$ non segue $x^2<y^2$. Provare che invece l'implicazione vale se $0\le x<y$.
\end{esercizio}
\risp{$-3<1$ ma $9>1$. Se $0\le x<y$, moltiplicando $x<y$ per $x\ge0$ e per $y>0$ si ha $x^2\le xy<y^2$.}

%-------------------------------------------------------------------
\section{Intervalli e valore assoluto}

\begin{esercizio}[\liv{1}]
Scrivere come intervalli (o unioni di intervalli) e rappresentare sulla retta:
\begin{gather*}
\{x\colon -2<x\le5\},\qquad \{x\colon x\ge3\}\cap\{x\colon x<7\},\\
\{x\colon x<0\}\cup\{x\colon x\ge1\},\qquad \R\setminus[0,1[.
\end{gather*}
\end{esercizio}
\risp{$]-2,5]$; \ $[3,7[$; \ $]-\infty,0[\,\cup\,[1,+\infty[$; \ $]-\infty,0[\,\cup\,[1,+\infty[$ (il terzo e il quarto coincidono).}

\begin{esercizio}[\liv{1}]
Scrivere con un'unica disuguaglianza che contiene un valore assoluto: $]1,5[$; \ $[-4,2]$; \ $]-\infty,-1[\,\cup\,]1,+\infty[$; \ un intervallo chiuso $[a,b]$ qualsiasi.
\end{esercizio}
\risp{$|x-3|<2$; \ $|x+1|\le3$; \ $|x|>1$; \ $\bigl|x-\frac{a+b}2\bigr|\le\frac{b-a}2$ (centro e semiampiezza).}

\begin{esercizio}[\liv{2}]
Un protocollo prescrive che il pH di una vasca sia $7{,}2$ con una tolleranza di $0{,}3$. Scrivere la condizione con il valore assoluto e come intervallo. Stesso esercizio per una temperatura di $20\,^\circ$C con tolleranza $1{,}5\,^\circ$C.
\end{esercizio}
\risp{$|\mathrm{pH}-7{,}2|\le0{,}3$, cio\`e $[6{,}9;\,7{,}5]$. $|T-20|\le1{,}5$, cio\`e $[18{,}5;\,21{,}5]$.}

\begin{esercizio}[\liv{2}]
Calcolare $|x-2|+|x+3|$ per $x=0$, $x=5$, $x=-4$. In tutti e tre i casi il risultato è almeno $5$: perché, pensando alle distanze sulla retta?
\end{esercizio}
\risp{$5$, $11$, $7$. $|x-2|+|x+3|$ è la somma delle distanze di $x$ da $2$ e da $-3$: non può essere minore della distanza fra $2$ e $-3$, che è $5$ (è la disuguaglianza triangolare). Vale esattamente $5$ quando $x$ sta fra $-3$ e $2$.}

\begin{esercizio}[\liv{1}]
Calcolare $|-3|$, \ $|2-5|$, \ $|\sqrt2-2|$, \ $|3-\pi|$. Verificare poi la disuguaglianza $|x+y|\le|x|+|y|$ per $x=3,\ y=4$ e per $x=3,\ y=-4$: in quale dei due casi vale l'uguaglianza?
\end{esercizio}
\risp{$3$; \ $3$; \ $2-\sqrt2$; \ $\pi-3$. Per $x=3$, $y=4$: $7=7$, uguaglianza. Per $x=3$, $y=-4$: $|{-1}|=1<7$. L'uguaglianza vale quando $x$ e $y$ hanno lo stesso segno.}

%-------------------------------------------------------------------
\section{Problemi di riepilogo}

\begin{esercizio}[\liv{2}]
In un monitoraggio si vogliono scegliere $3$ siti tra $10$ e, in ciascun sito scelto, campionare in $2$ delle $4$ stagioni. In quanti modi si pu\`o costruire il piano di campionamento?
\end{esercizio}
\risp{$\binom{10}{3}\cdot\binom42^{3}=120\cdot216=25\,920$.}

\begin{esercizio}[\liv{1}]
In una stazione si registrano le temperature $12{,}4$; $13{,}1$; $12{,}9$; $13{,}6\ ^\circ$C. Scrivere il più piccolo intervallo chiuso che le contiene tutte, e scriverlo nella forma $|T-c|\le\delta$.
\end{esercizio}
\risp{$[12{,}4;\,13{,}6]$; centro $c=13$, semiampiezza $\delta=0{,}6$: $|T-13|\le0{,}6$.}

\begin{esercizio}[\liv{3}]
Un'indagine su $200$ laghi rileva che il $40\,\%$ ha concentrazioni elevate di fosforo, il $25\,\%$ ha fioriture algali e $30$ laghi presentano entrambe le condizioni. Quanti laghi non presentano nessuna delle due condizioni? Quale percentuale dei laghi con fioriture algali ha concentrazioni elevate di fosforo?
\end{esercizio}
\risp{$80$ laghi con fosforo elevato, $50$ con fioriture; almeno una condizione: $80+50-30=100$; nessuna: $100$. Dei $50$ laghi con fioriture, $30$ hanno fosforo elevato: il $60\,\%$.}

%-------------------------------------------------------------------
\stampasoluzioni
   % Lezioni 1--3
% !TEX root =esercizi.tex
\chapter{Lezioni 4--6 (30 settembre -- 5 ottobre 2026)}

\begin{vuoto}[Argomenti svolti]{}
\begin{itemize}
\item[\bul] \textbf{Lezione 4 --- Equazioni e disequazioni di primo grado}. Principi di equivalenza; equazioni di primo grado e fratte, condizioni di esistenza. Disequazioni di primo grado; sistemi di disequazioni. Equazioni e disequazioni con il valore assoluto.
\item[\bul] \textbf{Lezione 5 --- Equazioni e disequazioni di secondo grado}. Prodotti notevoli e scomposizioni; legge di annullamento del prodotto. Formula risolutiva, discriminante, scomposizione del trinomio. Segno del trinomio e metodo della parabola. Disequazioni di secondo grado; segno di prodotti e quozienti (metodo dei due sistemi, tabella dei segni), anche con fattori di primo grado (Cap.~4, \S4.5); biquadratiche, sistemi. Notazione scientifica (Cap.~6, \S6.7) e percentuali (Cap.~7, \S7.1).
\item[\bul] \textbf{Lezione 6 --- Completezza, potenze e radici}. Maggioranti e minoranti, estremo superiore e inferiore, massimo e minimo; proprietà di Archimede; completezza di $\R$. Potenze a esponente intero e razionale, cenno all'esponente reale ($3^{\sqrt2}$ come estremo superiore); radici $n$-esime, $\sqrt{x^2}=|x|$. Disequazioni $x^n\gtrless a$.
\end{itemize}
\end{vuoto}

Unità di misura, ordini di grandezza ed equazioni e disequazioni con la radice quadrata non sono ancora stati svolti: i relativi esercizi sono rinviati alla Scheda~3. La caratterizzazione dell'estremo superiore e inferiore con $\eps$, che sarà svolta in L7, è richiamata nel riquadro all'inizio della sezione~2.8.

Il livello di difficoltà è indicato con \liv{1} (esercizio di base), \liv{2} (esercizio standard), \liv{3} (esercizio più impegnativo). In tutti gli esercizi la soluzione va scritta come intervallo o unione di intervalli.

%-------------------------------------------------------------------
\section{Equazioni e disequazioni di primo grado}

\begin{esercizio}[\liv{1}]
Risolvere: \ $2x+3=7$; \ $5-x=2x-4$; \ $\dfrac x3+1=\dfrac x2$.
\end{esercizio}
\risp{$x=2$; \ $x=3$; \ $1=\frac x6$, cioè $x=6$.}

\begin{esercizio}[\liv{1}]
Tradurre in un'equazione e risolvere: ``il doppio di un numero, diminuito di $5$, è uguale a $11$''.
\end{esercizio}
\risp{$2x-5=11$, $x=8$.}

\begin{esercizio}[\liv{1}]
Risolvere, dicendo se l'equazione è determinata, impossibile o indeterminata:
\[
3(x-2)-2(x+1)=x-8,\qquad \frac{x+1}{2}-\frac{x-1}{3}=1,\qquad 2(x-1)=2x+3 .
\]
\end{esercizio}
\risp{La prima si riduce a $x-8=x-8$: indeterminata. La seconda dà $3x+3-2x+2=6$, cioè $x=1$. La terza dà $-2=3$: impossibile.}

\begin{esercizio}[\liv{1}]
Risolvere e rappresentare sulla retta:
\[
5-2x\ge3x+20,\qquad \frac{x-2}{4}>\frac{2x+1}{3},\qquad 0{,}2x-1{,}5<0{,}5x .
\]
\end{esercizio}
\risp{$-5x\ge15$, cioè $x\le-3$; \ $3x-6>8x+4$, cioè $-5x>10$, cioè $x<-2$; \ $-1{,}5<0{,}3x$, cioè $x>-5$.}

\begin{esercizio}[\liv{2}]
Risolvere, scrivendo prima le condizioni di esistenza:
\[
\frac{3}{x-1}=\frac{2}{x+2},\qquad \frac{x}{x-2}-\frac{2}{x-2}=3 .
\]
\end{esercizio}
\risp{Prima: $x\ne1,-2$; $3x+6=2x-2$, $x=-8$, accettabile. Seconda: $x\ne2$; per $x\ne2$ il primo membro vale $\frac{x-2}{x-2}=1\ne3$: impossibile.}

\begin{esercizio}[\liv{2}]
Un bacino di raccolta contiene $12\,000$ L; da esso si prelevano $150$ L all'ora e vi affluiscono $90$ L all'ora. Dopo quante ore il contenuto scende sotto i $3000$ L? In quale intervallo di tempo il contenuto è compreso fra $6000$ e $9000$ L?
\end{esercizio}
\risp{Il contenuto è $12\,000-60t$ (con $t$ in ore). $12\,000-60t<3000$ dà $t>150$ h. $6000\le12\,000-60t\le9000$ dà $50\le t\le100$ h.}

%-------------------------------------------------------------------
\section{Sistemi di disequazioni}

\begin{esercizio}[\liv{1}]
Risolvere e rappresentare sulla retta:
\[
\begin{cases}x-1>0\\ x<4\end{cases}\qquad \begin{cases}2x\ge6\\ x-5\le0\end{cases}\qquad \begin{cases}x>2\\ x<-1\end{cases}
\]
\end{esercizio}
\risp{$]1,4[$; \ $[3,5]$; \ nessuna soluzione (le due semirette non si intersecano).}

\begin{esercizio}[\liv{1}]
Risolvere:
\[
\begin{cases}2x+1>3\\ 4-x\ge0\end{cases}\qquad
\begin{cases}3x-2\le x\\ x+5>2x\end{cases}\qquad
-1\le\frac{3-x}{2}<2 .
\]
\end{esercizio}
\risp{$]1,4]$; \ $x\le1$ e $x<5$, cioè $]-\infty,1]$; \ $-2\le3-x<4$, cioè $-1<x\le5$: $]-1,5]$.}

\begin{esercizio}[\liv{2}]
Risolvere il sistema $\begin{cases}\frac{x-3}{2}\ge x-5\\ 2(x+1)>3x-4\\ x\ge0\end{cases}$.
\end{esercizio}
\risp{La prima dà $x\le7$, la seconda $x<6$: con $x\ge0$ la soluzione è $[0,6[$.}

%-------------------------------------------------------------------
\section{Segno di prodotti e quozienti}

\begin{esercizio}[\liv{1}]
Per quali $x$ il quoziente $\dfrac{x-3}{x+2}$ è positivo? Negativo? Nullo? Per quali non è definito?
\end{esercizio}
\risp{Positivo per $x<-2$ oppure $x>3$; negativo per $-2<x<3$; nullo per $x=3$; non definito per $x=-2$.}

\begin{esercizio}[\liv{1}]
Risolvere $(x-1)(x-4)>0$ e $(x-1)(x-4)\le0$.
\end{esercizio}
\risp{$]-\infty,1[\,\cup\,]4,+\infty[$; \ $[1,4]$.}

\begin{esercizio}[\liv{1}]
Risolvere:
\[
\frac{x+3}{2-x}\ge0,\qquad (2x-1)(x+4)<0,\qquad \frac{3-x}{x+1}\le0 .
\]
\end{esercizio}
\risp{$[-3,2[$; \ $]-4,\frac12[$; \ $]-\infty,-1[\,\cup\,[3,+\infty[$.}

\begin{esercizio}[\liv{2}]
Risolvere, con la tabella dei segni:
\[
\frac{(x-1)(x+2)}{x-4}\ge0,\qquad \frac{2}{x-1}\ge1,\qquad \frac1{x-1}<\frac1{x+1} .
\]
\end{esercizio}
\risp{Prima: $[-2,1]\,\cup\,]4,+\infty[$. Seconda: $\frac{3-x}{x-1}\ge0$, cioè $]1,3]$. Terza: $\frac{2}{(x-1)(x+1)}<0$, cioè $(x-1)(x+1)<0$: $]-1,1[$.}

\begin{esercizio}[\liv{3}]
Si mescolano $x$ litri di una soluzione al $30\%$ con $5$ litri di una soluzione al $10\%$. Scrivere la concentrazione della miscela in funzione di $x$ e determinare per quali $x$ essa supera il $20\%$. Perché in questo caso si può moltiplicare per il denominatore senza studiarne il segno?
\end{esercizio}
\risp{$c(x)=\frac{0{,}3x+0{,}5}{x+5}$. Per $x\ge0$ il denominatore è positivo, quindi $c>0{,}2$ equivale a $0{,}3x+0{,}5>0{,}2x+1$, cioè $x>5$ L.}

%-------------------------------------------------------------------
\section{Valore assoluto}

\begin{esercizio}[\liv{1}]
Calcolare $|x-3|$ per $x=5$, $x=3$, $x=-1$. Scrivere poi $|x-3|$ senza il valore assoluto, distinguendo i casi $x\ge3$ e $x<3$.
\end{esercizio}
\risp{$2$, $0$, $4$. \ $|x-3|=x-3$ se $x\ge3$, \ $|x-3|=3-x$ se $x<3$.}

\begin{esercizio}[\liv{1}]
Risolvere: \ $|x|=4$; \ $|x|<2$; \ $|x|\ge1$; \ $|x|=-3$.
\end{esercizio}
\risp{$x=\pm4$; \ $]-2,2[$; \ $]-\infty,-1]\cup[1,+\infty[$; \ nessuna soluzione.}

\begin{esercizio}[\liv{1}]
Risolvere:
\begin{gather*}
|2x-3|=5,\qquad |x+1|<3,\\
|4-x|\ge1,\qquad |3x+2|\le-1 .
\end{gather*}
\end{esercizio}
\risp{$x=4$ oppure $x=-1$; \ $]-4,2[$; \ $]-\infty,3]\cup[5,+\infty[$; \ nessuna soluzione.}

\begin{esercizio}[\liv{2}]
Risolvere:
\begin{gather*}
|x-2|=|x+4|,\qquad |x|+|x-2|=4,\\
|x-1|>2x,\qquad |x-1|+|x+1|<4 .
\end{gather*}
\end{esercizio}
\risp{Prima: punto equidistante da $2$ e $-4$, $x=-1$. Seconda: per $x<0$, $2-2x=4$, $x=-1$; per $0\le x<2$, $2=4$, mai; per $x\ge2$, $2x-2=4$, $x=3$; soluzioni $-1$ e $3$. Terza: per $x\ge1$, $x-1>2x$ dà $x<-1$, incompatibile; per $x<1$, $1-x>2x$ dà $x<\frac13$; soluzione $]-\infty,\frac13[$. Quarta: $]-2,2[$.}

\begin{esercizio}[\liv{3}]
Una specie di pesci vive fra $8$ e $24\,^\circ$C. In un lago, la temperatura alla profondità $h$ (in metri, $0\le h\le20$) è $T(h)=22-0{,}8h$. Scrivere la condizione sulla temperatura con un valore assoluto e determinare a quali profondità la specie trova condizioni adatte.
\end{esercizio}
\risp{$|T-16|\le8$, cioè $|6-0{,}8h|\le8$: $-8\le6-0{,}8h\le8$, cioè $-2{,}5\le h\le17{,}5$. Con $0\le h\le20$: fra $0$ e $17{,}5$ m.}

%-------------------------------------------------------------------
\section{Equazioni di secondo grado e scomposizioni}

\begin{esercizio}[\liv{1}]
Scomporre in fattori:
\begin{gather*}
x^2-10x+25,\qquad 4x^2-1,\qquad x^3+27,\\
x^3-x^2-4x+4,\qquad 2x^2-7x+3 .
\end{gather*}
\end{esercizio}
\risp{$(x-5)^2$; \ $(2x-1)(2x+1)$; \ $(x+3)(x^2-3x+9)$; \ $(x-1)(x-2)(x+2)$; \ $(x-3)(2x-1)$.}

\begin{esercizio}[\liv{1}]
Risolvere, usando la formula risolutiva solo quando serve:
\begin{gather*}
x^2-x-12=0,\qquad 5x^2-20x=0,\qquad 9x^2-4=0,\\
x^2+6x+10=0,\qquad x^2-2\sqrt2\,x+2=0 .
\end{gather*}
\end{esercizio}
\risp{$4$ e $-3$; \ $0$ e $4$; \ $\pm\frac23$; \ nessuna soluzione ($\Delta=-4$); \ $(x-\sqrt2)^2=0$, cioè $x=\sqrt2$.}

\begin{esercizio}[\liv{2}]
(a) Risolvere $(x-1)(x+2)=4$ e $\dfrac{x}{x-1}+\dfrac1x=2$. (b) Per quali $k$ l'equazione $x^2-4x+k=0$ ha due soluzioni distinte? (c) Trovare due numeri di somma $10$ e prodotto $21$.
\end{esercizio}
\risp{(a) $x^2+x-6=0$: $2$ e $-3$. Con $x\ne0,1$: $x^2+x-1=2x^2-2x$, cioè $x^2-3x+1=0$, $x=\frac{3\pm\sqrt5}2$, entrambe accettabili. (b) $\Delta=16-4k>0$: $k<4$. (c) Sono le soluzioni di $t^2-10t+21=0$: $3$ e $7$.}

%-------------------------------------------------------------------
\section{Disequazioni di secondo grado}

\begin{esercizio}[\liv{1}]
Risolvere: \ $x^2<9$; \ $x^2\ge4$; \ $x^2-x\ge0$; \ $x^2+4>0$.
\end{esercizio}
\risp{$]-3,3[$; \ $]-\infty,-2]\cup[2,+\infty[$; \ $x(x-1)\ge0$: $]-\infty,0]\cup[1,+\infty[$; \ ogni $x\in\R$.}

\begin{esercizio}[\liv{1}]
Risolvere:
\begin{gather*}
x^2-3x-10\le0,\qquad 2x^2-x-1>0,\qquad -x^2+2x-1\ge0,\\
x^2+x+1>0,\qquad 4-x^2<0 .
\end{gather*}
\end{esercizio}
\risp{$[-2,5]$; \ $]-\infty,-\frac12[\,\cup\,]1,+\infty[$; \ $-(x-1)^2\ge0$, cioè $x=1$; \ ogni $x\in\R$; \ $]-\infty,-2[\,\cup\,]2,+\infty[$.}

\begin{esercizio}[\liv{2}]
Risolvere:
\begin{gather*}
\frac{x^2-4}{x+1}\ge0,\qquad \frac{x^2-5x+6}{x^2+x+1}<0,\\
(x-2)(x^2-2x-3)>0,\qquad \frac{x}{x-2}\ge\frac{3}{x+2} .
\end{gather*}
\end{esercizio}
\risp{Prima: $[-2,-1[\,\cup\,[2,+\infty[$. Seconda: il denominatore è sempre positivo: $]2,3[$. Terza: $(x-2)(x-3)(x+1)>0$: $]-1,2[\,\cup\,]3,+\infty[$. Quarta: $\frac{x^2-x+6}{(x-2)(x+2)}\ge0$, con numeratore sempre positivo: $]-\infty,-2[\,\cup\,]2,+\infty[$.}

\begin{esercizio}[\liv{2}]
Risolvere il sistema $\begin{cases}x^2-5x+4\le0\\ x^2-4>0\end{cases}$ e la disequazione biquadratica $x^4-13x^2+36\le0$.
\end{esercizio}
\risp{Sistema: $[1,4]$ intersecato con $]-\infty,-2[\,\cup\,]2,+\infty[$: $]2,4]$. Biquadratica: con $t=x^2$, $4\le t\le9$, cioè $2\le|x|\le3$: $[-3,-2]\cup[2,3]$.}

\begin{esercizio}[\liv{3}]
Un prato rettangolare ha perimetro $200$ m. Per quali lunghezze di un lato l'area supera $2100\ \mathrm{m^2}$? Qual è l'area massima possibile? (Completare il quadrato.)
\end{esercizio}
\risp{Con lati $x$ e $100-x$: $x(100-x)>2100$, cioè $x^2-100x+2100<0$, cioè $30<x<70$ m. Poiché $x(100-x)=2500-(x-50)^2$, l'area massima è $2500\ \mathrm{m^2}$, per il quadrato di lato $50$ m.}

%-------------------------------------------------------------------
\section{Valore assoluto e radici}

\begin{esercizio}[\liv{1}]
Scrivere le condizioni di esistenza di $\sqrt{3-x}$, \ $\sqrt{x^2-1}$, \ $\dfrac{1}{\sqrt{x+4}}$.
\end{esercizio}
\risp{$x\le3$; \ $x^2\ge1$, cioè $x\le-1$ oppure $x\ge1$; \ $x>-4$ (la radice deve esistere ed essere diversa da zero).}

\begin{esercizio}[\liv{1}]
Risolvere: \ $\sqrt{x-2}=3$; \ $\sqrt{x+1}=-2$; \ $\sqrt{x^2}=5$.
\end{esercizio}
\risp{$x-2=9$, $x=11$; \ nessuna soluzione (una radice non è mai negativa); \ $|x|=5$, cioè $x=\pm5$.}

\begin{esercizio}[\liv{2}]
Risolvere:
\[
|x^2-4|<3,\qquad |x^2-x|=2,\qquad x^2-5|x|+6<0 .
\]
\end{esercizio}
\risp{Prima: $1<x^2<7$, cioè $1<|x|<\sqrt7$: $]-\sqrt7,-1[\,\cup\,]1,\sqrt7[$. Seconda: $x^2-x=2$ dà $2$ e $-1$; $x^2-x=-2$ non ha soluzioni. Terza: con $t=|x|$, $2<t<3$: $]-3,-2[\,\cup\,]2,3[$.}

%-------------------------------------------------------------------
\section{Estremo superiore e completezza}

\begin{vuoto}[Caratterizzazione dell'estremo superiore e inferiore]{}
Sia $A\subseteq\R$ non vuoto. Un numero $s$ è l'estremo superiore di $A$ se e solo se
\begin{itemize}
\item[(i)] $x\le s$ per ogni $x\in A$ \hfill ($s$ è un maggiorante);
\item[(ii)] per ogni $\eps>0$ esiste $x\in A$ con $x>s-\eps$ \hfill ($s-\eps$ non è un maggiorante).
\end{itemize}
Analogamente, $i$ è l'estremo inferiore di $A$ se e solo se
\begin{itemize}
\item[(i)] $x\ge i$ per ogni $x\in A$ \hfill ($i$ è un minorante);
\item[(ii)] per ogni $\eps>0$ esiste $x\in A$ con $x<i+\eps$ \hfill ($i+\eps$ non è un minorante).
\end{itemize}
Nella (ii) l'ordine dei quantificatori conta: prima si fissa $\eps$, poi si cerca $x$, che in generale dipende da $\eps$. Negli esempi con $\frac1n$ la (ii) si verifica con la \emph{proprietà di Archimede}: per ogni $\eps>0$ esiste $n\in\N$ con $n>\frac1\eps$, cioè $\frac1n<\eps$.
\end{vuoto}

\begin{esercizio}[\liv{1}]
Per ciascun insieme dire se ha massimo, minimo, estremo superiore, estremo inferiore: \ $[2,5[$; \ $]0,+\infty[$; \ $\{1,2,3\}$.
\end{esercizio}
\risp{$[2,5[$: minimo $2$, estremo superiore $5$, nessun massimo. \ $]0,+\infty[$: estremo inferiore $0$, nessun minimo; non è limitato superiormente, quindi non ha estremo superiore in $\R$. \ $\{1,2,3\}$: minimo $1$, massimo $3$.}

\begin{esercizio}[\liv{1}]
Per ciascun insieme determinare maggioranti, estremo superiore e inferiore, e dire se esistono massimo e minimo:
\begin{gather*}
]-1,3],\qquad \Bigl\{\frac1n\colon n\in\N\Bigr\},\\
\{x\in\R\colon x^2\le4\},\qquad \Bigl\{(-1)^n+\frac1n\colon n\in\N\Bigr\}.
\end{gather*}
\end{esercizio}
\risp{$]-1,3]$: $\sup=\max=3$, $\inf=-1$, non minimo. $\{1/n\}$: $\sup=\max=1$, $\inf=0$, non minimo. $[-2,2]$: $\max=2$, $\min=-2$. Nell'ultimo i termini sono $0,\frac32,-\frac23,\frac54,\ldots$: $\sup=\max=\frac32$ (per $n=2$), $\inf=-1$, non minimo.}

\begin{esercizio}[\liv{1}]
Determinare, se esistono, massimo, minimo, estremo superiore ed estremo inferiore di:
\begin{gather*}
\N,\qquad \Z,\qquad \Bigl\{-\frac1n\colon n\in\N\Bigr\},\\
\Bigl\{\frac{n}{n+1}\colon n\in\N\Bigr\},\qquad \{x\in\R\colon |x-2|<1\},\qquad ]-\infty,3]\cup\{5\}.
\end{gather*}
\end{esercizio}
\risp{$\N$: $\min=1$, non limitato superiormente ($\sup\N=+\infty$, per Archimede). \ $\Z$: $\sup=+\infty$, $\inf=-\infty$. \ $\{-1,-\frac12,-\frac13,\ldots\}$: $\min=-1$, $\sup=0$, non massimo. \ $\{\frac12,\frac23,\frac34,\ldots\}$: $\min=\frac12$, $\sup=1$, non massimo. \ $]1,3[$: $\inf=1$, $\sup=3$, né minimo né massimo. \ $\max=5$, $\inf=-\infty$.}

\begin{esercizio}[\liv{1}]
Vero o falso?
(a) Ogni insieme finito non vuoto ha massimo.
(b) Se $A$ ha massimo, allora $\max A=\sup A$.
(c) Se $\sup A\in A$, allora $\sup A$ è il massimo di $A$.
(d) Ogni insieme non vuoto e limitato superiormente ha massimo.
(e) Se $M$ è un maggiorante di $A$, anche $M+1$ lo è.
(f) Un insieme può avere due estremi superiori diversi.
\end{esercizio}
\risp{(a) V. (b) V: il massimo è un maggiorante, e ogni altro maggiorante è $\ge$ di esso perché il massimo sta in $A$. (c) V. (d) F: $]0,1[$. (e) V. (f) F: il minimo dei maggioranti, se esiste, è uno solo.}

\begin{esercizio}[\liv{1}]
Una centralina registra la temperatura dell'aria ogni ora, per un mese. (a) L'insieme delle temperature registrate ha sempre massimo e minimo? Perché? (b) L'insieme $\{15-\frac1n\colon n\in\N\}$ descrive temperature che si avvicinano a $15\,^\circ$C senza raggiungerlo: determinarne massimo, minimo, estremo superiore ed estremo inferiore.
\end{esercizio}
\risp{(a) Sì: è un insieme finito (al più $31\cdot24=744$ valori), e fra un numero finito di numeri ce n'è sempre uno più grande e uno più piccolo di tutti. (b) $\min=14$ (per $n=1$), $\sup=15$, non massimo.}

\begin{esercizio}[\liv{2}]
Usando la caratterizzazione dell'estremo superiore, provare che $\sup\{1-\frac1n\colon n\in\N\}=1$. Quale proprietà di $\R$ si usa?
\end{esercizio}
\risp{(i) $1-\frac1n\le1$ per ogni $n$. (ii) Dato $\eps>0$, per la proprietà di Archimede esiste $n>\frac1\eps$, e allora $1-\frac1n>1-\eps$.}

\begin{esercizio}[\liv{2}]
Usando la caratterizzazione dell'estremo inferiore, provare che $\inf\{2+\frac3n\colon n\in\N\}=2$. L'insieme ha minimo? Ha massimo?
\end{esercizio}
\risp{(i) $2+\frac3n>2$ per ogni $n$. (ii) Dato $\eps>0$, per Archimede esiste $n>\frac3\eps$, e allora $2+\frac3n<2+\eps$. Non ha minimo ($2$ non è della forma $2+\frac3n$); ha massimo $5$, per $n=1$.}

\begin{esercizio}[\liv{2}]
Determinare estremo superiore ed estremo inferiore di $\{x^2\colon -1<x<2\}$ e di $\{\frac1x\colon x>0\}$, dicendo se sono massimo e minimo. (Disegnare il grafico di $x^2$ aiuta.)
\end{esercizio}
\risp{Al variare di $x$ in $]-1,2[$, $x^2$ descrive $[0,4[$: $\min=0$ (per $x=0$), $\sup=4$, non massimo. Al variare di $x>0$, $\frac1x$ descrive $]0,+\infty[$: $\inf=0$, non minimo; non limitato superiormente.}

\begin{esercizio}[\liv{2}]
Siano $A\subseteq B$ due insiemi non vuoti, con $B$ limitato superiormente. Provare che $A$ è limitato superiormente e che $\sup A\le\sup B$. Mostrare con un esempio che può essere $\sup A=\sup B$ anche se $A\ne B$.
\end{esercizio}
\risp{$\sup B$ è un maggiorante di $B$, quindi anche di $A$ (ogni elemento di $A$ sta in $B$). Allora $A$ è limitato superiormente, e $\sup A$, che è il più piccolo dei maggioranti di $A$, è $\le\sup B$. Esempio: $A=]0,1[$, $B=[0,1]$.}

\begin{esercizio}[\liv{2}]
Sia $A=\{x\in\Q\colon x>0,\ x^2<3\}$. Qual è l'estremo superiore di $A$ in $\R$? Appartiene a $\Q$? Che cosa dice questo esempio su $\Q$?
\end{esercizio}
\risp{$\sup A=\sqrt3$, che è irrazionale (esercizio della Scheda~1): come sottoinsieme di $\Q$, $A$ è limitato ma non ha estremo superiore in $\Q$. È l'esempio che mostra che $\Q$ non è completo.}

\begin{esercizio}[\liv{3}]
Sia $A$ non vuoto e limitato superiormente, e sia $-A=\{-x\colon x\in A\}$. Provare che $-A$ è limitato inferiormente e che $\inf(-A)=-\sup A$.
\end{esercizio}
\risp{Sia $s=\sup A$. (i) $x\le s$ per ogni $x\in A$ equivale a $-x\ge-s$: $-s$ è un minorante di $-A$. (ii) Dato $\eps>0$ esiste $x\in A$ con $x>s-\eps$, cioè $-x<-s+\eps$, con $-x\in-A$. Per la caratterizzazione dell'estremo inferiore, $\inf(-A)=-s$.}

%-------------------------------------------------------------------
\section{Potenze e radici}

\begin{esercizio}[\liv{1}]
Calcolare senza calcolatrice: \ $2^{-3}$, \ $\bigl(\tfrac12\bigr)^{-2}$, \ $10^{0}$, \ $4^{1/2}$, \ $8^{2/3}$.
\end{esercizio}
\risp{$\frac18$; \ $4$; \ $1$; \ $2$; \ $4$.}

\begin{esercizio}[\liv{1}]
Semplificare: \ $a^3a^5$, \ $(a^2)^4$, \ $\dfrac{a^7}{a^3}$ (con $a\ne0$), \ $(2a)^3$.
\end{esercizio}
\risp{$a^8$; \ $a^8$; \ $a^4$; \ $8a^3$.}

\begin{esercizio}[\liv{1}]
Calcolare:
\begin{gather*}
\bigl(2^3\cdot2^{-5}\bigr)^{-1},\qquad 27^{2/3},\qquad 16^{-3/4},\\
\sqrt{50}-\sqrt{18},\qquad \sqrt[3]{-64} .
\end{gather*}
\end{esercizio}
\risp{$4$; \ $9$; \ $\frac18$; \ $5\sqrt2-3\sqrt2=2\sqrt2$; \ $-4$.}

\begin{esercizio}[\liv{2}]
Semplificare, indicando le condizioni:
\begin{gather*}
\sqrt{x^2-2x+1},\qquad \sqrt[3]{(x-2)^3},\\
\sqrt{x^4},\qquad \bigl(x^{1/2}x^{1/3}\bigr)^6 .
\end{gather*}
\end{esercizio}
\risp{$|x-1|$; \ $x-2$; \ $x^2$; \ $x^5$, con $x>0$.}

\begin{esercizio}[\liv{2}]
Risolvere:
\begin{gather*}
x^3=125,\quad x^4=81,\quad x^6=-1,\quad x^3\le-8,\\
x^4\ge16,\quad x^5+1>0,\quad \sqrt[3]{2x-1}<1,\quad \sqrt[4]{x+2}\ge2 .
\end{gather*}
\end{esercizio}
\risp{$5$; \ $\pm3$; \ nessuna; \ $x\le-2$; \ $|x|\ge2$; \ $x>-1$; \ $2x-1<1$, $x<1$; \ $x+2\ge16$, $x\ge14$.}

\begin{esercizio}[\liv{3}]
Confrontare senza calcolatrice $\sqrt[3]{3}$ e $\sqrt2$; poi ordinare $2^{1/2}$, $3^{1/3}$, $5^{1/5}$.
\end{esercizio}
\risp{Elevando alla sesta (potenza crescente sui positivi): $9>8$, quindi $\sqrt[3]3>\sqrt2$. Elevando alla trentesima: $2^{15}=32\,768$, $3^{10}=59\,049$, $5^6=15\,625$, quindi $5^{1/5}<2^{1/2}<3^{1/3}$.}

%-------------------------------------------------------------------
\section{Notazione scientifica}

\begin{esercizio}[\liv{1}]
Scrivere in notazione scientifica: $0{,}000\,32$; \ $45\,000\,000$; \ $0{,}0605$; \ $1{,}0$.
\end{esercizio}
\risp{$3{,}2\cdot10^{-4}$; $4{,}5\cdot10^{7}$; $6{,}05\cdot10^{-2}$; $1{,}0\cdot10^{0}$.}

\begin{esercizio}[\liv{1}]
Calcolare, dando il risultato in notazione scientifica:
\begin{gather*}
(2{,}5\cdot10^{6})(4\cdot10^{-9}),\qquad \frac{6\cdot10^{3}}{1{,}5\cdot10^{-4}},\\
(3\cdot10^{5})^{2},\qquad \sqrt{1{,}6\cdot10^{9}}.
\end{gather*}
\end{esercizio}
\risp{$1\cdot10^{-2}$; \ $4\cdot10^{7}$; \ $9\cdot10^{10}$; \ $\sqrt{16\cdot10^{8}}=4\cdot10^{4}$.}

\begin{esercizio}[\liv{1}]
Indicare il numero di cifre significative di: $0{,}004\,50$; \ $4{,}50\cdot10^{3}$; \ $1{,}200\cdot10^{3}$; \ $0{,}1$. Perch\'e la scrittura ``$1200$'' \`e ambigua?
\end{esercizio}
\risp{$3$; $3$; $4$; $1$.}

\begin{esercizio}[\liv{2}]
Il raggio della Terra \`e circa $6{,}37\cdot10^{6}$ m. Calcolare in notazione scientifica, con tre cifre significative, la superficie $4\pi r^{2}$ e il volume $\frac43\pi r^{3}$ della Terra.
\end{esercizio}
\risp{$5{,}10\cdot10^{14}$ m$^{2}$; \ $1{,}08\cdot10^{21}$ m$^{3}$.}

\begin{esercizio}[\liv{2}]
Un batterio si divide ogni $20$ minuti. Partendo da un solo batterio, e supponendo che tutti si dividano, quanti batteri ci sono dopo $10$ ore? Dare il risultato in notazione scientifica. (Si usi $2^{10}\approx10^{3}$.)
\end{esercizio}
\risp{$2^{30}=(2^{10})^{3}\approx10^{9}$ (il valore esatto \`e $1\,073\,741\,824\approx1{,}07\cdot10^{9}$).}

%-------------------------------------------------------------------
\section{Percentuali}

\begin{esercizio}[\liv{1}]
Calcolare: il $35\,\%$ di $240$; lo $0{,}5\,\%$ di $3000$; la percentuale che $18$ rappresenta rispetto a $150$; il numero di cui $150$ \`e il $60\,\%$.
\end{esercizio}
\risp{$84$; $15$; $12\,\%$; $250$.}

\begin{esercizio}[\liv{1}]
Una popolazione passa da $1250$ a $1100$ individui. Calcolare la variazione assoluta e la variazione relativa.
\end{esercizio}
\risp{$-150$ individui; $-12\,\%$.}

\begin{esercizio}[\liv{2}]
Dopo una diminuzione del $25\,\%$, una popolazione conta $900$ individui. Quanti erano prima della diminuzione?
\end{esercizio}
\risp{$\frac{900}{0{,}75}=1200$ (non $900\cdot1{,}25=1125$).}

\begin{esercizio}[\liv{2}]
La copertura forestale di una provincia passa dal $32\,\%$ al $28\,\%$ del territorio. Esprimere la variazione in punti percentuali e in percentuale.
\end{esercizio}
\risp{$-4$ punti percentuali; $-12{,}5\,\%$.}

%-------------------------------------------------------------------
\section{Problemi di riepilogo}

\begin{esercizio}[\liv{1}]
Un'aula misura $8\times6\times3$ m. Calcolarne il volume in $\mathrm{m^3}$ e in litri, e scrivere il secondo risultato in notazione scientifica.
\end{esercizio}
\risp{$144\ \mathrm{m^3}=144\,000$ L $=1{,}44\cdot10^5$ L.}

\begin{esercizio}[\liv{2}]
Un lago di superficie $3{,}2\ \mathrm{km^2}$ e profondità media $4{,}5$ m viene raccolto, idealmente, in una vasca cubica. Quanto misura il lato? Esprimere il risultato in notazione scientifica con due cifre significative.
\end{esercizio}
\risp{Volume $1{,}44\cdot10^7\ \mathrm{m^3}$; lato $\sqrt[3]{1{,}44\cdot10^7}\ \mathrm m\approx2{,}4\cdot10^2$ m.}

%-------------------------------------------------------------------
\stampasoluzioni
   % Lezioni 4--6
% !TEX root =esercizi.tex
% ---------------------------------------------------------------------
% Scheda 3: Lezioni 7--9 (senza errori di misura, per scelta del docente). In L9 sono state svolte potenze, trasformazioni
% dei grafici, disequazioni irrazionali, funzioni esponenziali e logaritmiche
% e prime disequazioni. Ruffini, funzioni razionali e leggi di scala
% non sono stati svolti: restano negli appunti,
% ma per scelta del docente non ci sono esercizi su questi argomenti.
% ---------------------------------------------------------------------
\chapter{Lezioni 7--9 (7--12 ottobre 2026)}

\begin{vuoto}[Argomenti svolti]{}
\begin{itemize}
\item[\bul] \textbf{Lezione 7 --- Variazioni, logaritmi, piano cartesiano, retta}. Composizione di variazioni percentuali e tasso di incremento. Logaritmi: definizione, proprietà, equazioni esponenziali e logaritmiche, tempo di raddoppio, numero di cifre. Piano cartesiano, distanza e punto medio, circonferenza, regioni del piano. La retta: forme esplicita e implicita, coefficiente angolare; lettura grafica di equazioni e disequazioni. Modelli lineari e affini.
\item[\bul] \textbf{Lezione 8 --- Il concetto di funzione}. Funzione come sottoinsieme di $A\times B$; dominio, codominio, immagine, controimmagine; funzioni iniettive, suriettive, biiettive, e loro conteggio. Zeri, segno, monotonia, massimi e minimi; funzioni pari, dispari, periodiche, definite a tratti. Funzioni composte; funzione inversa.
\item[\bul] \textbf{Lezione 9 --- Potenze, disequazioni irrazionali, esponenziali e logaritmi}. Confronto fra potenze; $x^\alpha$ con $\alpha$ reale. Trasformazioni dei grafici: traslazioni, dilatazioni, simmetrie, valori assoluti. Equazioni e disequazioni irrazionali $\sqrt{A(x)}\gtrless B(x)$. Funzioni esponenziali e logaritmiche: grafici, monotonia, il logaritmo come funzione inversa. Prime disequazioni esponenziali e logaritmiche.
\end{itemize}
\end{vuoto}

Il livello di difficoltà è indicato con \liv{1} (esercizio di base), \liv{2} (esercizio standard), \liv{3} (esercizio più impegnativo).

%-------------------------------------------------------------------
\section{Percentuali}

\begin{esercizio}[\liv{2}]
Una popolazione cresce del $10\,\%$ all'anno per tre anni consecutivi. Qual \`e l'aumento percentuale complessivo?
\end{esercizio}
\risp{$1{,}1^{3}-1=33{,}1\,\%$ (non $30\,\%$).}

\begin{esercizio}[\liv{2}]
Un suolo contiene $0{,}2$ ppm (in massa) di mercurio. Quanti grammi di mercurio sono contenuti in $5$ t di suolo?
\end{esercizio}
\risp{$0{,}2\cdot10^{-6}\cdot5\cdot10^{6}$ g $=1$ g.}

\begin{esercizio}[\liv{2}]
Si vogliono preparare $400$ g di soluzione al $5\,\%$ m/m di cloruro di sodio. Quanti grammi di sale e quanti di acqua servono? Qual \`e la concentrazione se si aggiungono $100$ g di acqua?
\end{esercizio}
\risp{$20$ g di sale e $380$ g di acqua; $\frac{20}{500}=4\,\%$.}

\begin{esercizio}[\liv{3}]
Una specie aumenta del $50\,\%$ in un anno e diminuisce del $50\,\%$ l'anno successivo; un'altra diminuisce prima del $50\,\%$ e poi aumenta del $50\,\%$. Quale delle due si trova in una situazione migliore? Provare che un aumento del $p\,\%$ seguito da una diminuzione del $p\,\%$ produce sempre una diminuzione, e calcolarla.
\end{esercizio}
\risp{La stessa: in entrambi i casi il fattore \`e $1{,}5\cdot0{,}5=0{,}75$. In generale $(1+r)(1-r)-1=-r^{2}<0$ per $r\ne0$: si ha una diminuzione del $\frac{p^{2}}{100}\,\%$ (per $p=50$: del $25\,\%$).}

%-------------------------------------------------------------------
\section{Logaritmi}

\begin{esercizio}[\liv{1}]
Calcolare: \ $\log_2 16$; \ $\log_3 27$; \ $\log_{10}1000$; \ $\log_{10}0{,}01$; \ $\log_5 1$; \ $\log_7 7$.
\end{esercizio}
\risp{$4$; \ $3$; \ $3$; \ $-2$; \ $0$; \ $1$.}

\begin{esercizio}[\liv{1}]
Calcolare: \ $\log_2\frac18$; \ $\log_9 3$; \ $\log_{1/3}9$; \ $2^{\log_2 5}$; \ $\log_{10}4+\log_{10}25$; \ $\log_3 18-\log_3 2$.
\end{esercizio}
\risp{$-3$; \ $\frac12$; \ $-2$; \ $5$; \ $\log_{10}100=2$; \ $\log_3 9=2$.}

\begin{esercizio}[\liv{1}]
Risolvere: \ $2^x=32$; \ $3^x=\frac19$; \ $10^x=0{,}001$; \ $5^x=7$; \ $\log_2x=5$; \ $\log_{10}x=-2$.
\end{esercizio}
\risp{$5$; \ $-2$; \ $-3$; \ $x=\log_57$; \ $x=32$; \ $x=0{,}01$.}

\begin{esercizio}[\liv{2}]
Risolvere: \ $2\cdot3^x+1=19$; \ $4^{x-1}=8$; \ $7^{2x}=5$; \ $\log_3(2x-1)=2$; \ $\log_2x+\log_2(x-2)=3$.
\end{esercizio}
\risp{$3^x=9$, $x=2$; \ $2^{2x-2}=2^3$, $x=\frac52$; \ $x=\frac12\log_75$; \ $2x-1=9$, $x=5$; \ con $x>2$: $x(x-2)=8$, $x=4$ ($x=-2$ non rispetta la condizione).}

\begin{esercizio}[\liv{2}]
Quante cifre ha $2^{64}$, il numero (meno uno) dei chicchi della scacchiera? E $7^{20}$? ($\log_{10}2\approx0{,}3010$, $\log_{10}7\approx0{,}8451$.)
\end{esercizio}
\risp{$64\cdot0{,}3010\approx19{,}26$: $20$ cifre. $20\cdot0{,}8451\approx16{,}90$: $17$ cifre.}

\begin{esercizio}[\liv{3}]
Una coltura batterica raddoppia ogni $3$ ore. Dopo quanto tempo è diventata $1000$ volte più numerosa? E se invece cresce del $25\%$ ogni ora?
\end{esercizio}
\risp{$2^{t/3}=1000$: $t=3\log_21000=\frac{9}{\log_{10}2}\approx29{,}9$ ore. Con $1{,}25^t=1000$: $t=\frac{3}{\log_{10}1{,}25}\approx31{,}0$ ore.}

%-------------------------------------------------------------------
\section{Piano cartesiano e retta}

\begin{esercizio}[\liv{1}]
Dati $A=(1,-2)$ e $B=(4,2)$, calcolare la distanza $AB$, il punto medio del segmento $AB$ e l'equazione della retta per $A$ e $B$, in forma esplicita e implicita.
\end{esercizio}
\risp{$AB=\sqrt{3^2+4^2}=5$; punto medio $(\frac52,0)$; $m=\frac43$, $y=\frac43x-\frac{10}3$, cioè $4x-3y-10=0$.}

\begin{esercizio}[\liv{1}]
Scrivere l'equazione della retta per $P=(2,1)$ parallela alla retta $3x+y-4=0$, e quella della retta per $P$ perpendicolare a essa.
\end{esercizio}
\risp{La retta data ha $m=-3$. Parallela: $y=-3x+7$. Perpendicolare: $m=\frac13$, $y=\frac13x+\frac13$.}

\begin{esercizio}[\liv{2}]
Riconoscere che $x^2+y^2-4x+6y-3=0$ è una circonferenza e trovarne centro e raggio. Il punto $(5,1)$ è interno o esterno?
\end{esercizio}
\risp{Completando i quadrati, $(x-2)^2+(y+3)^2=16$: centro $(2,-3)$, raggio $4$. La distanza di $(5,1)$ dal centro è $\sqrt{9+16}=5>4$: il punto è esterno.}

\begin{esercizio}[\liv{2}]
Disegnare la regione del piano descritta dal sistema $y\le2x+1$, $y\ge0$, $x\le2$, e calcolarne l'area.
\end{esercizio}
\risp{È il triangolo di vertici $(-\frac12,0)$, $(2,0)$, $(2,5)$; area $\frac12\cdot\frac52\cdot5=\frac{25}4$.}

\begin{esercizio}[\liv{2}]
L'acqua gela a $0\,^\circ$C, cioè $32\,^\circ$F, e bolle a $100\,^\circ$C, cioè $212\,^\circ$F. Sapendo che il legame fra le due scale è affine, scrivere $F$ in funzione di $C$, dire che cosa rappresenta il coefficiente angolare (con la sua unità di misura) e trovare la temperatura che ha lo stesso valore nelle due scale. Si tratta di una proporzionalità diretta?
\end{esercizio}
\risp{$F=1{,}8\,C+32$; il coefficiente angolare è $1{,}8\ ^\circ\mathrm F$ per $^\circ$C: un grado Celsius corrisponde a $1{,}8$ gradi Fahrenheit. $C=1{,}8C+32$ dà $C=-40$. Non è una proporzionalità diretta, perché $q=32\ne0$: raddoppiare i gradi Celsius non raddoppia i Fahrenheit.}

\begin{esercizio}[\liv{3}]
Per quali valori di $k$ la retta $y=kx$ interseca la circonferenza di centro $(4,0)$ e raggio $2$?
\end{esercizio}
\risp{Sostituendo in $(x-4)^2+y^2=4$: $(1+k^2)x^2-8x+12=0$, che ha soluzioni reali se e solo se $\frac\Delta4=16-12(1+k^2)\ge0$, cioè $k^2\le\frac13$: $-\frac{1}{\sqrt3}\le k\le\frac1{\sqrt3}$. Per $k=\pm\frac1{\sqrt3}$ la retta è tangente.}

%-------------------------------------------------------------------
\section{Funzioni}

\begin{esercizio}[\liv{1}]
Determinare il dominio di
\[
f(x)=\frac{\sqrt{x-2}}{x-5},\qquad g(x)=\frac1{x^2-4x+3},\qquad h(x)=\sqrt{4-x^2},\qquad k(x)=\sqrt{x^2-x-6}.
\]
\end{esercizio}
\risp{$f$: $[2,5[\,\cup\,]5,+\infty[$; \ $g$: $\R\setminus\{1,3\}$; \ $h$: $[-2,2]$; \ $k$: $]-\infty,-2]\cup[3,+\infty[$.}

\begin{esercizio}[\liv{1}]
Quante sono le funzioni da $A=\{a,b,c\}$ in $B=\{1,2\}$? Quante sono suriettive? Quante iniettive? E da $B$ in $A$?
\end{esercizio}
\risp{Da $A$ in $B$: $2^3=8$ funzioni, di cui $8-2=6$ suriettive (si tolgono le due costanti) e nessuna iniettiva. Da $B$ in $A$: $3^2=9$ funzioni, di cui $3\cdot2=6$ iniettive e nessuna suriettiva.}

\begin{esercizio}[\liv{2}]
Sia $f(x)=x^2-2x$, $x\in\R$. (a) Determinare l'immagine di $f$. (b) Determinare le controimmagini $f^{-1}(\{3\})$ e $f^{-1}([0,3])$. (c) $f$ è iniettiva? (d) Su quale semiretta $f$ è strettamente crescente? Scrivere l'inversa della restrizione di $f$ a quella semiretta.
\end{esercizio}
\risp{(a) $f(x)=(x-1)^2-1$: immagine $[-1,+\infty[$. (b) $x^2-2x=3$: $\{-1,3\}$; $0\le x^2-2x\le3$: $[-1,0]\cup[2,3]$. (c) No: $f(-1)=f(3)$. (d) Su $[1,+\infty[$; da $y=(x-1)^2-1$ con $x\ge1$ si ricava $x=1+\sqrt{y+1}$, quindi l'inversa è $y\mapsto1+\sqrt{1+y}$, definita su $[-1,+\infty[$.}

\begin{esercizio}[\liv{2}]
Dire quali delle seguenti funzioni sono pari, quali dispari, quali né pari né dispari:
\[
x^3-x,\qquad x^2+|x|,\qquad x^2+x,\qquad \frac{x^2+1}{x},\qquad |x-1| .
\]
\end{esercizio}
\risp{Dispari; pari; né l'una né l'altra ($f(1)=2$, $f(-1)=0$); dispari; né l'una né l'altra ($f(1)=0$, $f(-1)=2$).}

\begin{esercizio}[\liv{2}]
Siano $f(x)=x^2$ e $g(x)=\sqrt{x-1}$. Scrivere $g\circ f$ e $f\circ g$, con i rispettivi domini.
\end{esercizio}
\risp{$(g\circ f)(x)=\sqrt{x^2-1}$, definita per $x\le-1$ oppure $x\ge1$. $(f\circ g)(x)=(\sqrt{x-1})^2=x-1$, ma definita solo per $x\ge1$: non è la funzione $x-1$ su tutto $\R$.}

\begin{esercizio}[\liv{3}]
Sia $f(x)=\dfrac{x}{1+|x|}$, $x\in\R$. (a) Dimostrare che $f$ è dispari. (b) Scrivendo $f(x)=1-\frac1{1+x}$ per $x\ge0$, dimostrare che $f$ è strettamente crescente su $[0,+\infty[$, e dedurre che lo è su tutto $\R$. (c) Determinare l'immagine, il suo estremo superiore e inferiore; sono massimo e minimo? (d) Scrivere l'inversa.
\end{esercizio}
\risp{(a) $f(-x)=\frac{-x}{1+|x|}=-f(x)$. (b) Se $0\le x_1<x_2$, allora $1+x_1<1+x_2$, $\frac1{1+x_1}>\frac1{1+x_2}$ e $f(x_1)<f(x_2)$; per la disparità $f$ è strettamente crescente anche su $]-\infty,0]$, e i valori sono negativi a sinistra di $0$ e positivi a destra, quindi lo è su $\R$. (c) Su $[0,+\infty[$ l'immagine è $[0,1[$ (per $y\in[0,1[$ si risolve $1-\frac1{1+x}=y$ con $x=\frac{y}{1-y}$); per disparità l'immagine è $]-1,1[$, con $\sup=1$ e $\inf=-1$, che non sono massimo né minimo. (d) $f^{-1}(y)=\dfrac{y}{1-|y|}$, per $y\in\,]-1,1[$.}

%-------------------------------------------------------------------
\section{Funzioni potenza}

\begin{esercizio}[\liv{1}]
Ordinare, senza calcolatrice, i numeri $x^{-1}$, $x^{1/3}$, $x^{3/2}$, $x^2$ per $x=\frac12$ e per $x=2$.
\end{esercizio}
\risp{Per $x=\frac12$: $x^2<x^{3/2}<x^{1/3}<x^{-1}$ (cioè $0{,}25<0{,}35<0{,}79<2$). Per $x=2$: $x^{-1}<x^{1/3}<x^{3/2}<x^2$ (cioè $0{,}5<1{,}26<2{,}83<4$). Per $x<1$ vince l'esponente minore, per $x>1$ il maggiore.}

\begin{esercizio}[\liv{2}]
Risolvere, ricordando il dominio di ciascuna potenza: \ $x^{3/2}\le8$; \ $x^{-2}>4$; \ $x^{1/3}>2$ (con $x>0$); \ $x^{-1/2}\ge3$.
\end{esercizio}
\risp{$[0,4]$; \ $x^2<\frac14$ con $x\ne0$: $]-\frac12,0[\,\cup\,]0,\frac12[$; \ $x>8$; \ $x>0$ e $\sqrt x\le\frac13$: $]0,\frac19]$.}

%-------------------------------------------------------------------
\section{Equazioni e disequazioni irrazionali}

\begin{vuoto}[Promemoria]{}
Con $A=A(x)$ e $B=B(x)$:
\[
\sqrt A\ge B\iff\begin{cases}B<0\\ A\ge0\end{cases}\ \text{oppure}\ \begin{cases}B\ge0\\ A\ge B^2\end{cases}
\qquad\qquad
\sqrt A\le B\iff\begin{cases}A\ge0\\ B\ge0\\ A\le B^2.\end{cases}
\]
Con le disuguaglianze strette si procede allo stesso modo. Elevare al quadrato è lecito solo quando i due membri sono non negativi.
\end{vuoto}

\begin{esercizio}[\liv{1}]
Risolvere: \ $\sqrt{x-1}=3$; \ $\sqrt{2x+3}=x$; \ $\sqrt{x+5}=x-1$.
\end{esercizio}
\risp{$x=10$; \ con $x\ge0$: $x^2-2x-3=0$, $x=3$ ($x=-1$ va scartata); \ con $x\ge1$: $x^2-3x-4=0$, $x=4$ ($x=-1$ va scartata).}

\begin{esercizio}[\liv{1}]
Risolvere: \ $\sqrt{x+2}<3$; \ $\sqrt{3-x}\ge2$; \ $\sqrt{x^2-1}>0$; \ $\sqrt{x-4}\ge-1$.
\end{esercizio}
\risp{$0\le x+2<9$: $[-2,7[$; \ $3-x\ge4$: $x\le-1$; \ $x^2-1>0$: $x<-1$ oppure $x>1$; \ il secondo membro è negativo, quindi basta la condizione di esistenza: $x\ge4$.}

\begin{esercizio}[\liv{2}]
Risolvere: \ $\sqrt{x+3}\le x+1$; \ $\sqrt{x^2-3x}<x-1$.
\end{esercizio}
\risp{Prima: $x\ge-3$, $x\ge-1$ e $x+3\le x^2+2x+1$, cioè $x^2+x-2\ge0$ ($x\le-2$ o $x\ge1$): $[1,+\infty[$. \ Seconda: $x\le0$ o $x\ge3$; $x>1$; $x^2-3x<x^2-2x+1$, cioè $x>-1$: $[3,+\infty[$.}

\begin{esercizio}[\liv{2}]
Risolvere: \ $\sqrt{x^2-4}\ge x-1$; \ $\sqrt{5-x}>x+1$.
\end{esercizio}
\risp{Prima. Sistema con $x-1<0$: $x<1$ e $x^2\ge4$, cioè $x\le-2$. Sistema con $x-1\ge0$: $x^2-4\ge x^2-2x+1$, cioè $x\ge\frac52$. Soluzione: $]-\infty,-2]\cup[\frac52,+\infty[$. \ Seconda. Sistema con $x+1<0$: $x<-1$ (e $x\le5$). Sistema con $x+1\ge0$: $5-x>x^2+2x+1$, cioè $x^2+3x-4<0$, $-4<x<1$, quindi $-1\le x<1$. Soluzione: $]-\infty,1[$.}

\begin{esercizio}[\liv{2}]
In acque basse la velocità di propagazione di un'onda è $v=\sqrt{gh}$, dove $h$ è la profondità e $g\approx9{,}8$ m/s$^2$. Per quali profondità la velocità è almeno $5$ m/s? Verificare che la formula è dimensionalmente corretta.
\end{esercizio}
\risp{$gh\ge25$: $h\ge\frac{25}{9{,}8}\approx2{,}6$ m. $gh$ ha dimensioni $LT^{-2}\cdot L=L^2T^{-2}$, e la sua radice $LT^{-1}$ è una velocità.}

\begin{esercizio}[\liv{3}]
Risolvere $\sqrt{2x-x^2}\ge x-1$, e interpretare il risultato graficamente, riconoscendo nel grafico di $\sqrt{2x-x^2}$ una semicirconferenza.
\end{esercizio}
\risp{Condizione: $0\le x\le2$. Sistema con $x<1$: $0\le x<1$. Sistema con $x\ge1$: $2x-x^2\ge x^2-2x+1$, cioè $2x^2-4x+1\le0$, $1-\frac{\sqrt2}2\le x\le1+\frac{\sqrt2}2$, quindi $1\le x\le1+\frac{\sqrt2}2$. Soluzione: $[0,1+\frac{\sqrt2}2]$. Graficamente: $y=\sqrt{2x-x^2}$ è la semicirconferenza superiore di $(x-1)^2+y^2=1$, e la retta $y=x-1$ la incontra nel punto di ascissa $1+\frac{\sqrt2}2$ (a $45^\circ$ dal centro); a sinistra di quel punto la semicirconferenza sta sopra la retta.}

\begin{esercizio}[\liv{2}]
Risolvere, scrivendo le condizioni necessarie:
\[
\sqrt{x+3}=x-3,\qquad \sqrt{2x+1}\le3,\qquad \sqrt{x^2-x-2}>x-3 .
\]
\end{esercizio}
\risp{Prima: $x\ge3$ e $x+3=(x-3)^2$, cioè $x^2-7x+6=0$: $x=6$ (la soluzione $x=1$ va scartata). Seconda: $0\le2x+1\le9$: $[-\frac12,4]$. Terza: primo sistema $x^2-x-2\ge0$ e $x<3$, cioè $x\le-1$ oppure $2\le x<3$; secondo sistema $x\ge3$ e $x^2-x-2>(x-3)^2$, cioè $x>\frac{11}5$, quindi $x\ge3$. Unendo: $]-\infty,-1]\cup[2,+\infty[$.}

\begin{esercizio}[\liv{3}]
Risolvere $\sqrt{x}+\sqrt{x-3}>3$ e $\dfrac{|x-1|-3}{\sqrt x-1}>0$.
\end{esercizio}
\risp{Prima: condizione $x\ge3$; si scrive $\sqrt{x-3}>3-\sqrt x$: se $3-\sqrt x<0$ (cioè $x>9$) è vera; altrimenti si eleva: $x-3>9-6\sqrt x+x$, cioè $\sqrt x>2$, $x>4$. Soluzione: $x>4$. Seconda: condizioni $x\ge0$, $x\ne1$. Numeratore positivo per $x>4$ (con $x\ge0$), denominatore positivo per $x>1$. Caso $(+/+)$: $x>4$; caso $(-/-)$: $0\le x<1$. Soluzione: $[0,1[\,\cup\,]4,+\infty[$.}

%-------------------------------------------------------------------
\section{Funzioni esponenziali e logaritmiche}

\begin{esercizio}[\liv{1}]
Disegnare i grafici di $2^x$, $(\frac13)^x$, $\log_3x$, $\log_{1/2}x$, indicando dominio, immagine, monotonia e i punti di coordinate intere più evidenti.
\end{esercizio}
\risp{$2^x$: $\R$, $]0,+\infty[$, crescente, per $(0,1)$ e $(1,2)$. $(\frac13)^x$: $\R$, $]0,+\infty[$, decrescente, per $(0,1)$ e $(-1,3)$. $\log_3x$: $]0,+\infty[$, $\R$, crescente, per $(1,0)$ e $(3,1)$. $\log_{1/2}x$: $]0,+\infty[$, $\R$, decrescente, per $(1,0)$ e $(2,-1)$.}

\begin{esercizio}[\liv{1}]
Determinare il dominio di: \ $\log_2(x-3)$; \ $\log(x^2-1)$; \ $\log_{10}\dfrac{x-1}{x+2}$; \ $\sqrt{2^x-4}$; \ $\dfrac1{\log x}$.
\end{esercizio}
\risp{$x>3$; \ $x<-1$ oppure $x>1$; \ $x<-2$ oppure $x>1$; \ $2^x\ge4$, cioè $x\ge2$; \ $x>0$ e $x\ne1$.}

\begin{esercizio}[\liv{2}]
Senza calcolatrice: (a) stabilire il segno di $\log_20{,}7$, $\log_{0{,}3}0{,}2$, $\log_{1/2}5$, $\log_75$; (b) stabilire quale è maggiore fra $2^{\sqrt3}$ e $2^{1{,}7}$, e fra $0{,}4^{-1}$ e $0{,}4^{-2}$.
\end{esercizio}
\risp{(a) Negativo, positivo, negativo, positivo: con base maggiore di $1$ il logaritmo è positivo se e solo se l'argomento è maggiore di $1$; con base minore di $1$ il contrario. (b) $\sqrt3>1{,}7$ (perché $3>2{,}89$) e la base è maggiore di $1$: $2^{\sqrt3}>2^{1{,}7}$. La base $0{,}4$ è minore di $1$ e $-2<-1$: $0{,}4^{-2}>0{,}4^{-1}$ (infatti $6{,}25>2{,}5$).}

\begin{esercizio}[\liv{2}]
Sia $f(x)=\log(x^2+1)$. Determinare dominio, segno e simmetrie. Su quali intervalli è crescente? Qual è il minimo? (Suggerimento: $f$ è composta di $x^2+1$ e di $\log$.)
\end{esercizio}
\risp{Dominio $\R$; $f(x)\ge0$, con $f(x)=0$ solo per $x=0$; pari. Su $[0,+\infty[$ la funzione $x^2+1$ è strettamente crescente e $\log$ anche, quindi la composta è strettamente crescente; per parità è strettamente decrescente su $]-\infty,0]$. Il minimo è $0$, in $x=0$; l'immagine è $[0,+\infty[$.}

%-------------------------------------------------------------------
\section{Trasformazioni dei grafici}

\begin{esercizio}[\liv{1}]
A partire dal grafico di $x^2$ disegnare i grafici di
\[
(x-3)^2,\qquad x^2-4,\qquad (x+1)^2+2,\qquad 1-x^2,\qquad 2x^2 ,
\]
indicando per ciascuno il vertice e l'immagine.
\end{esercizio}
\risp{Traslato a destra di $3$: vertice $(3,0)$, immagine $[0,+\infty[$. \ Traslato in basso di $4$: $(0,-4)$, $[-4,+\infty[$. \ A sinistra di $1$ e in alto di $2$: $(-1,2)$, $[2,+\infty[$. \ Simmetrico rispetto all'asse $x$ e traslato in alto di $1$: $(0,1)$, $]-\infty,1]$. \ Dilatato in verticale di un fattore $2$: $(0,0)$, $[0,+\infty[$.}

\begin{esercizio}[\liv{1}]
A partire dal grafico di $\sqrt x$ disegnare i grafici di $\sqrt{x+4}$, $\sqrt x-2$, $-\sqrt x$, $\sqrt{-x}$, $\sqrt{3-x}$, indicando dominio e immagine.
\end{esercizio}
\risp{$\sqrt{x+4}$: $[-4,+\infty[$, $[0,+\infty[$. \ $\sqrt x-2$: $[0,+\infty[$, $[-2,+\infty[$. \ $-\sqrt x$: $[0,+\infty[$, $]-\infty,0]$. \ $\sqrt{-x}$: $]-\infty,0]$, $[0,+\infty[$. \ $\sqrt{3-x}=\sqrt{-(x-3)}$: il grafico di $\sqrt{-x}$ traslato a destra di $3$; dominio $]-\infty,3]$, immagine $[0,+\infty[$.}

\begin{esercizio}[\liv{2}]
Disegnare i grafici di $|x^2-4|$, \ $x^2-4|x|$, \ $\bigl||x|-2\bigr|$, e dire per ciascuno se è pari.
\end{esercizio}
\risp{$|x^2-4|$: la parabola $x^2-4$ con l'arco fra $-2$ e $2$ ribaltato sopra l'asse; pari. \ $x^2-4|x|$ è $f(|x|)$ con $f(x)=x^2-4x=(x-2)^2-4$: si disegna la parte con $x\ge0$ (vertice $(2,-4)$) e la si ripete simmetricamente; pari. \ $|x|-2$ è una ``V'' con vertice $(0,-2)$; il valore assoluto ribalta la parte fra $-2$ e $2$: una ``W'' con zeri in $\pm2$ e massimo relativo $2$ in $0$; pari.}

\begin{esercizio}[\liv{2}]
A partire dai grafici di $2^x$ e $\log x$ disegnare i grafici di
\[
2^{|x|},\qquad |2^x-2|,\qquad \log(-x),\qquad \log(x+1)-1,\qquad |\log x| ,
\]
indicando dominio, immagine ed eventuali asintoti.
\end{esercizio}
\risp{$2^{|x|}$: $\R$, $[1,+\infty[$, pari, minimo in $0$. \ $|2^x-2|$: $\R$, $[0,+\infty[$, zero in $x=1$, asintoto orizzontale $y=2$ a sinistra. \ $\log(-x)$: $]-\infty,0[$, $\R$, simmetrico di $\log x$ rispetto all'asse $y$. \ $\log(x+1)-1$: $]-1,+\infty[$, $\R$, asintoto verticale $x=-1$, zero in $x=e-1$. \ $|\log x|$: $]0,+\infty[$, $[0,+\infty[$, zero in $x=1$.}

\begin{esercizio}[\liv{3}]
(a) Scrivere $x^2-6x+5$ nella forma $(x-h)^2+k$ e disegnare il grafico di $|x^2-6x+5|$. (b) Quante soluzioni ha l'equazione $|x^2-6x+5|=3$? Leggerlo dal grafico e poi verificarlo risolvendo l'equazione.
\end{esercizio}
\risp{(a) $(x-3)^2-4$: parabola di vertice $(3,-4)$ e zeri $1$ e $5$; l'arco fra $1$ e $5$ si ribalta, con massimo relativo $4$ in $x=3$. (b) La retta $y=3$ sta sotto il massimo relativo $4$, quindi taglia il grafico in $4$ punti. Infatti $(x-3)^2-4=3$ dà $x=3\pm\sqrt7$, e $(x-3)^2-4=-3$ dà $x=2$ e $x=4$.}

\begin{esercizio}[\liv{3}]
Al variare di $k\in\R$, quante soluzioni ha l'equazione $|x^2-4|=k$? (Ragionare sul grafico del primo membro e sulla retta orizzontale $y=k$.)
\end{esercizio}
\risp{Nessuna per $k<0$; due ($x=\pm2$) per $k=0$; quattro per $0<k<4$; tre per $k=4$ ($x=0$ e $x=\pm2\sqrt2$); due per $k>4$.}

%-------------------------------------------------------------------
\section{Disequazioni esponenziali e logaritmiche}

\begin{esercizio}[\liv{1}]
Risolvere: \ $2^x>16$; \ $3^{x-1}\le\frac19$; \ $(\frac12)^x<8$; \ $10^x\ge0{,}01$; \ $e^x>0$; \ $5^x<0$.
\end{esercizio}
\risp{$x>4$; \ $x-1\le-2$: $x\le-1$; \ $(\frac12)^x<(\frac12)^{-3}$, base minore di $1$: $x>-3$; \ $x\ge-2$; \ ogni $x\in\R$; \ nessuna soluzione.}

\begin{esercizio}[\liv{1}]
Risolvere: \ $\log_2x>3$; \ $\log_3x\le-1$; \ $\log_{1/2}x\ge2$; \ $\log x<0$.
\end{esercizio}
\risp{$x>8$; \ $0<x\le\frac13$; \ base minore di $1$: $0<x\le\frac14$; \ $0<x<1$.}

\begin{esercizio}[\liv{2}]
Risolvere: \ $2^{x^2-3x}<\frac14$; \ $0{,}5^{2x+1}>0{,}5^{x-2}$; \ $4^x-2^{x+1}-8\ge0$.
\end{esercizio}
\risp{$x^2-3x<-2$: $1<x<2$; \ base minore di $1$: $2x+1<x-2$, $x<-3$; \ con $t=2^x>0$: $t^2-2t-8\ge0$, cioè $t\ge4$ (la soluzione $t\le-2$ è incompatibile con $t>0$): $x\ge2$.}

\begin{esercizio}[\liv{2}]
Risolvere, scrivendo prima le condizioni di esistenza:
\begin{gather*}
\log_2(3x-1)\le3,\qquad \log_{1/3}(x-1)>-2,\\
\log(x^2-3)>0,\qquad \log_{10}(x+1)+\log_{10}(x-1)<\log_{10}3 .
\end{gather*}
\end{esercizio}
\risp{$x>\frac13$ e $3x-1\le8$: $]\frac13,3]$; \ $x>1$ e, con base minore di $1$, $x-1<9$: $]1,10[$; \ $x^2-3>1$, cioè $x<-2$ oppure $x>2$ (la condizione $x^2>3$ è compresa); \ $x>1$ e $x^2-1<3$: $]1,2[$.}

\begin{esercizio}[\liv{3}]
Risolvere: \ $(\log x)^2-\log x-2<0$; \ $2^x>3^{x-1}$.
\end{esercizio}
\risp{Con $t=\log x$ ($x>0$): $t^2-t-2<0$, cioè $-1<t<2$: $e^{-1}<x<e^2$. \ Prendendo $\log$ (base maggiore di $1$, verso conservato): $x\log2>(x-1)\log3$, cioè $x(\log3-\log2)<\log3$: $x<\dfrac{\log3}{\log\frac32}\approx2{,}71$. (Verifica: per $x=2$, $4>3$; per $x=3$, $8<9$.)}

%-------------------------------------------------------------------
\section{Misure e ordini di grandezza}

Esercizi su argomenti di L6 non ancora svolti a lezione: unità di misura e omogeneità dimensionale, ordini di grandezza.

\begin{esercizio}[\liv{1}]
Convertire: $3{,}2$ km in m; $450$ mL in L; $2{,}5$ h in s; $72$ km/h in m/s; $1{,}5$ m$^{2}$ in cm$^{2}$; $0{,}8$ m$^{3}$ in L.
\end{esercizio}
\risp{$3200$ m; $0{,}45$ L; $9000$ s; $20$ m/s; $15\,000$ cm$^{2}$; $800$ L.}

\begin{esercizio}[\liv{1}]
La densit\`a dell'acqua di mare \`e circa $1{,}03$ g/cm$^{3}$. Esprimerla in kg/m$^{3}$.
\end{esercizio}
\risp{$1030$ kg/m$^{3}$.}

\begin{esercizio}[\liv{2}]
Su un bacino di $25$ km$^{2}$ cadono $12$ mm di pioggia. Quanti metri cubi d'acqua sono caduti sul bacino?
\end{esercizio}
\risp{$0{,}012\ \text{m}\cdot25\cdot10^{6}\ \text{m}^{2}=3\cdot10^{5}$ m$^{3}$.}

\begin{esercizio}[\liv{2}]
Un fiume ha una portata di $3{,}6$ m$^{3}$/s. Quanti litri d'acqua passano in un'ora?
\end{esercizio}
\risp{$3{,}6\cdot3600=12\,960$ m$^{3}\approx1{,}30\cdot10^{7}$ L.}

\begin{esercizio}[\liv{2}]
La resa di un campo di frumento \`e di $5{,}2$ t/ha. Esprimerla in kg/m$^{2}$ ($1$ ha $=10^{4}$ m$^{2}$).
\end{esercizio}
\risp{$0{,}52$ kg/m$^{2}$.}

\begin{esercizio}[\liv{2}]
Siano $v$ una velocit\`a, $a$ un'accelerazione, $t$ un tempo, $m$ una massa, $s$ e $h$ lunghezze, $A$ un'area, $g$ l'accelerazione di gravit\`a, $Q$ una portata volumetrica. Quali delle seguenti formule sono dimensionalmente corrette?
\begin{gather*}
v=at,\qquad s=vt^{2},\qquad Q=Av,\\
h=\frac{v^{2}}{2g},\qquad s=\frac{m\,v}{t}.
\end{gather*}
\end{esercizio}
\risp{Sono corrette la prima, la terza e la quarta.}

\begin{esercizio}[\liv{2}]
Nella formula $x=x_{0}+v\,t+k\,t^{2}$, dove $x$, $x_{0}$ sono lunghezze, $v$ una velocit\`a e $t$ un tempo, quali dimensioni deve avere $k$?
\end{esercizio}
\risp{$LT^{-2}$.}

\begin{esercizio}[\liv{2}]
In un lago la concentrazione di ossigeno disciolto \`e $8$ mg/L. Esprimerla in mol/L, sapendo che la massa molare di O$_{2}$ \`e $32$ g/mol.
\end{esercizio}
\risp{$\frac{8\cdot10^{-3}}{32}=2{,}5\cdot10^{-4}$ mol/L.}

\begin{esercizio}[\liv{2}]
Il corpo umano contiene circa $3{,}7\cdot10^{13}$ cellule. Assumendo una massa media di $1$ ng per cellula, stimare la massa complessiva delle cellule. Il risultato \`e plausibile?
\end{esercizio}
\risp{$3{,}7\cdot10^{13}\cdot10^{-9}$ g $=3{,}7\cdot10^{4}$ g $=37$ kg: lo stesso ordine di grandezza della massa di una persona.}

\begin{esercizio}[\liv{3}]
Una citt\`a ha $10^{6}$ abitanti, ciascuno dei quali consuma in media $200$ litri d'acqua al giorno. Stimare il consumo annuo della citt\`a in litri e in metri cubi, e confrontarlo con il volume di un lago di $5$ km$^{2}$ di superficie e $20$ m di profondit\`a media.
\end{esercizio}
\risp{$2\cdot10^{2}\cdot10^{6}\cdot365\approx7{,}3\cdot10^{10}$ L $=7{,}3\cdot10^{7}$ m$^{3}$; il lago contiene $5\cdot10^{6}\cdot20=10^{8}$ m$^{3}$: il consumo annuo \`e circa il $73\,\%$ del volume del lago.}

%-------------------------------------------------------------------
\section{Problemi di riepilogo}

\begin{esercizio}[\liv{2}]
Un lago ha un volume di $2{,}4\cdot10^{7}$ m$^{3}$ e un inquinante vi \`e presente con concentrazione media di $0{,}05$ mg/L. Quanti chilogrammi di inquinante contiene il lago? Se la concentrazione diminuisce del $20\,\%$ all'anno, quale sar\`a la concentrazione dopo tre anni, e di quanto sar\`a diminuita in percentuale?
\end{esercizio}
\risp{$0{,}05\ \text{mg/L}\cdot2{,}4\cdot10^{10}$ L $=1{,}2\cdot10^{9}$ mg $=1{,}2\cdot10^{3}$ kg. Dopo tre anni: $0{,}05\cdot0{,}8^{3}=0{,}0256$ mg/L, con una diminuzione del $48{,}8\,\%$.}

\begin{esercizio}[\liv{2}]
Ogni anno si degrada il $30\,\%$ della quantità di un pesticida presente nel suolo. Dopo quanti anni ne resta meno del $10\,\%$ della quantità iniziale?
\end{esercizio}
\risp{$0{,}7^n<0{,}1$: $n>\dfrac{\log0{,}1}{\log0{,}7}\approx6{,}46$, cioè dopo $7$ anni.}

\begin{esercizio}[\liv{2}]
Una coltura batterica segue il modello $N(t)=500\cdot2^{t/4}$, con $t$ in ore. Qual è il tempo di raddoppio? Dopo quanto tempo la coltura supera i $10^5$ individui?
\end{esercizio}
\risp{$4$ ore. $2^{t/4}>200$: $t>4\log_2200\approx30{,}6$ ore.}

\begin{esercizio}[\liv{3}]
Due modelli per la stessa popolazione, con $t$ in anni: $A(t)=1000+150\,t$ e $B(t)=1000\cdot1{,}1^t$. Che tipo di crescita descrive ciascuno? Quale prevede la popolazione maggiore dopo $5$ anni? Dopo quanti anni interi il modello $B$ supera il modello $A$? (L'equazione $A(t)=B(t)$ non si risolve con i metodi del corso: procedere per tentativi, aiutandosi con un grafico.)
\end{esercizio}
\risp{$A$ è affine (aumento assoluto costante di $150$ individui all'anno), $B$ esponenziale (aumento relativo costante del $10\,\%$). Dopo $5$ anni $A=1750$, $B\approx1611$. Si ha $B(8)\approx2144<A(8)=2200$ e $B(9)\approx2358>A(9)=2350$: dal nono anno $B$ supera $A$, e da lì in poi la differenza cresce sempre.}

%-------------------------------------------------------------------
\stampasoluzioni

   % Lezioni 7--9
% !TEX root =esercizi.tex
% ---------------------------------------------------------------------
% Scheda 4: Lezioni 10--12. In L10 esponenziali e logaritmi come funzioni
% (completati) e funzioni trigonometriche; in L11 equazioni e disequazioni
% trigonometriche, inverse, campi di esistenza di riepilogo; in L12
% successioni, Fibonacci e rapporto aureo, successioni divergenti.
% Non svolti in L12 (e quindi senza esercizi qui): ricorrenza affine,
% monotonia e limitatezza (riprese in L13).
% ---------------------------------------------------------------------
\chapter{Lezioni 10--12 (14--19 ottobre 2026)}

\begin{vuoto}[Argomenti svolti]{}
\begin{itemize}
\item[\bul] \textbf{Lezione 10 --- Esponenziali e logaritmi; funzioni trigonometriche}. Funzioni esponenziali e logaritmiche: grafici, monotonia, crescita e decadimento, tempo di raddoppio e di dimezzamento, il numero $e$ e il logaritmo naturale, $N(t)=N_0e^{kt}$. Radianti; seno, coseno, tangente sulla circonferenza goniometrica; grafici.
\item[\bul] \textbf{Lezione 11 --- Trigonometria}. Valori notevoli e simmetrie; equazioni e disequazioni trigonometriche elementari, lette sulla circonferenza e sul grafico; funzioni trigonometriche inverse; modelli periodici $m+A\sin(\omega t+\varphi)$. Campi di esistenza di funzioni composte, come riepilogo di tutte le disequazioni.
\item[\bul] \textbf{Lezione 12 --- Successioni}. Successioni come funzioni su $\N$: elenco, formula esplicita, legge ricorsiva, grafico. La successione di Fibonacci, i rapporti fra termini consecutivi, il rapporto aureo e la sezione aurea. Successioni divergenti e confronto delle loro velocità.
\end{itemize}
\end{vuoto}

Il livello di difficoltà è indicato con \liv{1} (esercizio di base), \liv{2} (esercizio standard), \liv{3} (esercizio più impegnativo).

%-------------------------------------------------------------------
\section{Crescita e decadimento esponenziale}

\begin{esercizio}[\liv{1}]
Lo iodio-131, usato in medicina nucleare, ha un tempo di dimezzamento di circa $8$ giorni. Partendo da $100$ mg, quanto ne resta dopo $24$ giorni? E dopo $30$ giorni? Dopo quanto tempo ne resta $1$ mg?
\end{esercizio}
\risp{$N(t)=100\cdot2^{-t/8}$. Dopo $24$ giorni (tre dimezzamenti) $12{,}5$ mg; dopo $30$ giorni $100\cdot2^{-3{,}75}\approx7{,}4$ mg. $2^{-t/8}=0{,}01$ dà $t=8\log_2100\approx53$ giorni.}

\begin{esercizio}[\liv{2}]
Una popolazione cresce del $3\,\%$ all'anno. Calcolare il tempo di raddoppio esatto e quello dato dalla regola del $70$. Scrivere il modello nella forma $N_0e^{kt}$: quanto vale il tasso istantaneo $k$?
\end{esercizio}
\risp{$T=\dfrac{\log2}{\log1{,}03}\approx23{,}4$ anni; regola del $70$: $\frac{70}3\approx23{,}3$ anni. $k=\log1{,}03\approx0{,}0296$, cioè circa $2{,}96\,\%$, un po' meno del $3\,\%$.}

\begin{esercizio}[\liv{2}]
Il carbonio-14 ha tempo di dimezzamento di $5730$ anni. Un reperto contiene il $40\,\%$ del carbonio-14 che conteneva in origine. Quanti anni ha?
\end{esercizio}
\risp{$2^{-t/5730}=0{,}4$: $t=5730\log_2\frac1{0{,}4}=5730\,\dfrac{\log2{,}5}{\log2}\approx7{,}6\cdot10^3$ anni.}

\begin{esercizio}[\liv{2}]
Scrivere $N(t)=200\cdot0{,}85^{t}$ nella forma $N_0e^{kt}$. Si tratta di crescita o di decadimento? Qual è il tempo di dimezzamento?
\end{esercizio}
\risp{$k=\log0{,}85\approx-0{,}163$: decadimento, $N(t)=200\,e^{-0{,}163\,t}$. Il tempo di dimezzamento è $\dfrac{\log2}{-\log0{,}85}\approx4{,}27$.}

\begin{esercizio}[\liv{2}]
Il pH di una soluzione è $-\log_{10}[\mathrm H^+]$, con la concentrazione in mol/L. Calcolare il pH di una soluzione con $[\mathrm H^+]=3{,}2\cdot10^{-5}$ mol/L. Quante volte è più grande la concentrazione di ioni idrogeno in una soluzione a pH $4$ rispetto a una a pH $6$?
\end{esercizio}
\risp{$\mathrm{pH}=5-\log_{10}3{,}2\approx4{,}49$. A pH $4$ la concentrazione è $10^{-4}$, a pH $6$ è $10^{-6}$: è $100$ volte più grande.}

\begin{esercizio}[\liv{3}]
Due colture batteriche, con $t$ in ore: $A(t)=1000\cdot2^{t/3}$ e $B(t)=8000\cdot2^{t/5}$. Quale raddoppia più in fretta? Da quale istante in poi $A$ supera $B$?
\end{esercizio}
\risp{$A$ raddoppia in $3$ ore, $B$ in $5$. $A(t)>B(t)$ equivale a $2^{t/3-t/5}>8=2^3$, cioè $\frac{2t}{15}>3$: $t>22{,}5$ ore.}

%-------------------------------------------------------------------
\section{Radianti e valori delle funzioni trigonometriche}

\begin{esercizio}[\liv{1}]
Esprimere in radianti $30^\circ$, $135^\circ$, $270^\circ$, $-60^\circ$, e in gradi $\frac{5\pi}6$, $\frac{7\pi}4$, $2$ rad.
\end{esercizio}
\risp{$\frac\pi6$, $\frac{3\pi}4$, $\frac{3\pi}2$, $-\frac\pi3$; \ $150^\circ$, $315^\circ$, $\frac{360^\circ}{\pi}\approx114{,}6^\circ$.}

\begin{esercizio}[\liv{1}]
Calcolare: \ $\sin\frac{5\pi}6$; \ $\cos\frac{4\pi}3$; \ $\tan\frac{3\pi}4$; \ $\sin\bigl(-\frac\pi4\bigr)$; \ $\cos\frac{11\pi}6$; \ $\tan\frac{7\pi}6$; \ $\sin\frac{13\pi}3$.
\end{esercizio}
\risp{$\frac12$; \ $-\frac12$; \ $-1$; \ $-\frac{\sqrt2}2$; \ $\frac{\sqrt3}2$; \ $\frac{\sqrt3}3$; \ $\frac{13\pi}3=4\pi+\frac\pi3$, quindi $\frac{\sqrt3}2$.}

\begin{esercizio}[\liv{2}]
Sapendo che $\sin x=\frac35$ e $\frac\pi2<x<\pi$, calcolare $\cos x$, $\tan x$, $\sin(\pi-x)$, $\cos(x+\pi)$.
\end{esercizio}
\risp{Nel secondo quadrante il coseno è negativo: $\cos x=-\sqrt{1-\frac9{25}}=-\frac45$; $\tan x=-\frac34$; $\sin(\pi-x)=\sin x=\frac35$; $\cos(x+\pi)=-\cos x=\frac45$.}

\begin{esercizio}[\liv{2}]
Il raggio della Terra è circa $6370$ km. Quanto è lungo un arco di meridiano di $1^\circ$? Napoli si trova a una latitudine di circa $40{,}85^\circ$ nord: quanto dista dall'equatore, misurando lungo il meridiano?
\end{esercizio}
\risp{L'arco è lungo $r\theta$ con $\theta$ in radianti: $6370\cdot\frac\pi{180}\approx111$ km; $6370\cdot40{,}85\cdot\frac\pi{180}\approx4540$ km.}

%-------------------------------------------------------------------
\section{Grafici e modelli periodici}

\begin{esercizio}[\liv{1}]
Determinare periodo, ampiezza e immagine di: \ $3\sin(2x)$; \ $1+\cos\frac x2$; \ $1-2\sin(\pi x)$.
\end{esercizio}
\risp{Periodo $\pi$, ampiezza $3$, immagine $[-3,3]$; \ periodo $4\pi$, ampiezza $1$, immagine $[0,2]$; \ periodo $2$, ampiezza $2$, immagine $[-1,3]$.}

\begin{esercizio}[\liv{2}]
A partire dal grafico di $\sin x$ disegnare i grafici di $\sin\bigl(x-\frac\pi4\bigr)$, $2\cos x-1$, $|\sin x|$, $\sin|x|$. Quali sono periodici, e con quale periodo?
\end{esercizio}
\risp{$\sin(x-\frac\pi4)$: traslato a destra di $\frac\pi4$, periodo $2\pi$. $2\cos x-1$: dilatato verticalmente di $2$ e abbassato di $1$, periodo $2\pi$, immagine $[-3,1]$. $|\sin x|$: le parti negative ribaltate, periodo $\pi$. $\sin|x|$: per $x\ge0$ coincide con $\sin x$, ed è pari; non è periodica.}

\begin{esercizio}[\liv{2}]
In un porto l'altezza della marea, in metri, è descritta da $h(t)=1{,}2+0{,}8\sin\dfrac{2\pi t}{12{,}4}$, con $t$ in ore. Determinare altezza massima e minima e il periodo. In quale intervallo di tempo, fra $t=0$ e $t=12{,}4$, l'altezza supera $1{,}6$ m? Per quante ore?
\end{esercizio}
\risp{Massimo $2$ m, minimo $0{,}4$ m, periodo $12{,}4$ ore. $h>1{,}6$ equivale a $\sin\frac{2\pi t}{12{,}4}>\frac12$, cioè $\frac\pi6<\frac{2\pi t}{12{,}4}<\frac{5\pi}6$: $\frac{12{,}4}{12}<t<\frac{5\cdot12{,}4}{12}$, cioè $1{,}03<t<5{,}17$, per circa $4{,}1$ ore.}

\begin{esercizio}[\liv{3}]
In una località la temperatura media giornaliera ha un massimo di $26\,^\circ$C il giorno $200$ dell'anno e un minimo di $6\,^\circ$C. Scrivere un modello $T(t)=m+A\sin\bigl(\frac{2\pi}{365}(t-t_0)\bigr)$ e dire in quale giorno si ha il minimo.
\end{esercizio}
\risp{$m=\frac{26+6}2=16$, $A=\frac{26-6}2=10$. Il seno vale $1$ quando l'argomento è $\frac\pi2$, cioè $t-t_0=\frac{365}4$: $t_0=200-91{,}25=108{,}75$. $T(t)=16+10\sin\bigl(\frac{2\pi}{365}(t-108{,}75)\bigr)$. Il minimo è mezzo periodo dopo il massimo: $200+182{,}5-365=17{,}5$, cioè verso il 17--18 gennaio.}

%-------------------------------------------------------------------
\section{Equazioni trigonometriche}

\begin{esercizio}[\liv{1}]
Risolvere in $\R$: \ $\sin x=\frac{\sqrt3}2$; \ $\cos x=-\frac12$; \ $\tan x=1$; \ $\sin x=2$; \ $\cos x=0$.
\end{esercizio}
\risp{$x=\frac\pi3+2k\pi$ oppure $x=\frac{2\pi}3+2k\pi$; \ $x=\pm\frac{2\pi}3+2k\pi$; \ $x=\frac\pi4+k\pi$; \ nessuna soluzione; \ $x=\frac\pi2+k\pi$ ($k\in\Z$).}

\begin{esercizio}[\liv{2}]
Risolvere in $[0,2\pi]$: \ $\sin2x=\frac12$; \ $\cos\bigl(x-\frac\pi3\bigr)=1$; \ $2\sin^2x-\sin x-1=0$.
\end{esercizio}
\risp{$2x=\frac\pi6+2k\pi$ oppure $2x=\frac{5\pi}6+2k\pi$: $x=\frac\pi{12},\frac{5\pi}{12},\frac{13\pi}{12},\frac{17\pi}{12}$. \ $x-\frac\pi3=2k\pi$: $x=\frac\pi3$. \ Con $t=\sin x$: $2t^2-t-1=0$, $t=1$ o $t=-\frac12$: $x=\frac\pi2,\frac{7\pi}6,\frac{11\pi}6$.}

\begin{esercizio}[\liv{2}]
Risolvere in $\R$: \ $\sin x=\cos x$; \ $2\cos^2x+3\sin x-3=0$.
\end{esercizio}
\risp{Se $\cos x=0$ allora $\sin x=\pm1\ne0$, quindi si può dividere: $\tan x=1$, $x=\frac\pi4+k\pi$. \ Con $\cos^2x=1-\sin^2x$: $2\sin^2x-3\sin x+1=0$, $\sin x=1$ o $\sin x=\frac12$: $x=\frac\pi2+2k\pi$, $x=\frac\pi6+2k\pi$, $x=\frac{5\pi}6+2k\pi$.}

\begin{esercizio}[\liv{2}]
Risolvere $\sin x=0{,}3$ in $[0,2\pi]$, usando la funzione $\arcsin$, e dare i valori approssimati.
\end{esercizio}
\risp{$x=\arcsin0{,}3\approx0{,}305$ e $x=\pi-\arcsin0{,}3\approx2{,}837$.}

%-------------------------------------------------------------------
\section{Disequazioni trigonometriche}

\begin{esercizio}[\liv{1}]
Risolvere in $[0,2\pi]$: \ $\sin x\ge\frac12$; \ $\cos x<\frac{\sqrt2}2$; \ $\tan x>1$.
\end{esercizio}
\risp{$[\frac\pi6,\frac{5\pi}6]$; \ $]\frac\pi4,\frac{7\pi}4[$; \ $]\frac\pi4,\frac\pi2[\,\cup\,]\frac{5\pi}4,\frac{3\pi}2[$.}

\begin{esercizio}[\liv{2}]
Risolvere in $[0,2\pi]$: \ $2\cos x+1\le0$; \ $\sin^2x<\frac14$.
\end{esercizio}
\risp{$\cos x\le-\frac12$: $[\frac{2\pi}3,\frac{4\pi}3]$. \ $|\sin x|<\frac12$: $[0,\frac\pi6[\,\cup\,]\frac{5\pi}6,\frac{7\pi}6[\,\cup\,]\frac{11\pi}6,2\pi]$.}

\begin{esercizio}[\liv{2}]
Risolvere in $[0,2\pi]$ la disequazione $\sin2x>0$, ponendo $t=2x$ (attenzione all'intervallo in cui varia $t$).
\end{esercizio}
\risp{$t\in[0,4\pi]$ e $\sin t>0$ per $t\in\,]0,\pi[\,\cup\,]2\pi,3\pi[$: $x\in\,]0,\frac\pi2[\,\cup\,]\pi,\frac{3\pi}2[$.}

\begin{esercizio}[\liv{3}]
Risolvere in $[0,2\pi]$: \ $2\sin^2x-3\sin x+1>0$.
\end{esercizio}
\risp{Con $t=\sin x$: $(2t-1)(t-1)>0$, cioè $t<\frac12$ oppure $t>1$; la seconda è impossibile. Quindi $\sin x<\frac12$: $[0,\frac\pi6[\,\cup\,]\frac{5\pi}6,2\pi]$.}

%-------------------------------------------------------------------
\section{Funzioni trigonometriche inverse}

\begin{esercizio}[\liv{1}]
Calcolare: \ $\arcsin\frac12$; \ $\arccos\bigl(-\frac{\sqrt3}2\bigr)$; \ $\arctan(-1)$; \ $\arcsin(-1)$; \ $\arccos0$; \ $\arctan\sqrt3$.
\end{esercizio}
\risp{$\frac\pi6$; \ $\frac{5\pi}6$; \ $-\frac\pi4$; \ $-\frac\pi2$; \ $\frac\pi2$; \ $\frac\pi3$.}

\begin{esercizio}[\liv{2}]
Calcolare:
\[
\arcsin\Bigl(\sin\frac{3\pi}4\Bigr),\quad \arccos\Bigl(\cos\Bigl(-\frac\pi6\Bigr)\Bigr),\quad \arctan\Bigl(\tan\frac{5\pi}4\Bigr),\quad \sin\Bigl(\arccos\frac13\Bigr),\quad \tan\Bigl(\arcsin\frac35\Bigr).
\]
\end{esercizio}
\risp{$\frac\pi4$ (non $\frac{3\pi}4$, che non sta in $[-\frac\pi2,\frac\pi2]$); \ $\frac\pi6$; \ $\frac\pi4$; \ $\sqrt{1-\frac19}=\frac{2\sqrt2}3$ (positivo, perché $\arccos\frac13\in[0,\pi]$); \ con $\sin y=\frac35$ e $y\in[-\frac\pi2,\frac\pi2]$ si ha $\cos y=\frac45$, quindi $\frac34$.}

\begin{esercizio}[\liv{2}]
Determinare il dominio di: \ $\arcsin(2x-1)$; \ $\arccos\frac1x$; \ $\arcsin(x^2-1)$.
\end{esercizio}
\risp{$-1\le2x-1\le1$: $[0,1]$; \ $\bigl|\frac1x\bigr|\le1$: $|x|\ge1$; \ $0\le x^2\le2$: $[-\sqrt2,\sqrt2]$.}

\begin{esercizio}[\liv{2}]
Risolvere: \ $\arctan x>\frac\pi3$; \ $\arcsin x\le\frac\pi6$; \ $\arccos x<\frac\pi4$.
\end{esercizio}
\risp{$\arctan$ è crescente: $x>\sqrt3$. \ $\arcsin$ è crescente su $[-1,1]$: $[-1,\frac12]$. \ $\arccos$ è decrescente: $x>\frac{\sqrt2}2$, con $x\le1$: $]\frac{\sqrt2}2,1]$.}

%-------------------------------------------------------------------
\section{Campi di esistenza}

\begin{vuoto}[Promemoria]{}
Per determinare il dominio di una funzione composta si scrivono, e si mettono a sistema, le condizioni di ciascun ``pezzo'':
\begin{itemize}
\item[\bul] denominatori diversi da $0$;
\item[\bul] argomenti delle radici di indice pari maggiori o uguali a $0$;
\item[\bul] argomenti dei logaritmi strettamente positivi;
\item[\bul] argomenti di $\arcsin$ e $\arccos$ compresi fra $-1$ e $1$;
\item[\bul] argomenti della tangente diversi da $\frac\pi2+k\pi$.
\end{itemize}
\end{vuoto}

\begin{esercizio}[\liv{1}]
Determinare il dominio di: \ $\sqrt{x^2-5x+6}$; \ $\log(4-x^2)$; \ $\dfrac1{e^x-1}$; \ $\sqrt{\log x}$; \ $\dfrac{x+1}{\sqrt{9-x^2}}$.
\end{esercizio}
\risp{$x\le2$ oppure $x\ge3$; \ $]-2,2[$; \ $x\ne0$; \ $\log x\ge0$: $x\ge1$; \ $9-x^2>0$: $]-3,3[$.}

\begin{esercizio}[\liv{2}]
Determinare il dominio di: \ $\sqrt{\dfrac{x-1}{x+3}}$; \ $\log(|x|-2)$; \ $\sqrt{1-\log_2x}$; \ $\log_{1/3}(x^2-x)$; \ $\dfrac{\sqrt{2^x-8}}{x-5}$.
\end{esercizio}
\risp{$x<-3$ oppure $x\ge1$; \ $|x|>2$: $x<-2$ oppure $x>2$; \ $x>0$ e $\log_2x\le1$: $]0,2]$; \ $x^2-x>0$: $x<0$ oppure $x>1$; \ $2^x\ge8$ e $x\ne5$: $[3,5[\,\cup\,]5,+\infty[$.}

\begin{esercizio}[\liv{2}]
Determinare il dominio di:
\[
\sqrt{\log(x-2)},\qquad \log(\log x),\qquad \sqrt{\log_{1/2}(x+1)},\qquad \frac1{\sqrt{x+4}-2},\qquad \sqrt{4^x-3\cdot2^x+2}.
\]
\end{esercizio}
\risp{$x>2$ e $\log(x-2)\ge0$, cioè $x-2\ge1$: $[3,+\infty[$; \ $x>0$ e $\log x>0$: $x>1$; \ $x>-1$ e, con base minore di $1$, $x+1\le1$: $]-1,0]$; \ $x\ge-4$ e $\sqrt{x+4}\ne2$, cioè $x\ne0$: $[-4,0[\,\cup\,]0,+\infty[$; \ con $t=2^x>0$: $t^2-3t+2\ge0$, cioè $t\le1$ o $t\ge2$: $x\le0$ oppure $x\ge1$.}

\begin{esercizio}[\liv{2}]
Determinare il dominio di: \ $\sqrt{\sin x}$; \ $\log(\cos x)$; \ $\dfrac1{2\sin x-1}$; \ $\sqrt{1-2\cos x}$ (in $[0,2\pi]$); \ $\arcsin(\log x)$.
\end{esercizio}
\risp{$\sin x\ge0$: $[2k\pi,\pi+2k\pi]$; \ $\cos x>0$: $]-\frac\pi2+2k\pi,\frac\pi2+2k\pi[$; \ $\sin x\ne\frac12$: $x\ne\frac\pi6+2k\pi$ e $x\ne\frac{5\pi}6+2k\pi$; \ $\cos x\le\frac12$: $[\frac\pi3,\frac{5\pi}3]$; \ $x>0$ e $-1\le\log x\le1$: $[\frac1e,e]$ ($k\in\Z$).}

\begin{esercizio}[\liv{3}]
Determinare il dominio di: \ $\sqrt{(2-x)\log x}$; \ $\sqrt{\sqrt{x+3}-x-1}$; \ $\log\dfrac{x^2-4}{x-1}$; \ $\sqrt{\arctan x-\frac\pi4}$.
\end{esercizio}
\risp{Prima: $x>0$ e $(2-x)\log x\ge0$; i due fattori hanno lo stesso segno (o uno è nullo) solo per $1\le x\le2$: $[1,2]$. \ Seconda: $x\ge-3$ e $\sqrt{x+3}\ge x+1$; se $x<-1$ è vera; se $x\ge-1$ si eleva al quadrato: $x^2+x-2\le0$, $-2\le x\le1$. Soluzione $[-3,1]$. \ Terza: $\frac{(x-2)(x+2)}{x-1}>0$: $]-2,1[\,\cup\,]2,+\infty[$. \ Quarta: $\arctan x\ge\frac\pi4$: $x\ge1$.}

\begin{esercizio}[\liv{3}]
Determinare il dominio di $f(x)=\sqrt{\log_{1/2}(x^2-3x+3)}$.
\end{esercizio}
\risp{$x^2-3x+3>0$ per ogni $x$ (discriminante negativo). Con base minore di $1$, $\log_{1/2}(x^2-3x+3)\ge0$ equivale a $x^2-3x+3\le1$, cioè $x^2-3x+2\le0$: $[1,2]$.}

%-------------------------------------------------------------------
\section{Successioni}

\begin{esercizio}[\liv{1}]
Scrivere i primi cinque termini e disegnare il grafico di: \ $a_n=\dfrac{n-1}n$ ($n\ge1$); \ $b_n=(-1)^n\,n$ ($n\ge1$); \ $c_n=1+\dfrac{(-1)^n}n$ ($n\ge1$); \ $d_0=1$, $d_{n+1}=3d_n-1$.
\end{esercizio}
\risp{$0,\frac12,\frac23,\frac34,\frac45$; \ $-1,2,-3,4,-5$; \ $0,\frac32,\frac23,\frac54,\frac45$; \ $1,2,5,14,41$. I grafici sono insiemi di punti isolati, da non unire.}

\begin{esercizio}[\liv{1}]
Trovare una formula esplicita per una successione che inizia con: \ $5,8,11,14,\dots$; \ $3,-6,12,-24,\dots$; \ $\frac12,\frac23,\frac34,\frac45,\dots$; \ $1,8,27,64,\dots$
\end{esercizio}
\risp{$5+3n$ ($n\ge0$); \ $3(-2)^n$ ($n\ge0$); \ $\frac n{n+1}$ ($n\ge1$); \ $n^3$ ($n\ge1$). Non sono le uniche possibili: un numero finito di termini non determina la successione.}

\begin{esercizio}[\liv{2}]
Scrivere i numeri di Fibonacci fino a $F_{15}$ e calcolare i rapporti $F_{n+1}/F_n$ per $n=10,\dots,14$, confrontandoli con $\varphi=\frac{1+\sqrt5}2$. Verificare che $F_{n+1}F_{n-1}-F_n^2=(-1)^n$ per $n=2,\dots,6$.
\end{esercizio}
\risp{$1,1,2,3,5,8,13,21,34,55,89,144,233,377,610$. Rapporti: $1{,}618182$; $1{,}617978$; $1{,}618056$; $1{,}618026$; $1{,}618037$, alternativamente sopra e sotto $\varphi\approx1{,}618034$. Per esempio $F_3F_1-F_2^2=2-1=1$, $F_4F_2-F_3^2=3-4=-1$, $F_5F_3-F_4^2=10-9=1$.}

\begin{esercizio}[\liv{2}]
Un segmento di $10$ cm è diviso in sezione aurea: quanto sono lunghe le due parti? Un rettangolo aureo ha il lato minore di $5$ cm: quanto è lungo il lato maggiore? Verificare con la calcolatrice che $\varphi^2=\varphi+1$ e $\frac1\varphi=\varphi-1$.
\end{esercizio}
\risp{La parte maggiore è $\frac{10}\varphi\approx6{,}18$ cm, la minore $\frac{10}{\varphi^2}\approx3{,}82$ cm. Il lato maggiore è $5\varphi\approx8{,}09$ cm. $\varphi^2\approx2{,}618=\varphi+1$; $\frac1\varphi\approx0{,}618=\varphi-1$.}

\begin{esercizio}[\liv{2}]
Ordinare le successioni $n^3$, $2^n$, $\log n$, $n!$, $\sqrt n$, $3^n$, $n^n$ calcolandone i valori per $n=10$. L'ordine trovato è quello ``per $n$ grande''? Confrontare $n^3$ e $2^n$ per $n=9$ e $n=10$.
\end{esercizio}
\risp{Per $n=10$: $\log n\approx2{,}3<\sqrt n\approx3{,}2<n^3=1000<2^n=1024<3^n=59\,049<n!\approx3{,}6\cdot10^6<n^n=10^{10}$. È l'ordine per $n$ grande, ma per $n=9$ si ha ancora $n^3=729>2^n=512$: il sorpasso di $2^n$ su $n^3$ avviene proprio a $n=10$.}

\begin{esercizio}[\liv{3}]
Sia $a_1=a_2=1$ e $a_{n+1}=a_n+2a_{n-1}$. (a) Scrivere i primi dieci termini e i rapporti $r_n=a_{n+1}/a_n$. (b) Mostrare che $r_n=1+\dfrac2{r_{n-1}}$, e trovare il candidato per il numero a cui si avvicinano i rapporti, come per Fibonacci.
\end{esercizio}
\risp{(a) $1,1,3,5,11,21,43,85,171,341$; rapporti $1$; $3$; $1{,}667$; $2{,}2$; $1{,}909$; $2{,}048$; $1{,}977$; $2{,}012$; $1{,}994$. (b) Dividendo per $a_n$: $r_n=1+2\frac{a_{n-1}}{a_n}=1+\frac2{r_{n-1}}$. Il candidato soddisfa $x=1+\frac2x$, cioè $x^2-x-2=0$, e la radice positiva è $x=2$.}

\begin{esercizio}[\liv{3}]
Il perimetro del poligono regolare di $n$ lati inscritto in una circonferenza di diametro $1$ è $p_n=n\sin\frac\pi n$. Calcolare esattamente $p_4$, $p_6$ e, con la formula di sottrazione $\sin(\alpha-\beta)=\sin\alpha\cos\beta-\cos\alpha\sin\beta$ applicata a $\frac\pi{12}=\frac\pi3-\frac\pi4$, anche $p_{12}$. Confrontare con $\pi$.
\end{esercizio}
\risp{$p_4=4\cdot\frac{\sqrt2}2=2\sqrt2\approx2{,}83$; $p_6=6\cdot\frac12=3$; $\sin\frac\pi{12}=\frac{\sqrt3}2\cdot\frac{\sqrt2}2-\frac12\cdot\frac{\sqrt2}2=\frac{\sqrt6-\sqrt2}4$, quindi $p_{12}=3(\sqrt6-\sqrt2)\approx3{,}106$. Tutti minori di $\pi\approx3{,}1416$, e sempre più vicini.}

%-------------------------------------------------------------------
\stampasoluzioni
   % Lezioni 10--12
% !TEX root =esercizi.tex
% ---------------------------------------------------------------------
% Scheda 5: Lezioni 13--15 (limiti di successioni). Per scelta del docente
% molti esercizi sui campi di esistenza (riepilogo) e sui limiti, con tutti
% i casi svolti a lezione. I limiti notevoli sono in forma generale:
% sin(a_n)/a_n con a_n -> 0, (1+1/a_n)^{a_n} con a_n -> +-infinito.
% ---------------------------------------------------------------------
\chapter{Lezioni 13--15 (21--26 ottobre 2026)}

\begin{vuoto}[Argomenti svolti]{}
\begin{itemize}
\item[\bul] \textbf{Lezione 13 --- Limiti di successioni}. Definizione di limite finito e infinito; successioni convergenti, divergenti, irregolari. Unicità del limite, limitatezza delle successioni convergenti, permanenza del segno. Teorema del confronto. Il limite $\frac{\sin a_n}{a_n}\to1$ per $a_n\to0$ e le sue conseguenze. Algebra dei limiti finiti.
\item[\bul] \textbf{Lezione 14 --- Il calcolo dei limiti}. Aritmetica dell'infinito. Le forme indeterminate $+\infty-\infty$, $0\cdot\infty$, $\frac\infty\infty$, $\frac00$; quozienti di polinomi, razionalizzazione. Gerarchia degli infiniti.
\item[\bul] \textbf{Lezione 15 --- Successioni monotone e il numero $e$}. Teorema sulle successioni monotone. Il numero $e$; $\bigl(1+\frac1{a_n}\bigr)^{a_n}\to e$ per $a_n\to\pm\infty$. Esponenziale e logaritmo lungo le successioni; $a_n^{b_n}=e^{b_n\ln a_n}$; forme determinate e forme indeterminate $1^\infty$, $0^0$, $\infty^0$.
\end{itemize}
\end{vuoto}

Il livello di difficoltà è indicato con \liv{1} (esercizio di base), \liv{2} (esercizio standard), \liv{3} (esercizio più impegnativo).

\begin{vuoto}[Promemoria: i limiti notevoli in forma generale]{}
Se $a_n\to0$ con $a_n\ne0$:
\[
\frac{\sin a_n}{a_n}\to1,\qquad \frac{\tan a_n}{a_n}\to1,\qquad \frac{1-\cos a_n}{a_n^2}\to\frac12 .
\]
Se $a_n\to+\infty$ oppure $a_n\to-\infty$:
\[
\Bigl(1+\frac1{a_n}\Bigr)^{a_n}\to e .
\]
Per usarli si cerca, dentro l'espressione, la successione $a_n$ giusta, e la si fa comparire anche dove manca, moltiplicando e dividendo. Per esempio $n\sin\frac3n=3\cdot\dfrac{\sin(3/n)}{3/n}\to3$, con $a_n=\frac3n$.
\end{vuoto}

%-------------------------------------------------------------------
\section{Campi di esistenza}

\begin{esercizio}[\liv{1}]
Determinare il dominio di: \ $\sqrt{x^2-4x}$; \ $\ln(x^2-x-6)$; \ $\dfrac1{2^x-4}$; \ $\sqrt{3-|x|}$; \ $\log_3(5-x)+\sqrt x$.
\end{esercizio}
\risp{$x\le0$ oppure $x\ge4$; \ $x<-2$ oppure $x>3$; \ $x\ne2$; \ $[-3,3]$; \ $x\ge0$ e $x<5$: $[0,5[$.}

\begin{esercizio}[\liv{2}]
Determinare il dominio di:
\[
\sqrt{\frac{x+2}{1-x}},\qquad \ln\frac{x-1}{x^2-4},\qquad \sqrt{\ln(2-x)},\qquad \frac1{\ln(x+3)},\qquad \sqrt{e^{2x}-3e^x+2}.
\]
\end{esercizio}
\risp{$[-2,1[$; \ $\frac{x-1}{(x-2)(x+2)}>0$: $]-2,1[\,\cup\,]2,+\infty[$; \ $\ln(2-x)\ge0$, cioè $2-x\ge1$: $x\le1$; \ $x>-3$ e $\ln(x+3)\ne0$, cioè $x\ne-2$; \ con $t=e^x>0$: $t^2-3t+2\ge0$, cioè $t\le1$ o $t\ge2$: $x\le0$ oppure $x\ge\ln2$.}

\begin{esercizio}[\liv{2}]
Determinare il dominio di: \ $\sqrt{2\cos x-1}$ (in $[0,2\pi]$); \ $\ln(\sin x)$; \ $\arcsin\dfrac{x-1}2$; \ $\arccos(\ln x)$; \ $\dfrac1{\tan x-1}$; \ $\sqrt{\arcsin x}$.
\end{esercizio}
\risp{$\cos x\ge\frac12$: $[0,\frac\pi3]\cup[\frac{5\pi}3,2\pi]$; \ $\sin x>0$: $]2k\pi,\pi+2k\pi[$; \ $-1\le\frac{x-1}2\le1$: $[-1,3]$; \ $-1\le\ln x\le1$: $[\frac1e,e]$; \ $x\ne\frac\pi2+k\pi$ (dove la tangente non esiste) e $x\ne\frac\pi4+k\pi$; \ $\arcsin x\ge0$: $[0,1]$ ($k\in\Z$).}

\begin{esercizio}[\liv{3}]
Determinare il dominio di:
\[
\sqrt{\log_{1/2}(x^2-2x)},\qquad \ln\bigl(\sqrt{x+1}-x+1\bigr),\qquad \sqrt{\frac{\ln x-1}{x-3}},\qquad \arcsin(2^x-1),\qquad \frac{\sqrt{|x-1|-2}}{\ln x}.
\]
\end{esercizio}
\risp{Primo: $x^2-2x>0$ e, con base minore di $1$, $x^2-2x\le1$; cioè ($x<0$ o $x>2$) e $1-\sqrt2\le x\le1+\sqrt2$: $[1-\sqrt2,0[\,\cup\,]2,1+\sqrt2]$. \ Secondo: $x\ge-1$ e $\sqrt{x+1}>x-1$; vera se $x<1$; per $x\ge1$ si eleva: $x^2-3x<0$, $0<x<3$. Dominio $[-1,3[$. \ Terzo: $x>0$, $x\ne3$, quoziente $\ge0$: il numeratore è $\le0$ per $x\le e$ e il denominatore è negativo per $x<3$; dominio $]0,e]\cup\,]3,+\infty[$. \ Quarto: $0\le2^x\le2$, cioè $x\le1$. \ Quinto: $|x-1|\ge2$ ($x\le-1$ o $x\ge3$), $x>0$, $x\ne1$: $[3,+\infty[$.}

\begin{esercizio}[\liv{3}]
Determinare il dominio di $f(x)=\ln\bigl(1-\sqrt{1-x^2}\bigr)$.
\end{esercizio}
\risp{Serve $1-x^2\ge0$, cioè $|x|\le1$, e $1-\sqrt{1-x^2}>0$, cioè $\sqrt{1-x^2}<1$, cioè $x\ne0$: $[-1,0[\,\cup\,]0,1]$.}

\begin{esercizio}[\liv{2}]
Determinare dominio e segno di $f(x)=\dfrac{e^x-2}{\ln(x+1)}$.
\end{esercizio}
\risp{Dominio: $x>-1$ e $x\ne0$. Il numeratore è positivo per $x>\ln2$, il denominatore per $x>0$. Quindi $f>0$ in $]-1,0[$ e in $]\ln2,+\infty[$, $f<0$ in $]0,\ln2[$, $f=0$ in $x=\ln2$.}

%-------------------------------------------------------------------
\section{La definizione di limite}

\begin{esercizio}[\liv{1}]
Verificare con la definizione, trovando l'indice $N$: (a) $\dfrac{3n+2}{n+1}\to3$, con $\varepsilon=10^{-2}$; \ (b) $\dfrac2{\sqrt n}\to0$; \ (c) $n^2+1\to+\infty$.
\end{esercizio}
\risp{(a) $\bigl|\frac{3n+2}{n+1}-3\bigr|=\frac1{n+1}<\varepsilon$ per $n>\frac1\varepsilon-1$: con $\varepsilon=10^{-2}$, $N=100$. (b) $\frac2{\sqrt n}<\varepsilon$ per $n>\frac4{\varepsilon^2}$. (c) $n^2+1>M$ per $n>\sqrt M$ (se $M>1$).}

\begin{esercizio}[\liv{2}]
Dire se sono convergenti, divergenti o irregolari: \ $(-1)^n\dfrac n{n+1}$; \ $\cos\dfrac{n\pi}2$; \ $\dfrac{(-1)^n}{n^2}$; \ $n+(-1)^n$; \ $(-1)^nn^2$.
\end{esercizio}
\risp{Irregolare (i termini pari si avvicinano a $1$, i dispari a $-1$); \ irregolare (vale $0,-1,0,1,\dots$); \ converge a $0$ (confronto con $\pm\frac1{n^2}$); \ diverge a $+\infty$ ($n+(-1)^n\ge n-1$); \ irregolare (non limitata, di segno alterno).}

%-------------------------------------------------------------------
\section{Confronto e limiti trigonometrici}

\begin{esercizio}[\liv{1}]
Calcolare usando il teorema del confronto: \ $\dfrac{\sin n}{n}$; \ $\dfrac{n+\cos n}{n}$; \ $\dfrac{(-1)^n+3}{\sqrt n}$; \ $\dfrac{\lfloor\sqrt n\rfloor}{\sqrt n}$.
\end{esercizio}
\risp{$0$ ($|\sin n|\le1$); \ $1$ ($\frac{n-1}n\le\cdot\le\frac{n+1}n$); \ $0$ (fra $\frac2{\sqrt n}$ e $\frac4{\sqrt n}$); \ $1$ ($\sqrt n-1<\lfloor\sqrt n\rfloor\le\sqrt n$, quindi $1-\frac1{\sqrt n}<\cdot\le1$).}

\begin{esercizio}[\liv{1}]
Calcolare, riconoscendo la successione $a_n\to0$: \ $n\sin\dfrac2n$; \ $n^2\sin\dfrac1{n^2}$; \ $\dfrac{\sin(1/n)}{\sin(2/n)}$; \ $n\tan\dfrac3n$; \ $\dfrac{\sin(3/n)}{\tan(5/n)}$.
\end{esercizio}
\risp{$2$ ($a_n=\frac2n$); \ $1$ ($a_n=\frac1{n^2}$); \ $\frac{\sin(1/n)}{1/n}\cdot\frac{2/n}{\sin(2/n)}\cdot\frac12\to\frac12$; \ $3$; \ $\frac{\sin(3/n)}{3/n}\cdot\frac{5/n}{\tan(5/n)}\cdot\frac35\to\frac35$.}

\begin{esercizio}[\liv{2}]\label{es:s5cos}
(Il limite di $1-\cos a_n$, rivisto.) Sia $a_n\to0$ con $a_n\ne0$. (a) Moltiplicando numeratore e denominatore per $1+\cos a_n$, mostrare che
\[
\frac{1-\cos a_n}{a_n^2}=\Bigl(\frac{\sin a_n}{a_n}\Bigr)^2\cdot\frac1{1+\cos a_n},
\]
e dedurre che $\dfrac{1-\cos a_n}{a_n^2}\to\dfrac12$. (b) Calcolare $n^2\Bigl(1-\cos\dfrac1n\Bigr)$, \ $n\Bigl(1-\cos\dfrac1n\Bigr)$, \ $n^2\Bigl(1-\cos\dfrac3n\Bigr)$.
\end{esercizio}
\risp{(a) $(1-\cos a_n)(1+\cos a_n)=1-\cos^2a_n=\sin^2a_n$; poiché $\cos a_n\to1$, il secondo fattore tende a $\frac12$ e il primo a $1$. \ (b) $\frac12$ ($a_n=\frac1n$); \ $\frac{1-\cos(1/n)}{(1/n)^2}\cdot\frac1n\to\frac12\cdot0=0$; \ $9\cdot\frac{1-\cos(3/n)}{(3/n)^2}\to\frac92$.}

\begin{esercizio}[\liv{2}]
Calcolare: \ $n\sin\dfrac1{n^2}$; \ $\dfrac{1-\cos(2/n)}{\sin^2(1/n)}$; \ $\sqrt n\,\sin\dfrac1n$; \ $\dfrac{\tan(1/n)-\sin(1/n)}{\sin(1/n)}$.
\end{esercizio}
\risp{$\frac{\sin(1/n^2)}{1/n^2}\cdot\frac1n\to0$; \ $\frac{1-\cos(2/n)}{(2/n)^2}\cdot4\cdot\bigl(\frac{1/n}{\sin(1/n)}\bigr)^2\to\frac12\cdot4=2$; \ $\frac{\sin(1/n)}{1/n}\cdot\frac1{\sqrt n}\to0$; \ $\frac1{\cos(1/n)}-1\to0$.}

\begin{esercizio}[\liv{2}]
(Il perimetro dei poligoni, L12.) Il perimetro del poligono regolare di $n$ lati inscritto in una circonferenza di diametro $1$ è $p_n=n\sin\frac\pi n$. Dimostrare che $p_n\to\pi$.
\end{esercizio}
\risp{$p_n=\pi\cdot\dfrac{\sin(\pi/n)}{\pi/n}$, con $a_n=\frac\pi n\to0$: $p_n\to\pi\cdot1=\pi$. La lunghezza della circonferenza di diametro $1$ è davvero il limite dei perimetri dei poligoni inscritti.}

\begin{esercizio}[\liv{3}]
Calcolare $\lim n^3\Bigl(\sin\dfrac1n-\tan\dfrac1n\Bigr)$. (Suggerimento: $\sin a-\tan a=\sin a\cdot\dfrac{\cos a-1}{\cos a}$.)
\end{esercizio}
\risp{Con $a_n=\frac1n$: $n^3(\sin a_n-\tan a_n)=\frac{\sin a_n}{a_n}\cdot\frac{\cos a_n-1}{a_n^2}\cdot\frac1{\cos a_n}\to1\cdot\bigl(-\frac12\bigr)\cdot1=-\frac12$.}

%-------------------------------------------------------------------
\section{Algebra dei limiti e forme indeterminate}

\begin{esercizio}[\liv{1}]
Calcolare (non ci sono forme indeterminate): \ $\Bigl(2+\dfrac1n\Bigr)\Bigl(3-\dfrac2{n^2}\Bigr)$; \ $n+\dfrac1n$; \ $\dfrac{1-n}{2+\frac1n}$; \ $\dfrac1{n^2+\sin n}$; \ $\dfrac{3+\frac{(-1)^n}n}{1+\frac1{\sqrt n}}$; \ $n^2\Bigl(2+\dfrac1n\Bigr)$.
\end{esercizio}
\risp{$6$; \ $+\infty$; \ $-\infty$ ($\frac{-\infty}2$); \ $0$ ($\frac1{+\infty}$); \ $3$; \ $+\infty$.}

\begin{esercizio}[\liv{1}]
Calcolare: \ $\dfrac{2n^3-n+4}{5n^3+n^2}$; \ $\dfrac{n^2+1}{n^3-2}$; \ $\dfrac{1-4n^4}{n^3+2n}$; \ $\dfrac{\sqrt n+2n}{3n-\sqrt n}$; \ $\dfrac{(n+1)^3-n^3}{n^2}$.
\end{esercizio}
\risp{$\frac25$; \ $0$; \ $-\infty$; \ $\frac23$ (si raccoglie $n$); \ $\frac{3n^2+3n+1}{n^2}\to3$.}

\begin{esercizio}[\liv{2}]
Calcolare, razionalizzando: \ $\sqrt{n^2+4n}-n$; \ $\sqrt{n+1}-\sqrt{n-1}$; \ $n\bigl(\sqrt{n^2+1}-n\bigr)$; \ $\sqrt{n^2+n}-\sqrt{n^2-n}$.
\end{esercizio}
\risp{$\frac{4n}{\sqrt{n^2+4n}+n}\to2$; \ $\frac2{\sqrt{n+1}+\sqrt{n-1}}\to0$; \ $\frac{n}{\sqrt{n^2+1}+n}\to\frac12$; \ $\frac{2n}{\sqrt{n^2+n}+\sqrt{n^2-n}}\to1$.}

\begin{esercizio}[\liv{3}]
Calcolare $\lim\bigl(\sqrt[3]{n^3+n^2}-n\bigr)$. (Suggerimento: $a-b=\dfrac{a^3-b^3}{a^2+ab+b^2}$.)
\end{esercizio}
\risp{Con $a=\sqrt[3]{n^3+n^2}$ e $b=n$: $a^3-b^3=n^2$ e $a^2+ab+b^2=n^2\bigl(\sqrt[3]{(1+\frac1n)^2}+\sqrt[3]{1+\frac1n}+1\bigr)$, quindi il limite è $\frac13$.}

\begin{esercizio}[\liv{2}]
Calcolare: \ $n^2\Bigl(\dfrac1n-\dfrac1{n+1}\Bigr)$; \ $n-\ln n$; \ $2^n-n^{10}$; \ $\ln(n+1)-\ln n$; \ $\ln(2n^2+1)-2\ln n$.
\end{esercizio}
\risp{$\frac{n^2}{n(n+1)}\to1$; \ $n\bigl(1-\frac{\ln n}n\bigr)\to+\infty$; \ $2^n\bigl(1-\frac{n^{10}}{2^n}\bigr)\to+\infty$; \ $\ln\bigl(1+\frac1n\bigr)\to\ln1=0$; \ $\ln\frac{2n^2+1}{n^2}\to\ln2$.}

%-------------------------------------------------------------------
\section{Gerarchia degli infiniti}

\begin{esercizio}[\liv{1}]
Calcolare: \ $\dfrac{n^5+3^n}{2^n+4^n}$; \ $\dfrac{n!+2^n}{n!-n^3}$; \ $\dfrac{(\ln n)^2}{\sqrt n}$; \ $\dfrac{n^n}{n!}$; \ $\dfrac{5^n+n!}{n^n}$; \ $\dfrac{n^2\,2^n}{3^n}$.
\end{esercizio}
\risp{Si divide per $4^n$: $\frac{n^5/4^n+(3/4)^n}{(1/2)^n+1}\to0$; \ $1$; \ $\bigl(\frac{\ln n}{n^{1/4}}\bigr)^2\to0$; \ $+\infty$ (reciproco di $\frac{n!}{n^n}\to0^+$); \ $\frac{5^n}{n!}\cdot\frac{n!}{n^n}+\frac{n!}{n^n}\to0$; \ $\frac{n^2}{(3/2)^n}\to0$.}

\begin{esercizio}[\liv{2}]
Ordinare dalla più lenta alla più veloce, per $n$ grande: \ $n\ln n$, \ $1{,}1^n$, \ $\sqrt n$, \ $n!$, \ $3\ln n$, \ $n^n$, \ $n^2$. Giustificare ogni confronto con un limite.
\end{esercizio}
\risp{$3\ln n\ll\sqrt n\ll n\ln n\ll n^2\ll1{,}1^n\ll n!\ll n^n$. Per esempio $\frac{n\ln n}{n^2}=\frac{\ln n}n\to0$ e $\frac{\sqrt n}{n\ln n}=\frac1{\sqrt n\ln n}\to0$; gli altri sono casi della gerarchia.}

%-------------------------------------------------------------------
\section{Successioni monotone}

\begin{esercizio}[\liv{2}]
Sia $a_1=1$ e $a_{n+1}=\sqrt{6+a_n}$. (a) Dimostrare per induzione che $1\le a_n<3$ e che $(a_n)$ è strettamente crescente. (b) Dedurre che converge. (c) Mostrare che $3-a_{n+1}=\dfrac{3-a_n}{3+a_{n+1}}\le\dfrac{3-a_n}3$ e calcolare il limite.
\end{esercizio}
\risp{(a) Se $1\le a_n<3$, allora $\sqrt7\le a_{n+1}<\sqrt9=3$. Inoltre $a_{n+1}>a_n$ equivale a $6+a_n>a_n^2$, cioè $(a_n-3)(a_n+2)<0$, vero per $0\le a_n<3$. (b) Crescente e limitata: converge. (c) $3-\sqrt{6+a_n}=\frac{9-(6+a_n)}{3+a_{n+1}}$; quindi $0<3-a_n\le\frac2{3^{n-1}}\to0$ e il limite è $3$ (i termini: $2{,}646$; $2{,}940$; $2{,}990$; $2{,}998$; \dots).}

\begin{esercizio}[\liv{2}]
Sia $b_n=\dfrac{2^n}{n!}$. (a) Mostrare che $b_{n+1}=\dfrac2{n+1}\,b_n$ e dedurre che $(b_n)$ è decrescente. (b) Dedurre che converge. (c) Passando al limite nella relazione del punto (a), trovare il limite.
\end{esercizio}
\risp{(a) $\frac2{n+1}\le1$ per $n\ge1$: decrescente ($b_1=b_2=2$, poi strettamente). (b) Decrescente e limitata inferiormente da $0$: converge a un $\ell\ge0$. (c) $b_{n+1}\to\ell$ e $\frac2{n+1}b_n\to0\cdot\ell=0$: per l'unicità del limite $\ell=0$.}

\begin{esercizio}[\liv{3}]
Sia $c_n=\bigl(1+\frac1n\bigr)^{n+1}$. (a) Mostrare che $\dfrac{c_{n-1}}{c_n}=\Bigl(1+\dfrac1{n^2-1}\Bigr)^n\cdot\dfrac n{n+1}$ per $n\ge2$. (b) Usando la disuguaglianza di Bernoulli, dedurre che $c_{n-1}>c_n$: la successione è decrescente. (c) Dedurre che $c_n\to e$ e che $\bigl(1+\frac1n\bigr)^n<e<\bigl(1+\frac1n\bigr)^{n+1}$ per ogni $n$. Che cosa si ottiene per $n=10$?
\end{esercizio}
\risp{(a) $c_{n-1}=\bigl(\frac n{n-1}\bigr)^n$, $c_n=\bigl(\frac{n+1}n\bigr)^{n+1}$, e $\frac n{n-1}\cdot\frac n{n+1}=\frac{n^2}{n^2-1}$. (b) Per Bernoulli $\bigl(1+\frac1{n^2-1}\bigr)^n\ge1+\frac n{n^2-1}>1+\frac1n$, quindi $\frac{c_{n-1}}{c_n}>\frac{n+1}n\cdot\frac n{n+1}=1$. (c) $c_n=e_n\bigl(1+\frac1n\bigr)\to e\cdot1=e$; $e_n$ cresce verso $e$ e $c_n$ decresce verso $e$, quindi $e_n<e<c_n$. Per $n=10$: $2{,}5937<e<2{,}8531$.}

%-------------------------------------------------------------------
\section{Il numero $e$ e le forme esponenziali}

\begin{esercizio}[\liv{1}]
Calcolare, riconducendosi a $\bigl(1+\frac1{a_n}\bigr)^{a_n}\to e$: \ $\Bigl(1+\dfrac2n\Bigr)^n$; \ $\Bigl(1-\dfrac1n\Bigr)^n$; \ $\Bigl(1+\dfrac1{n^2}\Bigr)^{n^2}$; \ $\Bigl(1+\dfrac1{3n}\Bigr)^n$; \ $\Bigl(\dfrac{n+3}n\Bigr)^{2n}$; \ $\Bigl(1+\dfrac1n\Bigr)^{n/2}$.
\end{esercizio}
\risp{$\bigl[(1+\frac1{n/2})^{n/2}\bigr]^2\to e^2$; \ $a_n=-n$: $\bigl[(1+\frac1{-n})^{-n}\bigr]^{-1}\to e^{-1}$; \ $e$ ($a_n=n^2$); \ $e^{1/3}$; \ $\bigl[(1+\frac1{n/3})^{n/3}\bigr]^6\to e^6$; \ $\sqrt e$.}

\begin{esercizio}[\liv{2}]
Calcolare: \ $\Bigl(\dfrac{n+1}{n-1}\Bigr)^n$; \ $\Bigl(\dfrac{n^2+1}{n^2}\Bigr)^n$; \ $\Bigl(1+\sin\dfrac1n\Bigr)^n$; \ $\Bigl(1+\dfrac1n\Bigr)^{\sqrt n}$; \ $\Bigl(1+\dfrac1{\sqrt n}\Bigr)^n$.
\end{esercizio}
\risp{$\frac{n+1}{n-1}=1+\frac1{a_n}$ con $a_n=\frac{n-1}2$, ed $n=a_n\cdot\frac{2n}{n-1}$: il limite è $e^2$; \ $\bigl[(1+\frac1{n^2})^{n^2}\bigr]^{1/n}\to e^0=1$; \ con $a_n=\frac1{\sin(1/n)}\to+\infty$: $(1+\frac1{a_n})^{a_n\cdot n\sin(1/n)}\to e^1=e$; \ $\bigl[(1+\frac1n)^n\bigr]^{1/\sqrt n}\to e^0=1$; \ $\bigl[(1+\frac1{\sqrt n})^{\sqrt n}\bigr]^{\sqrt n}$: la base tende a $e>1$ e l'esponente a $+\infty$, quindi $+\infty$.}

\begin{esercizio}[\liv{3}]
Calcolare $\lim\Bigl(\cos\dfrac1n\Bigr)^{n^2}$. (Suggerimento: $\cos\frac1n=1+\frac1{a_n}$ con $a_n=\frac1{\cos(1/n)-1}\to-\infty$.)
\end{esercizio}
\risp{$\bigl(\cos\frac1n\bigr)^{n^2}=\bigl(1+\frac1{a_n}\bigr)^{a_n\cdot b_n}$ con $b_n=n^2\bigl(\cos\frac1n-1\bigr)\to-\frac12$: il limite è $e^{-1/2}\approx0{,}607$.}

\begin{esercizio}[\liv{2}]
Calcolare (forme $0^0$ e $\infty^0$), scrivendo $a_n^{b_n}=e^{b_n\ln a_n}$: \ $\sqrt[n]n$; \ $\sqrt[n]{n^2+3^n}$; \ $\Bigl(\dfrac1n\Bigr)^{1/\ln n}$; \ $(\ln n)^{1/n}$; \ $\sqrt[n]{2^n+5^n}$; \ $\bigl(n^n\bigr)^{1/n^2}$.
\end{esercizio}
\risp{$e^{\frac{\ln n}n}\to e^0=1$; \ $3\sqrt[n]{1+\frac{n^2}{3^n}}\to3$ (o con il confronto: $3^n\le n^2+3^n\le2\cdot3^n$); \ $e^{-\frac{\ln n}{\ln n}}=e^{-1}$: una forma $0^0$ che vale $\frac1e$; \ $e^{\frac{\ln\ln n}n}\to1$; \ $5$; \ $n^{1/n}\to1$.}

\begin{esercizio}[\liv{1}]
Calcolare (forme determinate): \ $\Bigl(\dfrac12+\dfrac1n\Bigr)^n$; \ $\Bigl(2-\dfrac1n\Bigr)^n$; \ $\Bigl(\dfrac1n\Bigr)^n$; \ $\Bigl(3+\dfrac1n\Bigr)^{-n}$; \ $\Bigl(\dfrac{n+1}n\Bigr)^{-n}$. Quale delle cinque è invece una forma indeterminata?
\end{esercizio}
\risp{$0$ (base verso $\frac12<1$, esponente verso $+\infty$); \ $+\infty$; \ $0$ ($0^{+\infty}$); \ $0$; \ l'ultima è una forma $1^\infty$: $\bigl[(1+\frac1n)^n\bigr]^{-1}\to\frac1e$.}

%-------------------------------------------------------------------
\section{Riepilogo}

\begin{esercizio}[\liv{2}]
Calcolare:
\[
\frac{n^2\sin(1/n)+\ln n}{n+1};\qquad \frac{3^{n+1}+2^n}{3^n+n^3};\qquad n\bigl(\ln(n+1)-\ln n\bigr);\qquad \sqrt n\bigl(\sqrt{n+1}-\sqrt n\bigr).
\]
\end{esercizio}
\risp{$n^2\sin\frac1n=n\cdot\frac{\sin(1/n)}{1/n}$, quindi il quoziente è $\frac{n\cdot\frac{\sin(1/n)}{1/n}+\ln n}{n+1}\to1$; \ $\frac{3+(2/3)^n}{1+n^3/3^n}\to3$; \ $\ln\bigl(1+\frac1n\bigr)^n\to\ln e=1$; \ $\frac{\sqrt n}{\sqrt{n+1}+\sqrt n}\to\frac12$.}

\begin{esercizio}[\liv{3}]
Al variare del parametro $a>0$, calcolare $\lim\dfrac{a^n+n^2}{2^n+n}$.
\end{esercizio}
\risp{Se $a>2$: $+\infty$ (si divide per $2^n$: $(\frac a2)^n\to+\infty$). Se $a=2$: $1$. Se $0<a<2$: $0$ (anche per $a\le1$, dove $a^n$ resta limitato e $\frac{n^2}{2^n}\to0$).}

\begin{esercizio}[\liv{3}]
Vero o falso? Motivare. (a) Se $a_n\to1$ e $b_n\to+\infty$, allora $a_n^{b_n}\to1$. (b) Se $a_n\to\frac12$ e $b_n\to+\infty$, allora $a_n^{b_n}\to0$. (c) Se $\frac{a_n}{b_n}\to1$ e $b_n\to+\infty$, allora $a_n-b_n\to0$. (d) Se $a_n\to0$, allora $\frac{\sin a_n}{a_n}\to1$.
\end{esercizio}
\risp{(a) Falso: $(1+\frac1n)^n\to e$. (b) Vero: è una forma determinata ($b_n\ln a_n\to-\infty$). (c) Falso: $a_n=n+1$, $b_n=n$, la differenza è $1$; con $a_n=n+\sqrt n$ diverge. (d) Falso così come è scritto: serve $a_n\ne0$ (altrimenti il quoziente non è definito), ma con $a_n\ne0$ è vero.}

%-------------------------------------------------------------------
\stampasoluzioni
   % Lezioni 13--15
% !TEX root =esercizi.tex
% ---------------------------------------------------------------------
% Scheda 6: Lezioni 16--18 (limiti di funzioni, continuita' e zeri,
% derivata). La regola della funzione composta non e' ancora stata fatta
% (L19): negli esercizi sulle derivate compaiono solo somme, prodotti e
% quozienti di funzioni elementari. Discontinuita' secondo la convenzione
% del docente: anche nei punti in cui f non e' definita ma lo e' vicino.
% ---------------------------------------------------------------------
\chapter{Lezioni 16--18 (28 ottobre -- 2 novembre 2026)}

\begin{vuoto}[Argomenti svolti]{}
\begin{itemize}
\item[\bul] \textbf{Lezione 16 --- Limiti di funzioni}. Definizione di limite con le successioni; legame con la definizione con $\varepsilon$ e $\delta$. La catena dei limiti notevoli: da $\frac{\sin a_n}{a_n}$ e dal numero $e$ fino a $\log(1+a_n)$, $e^{a_n}$, $(1+a_n)^\alpha$. Regole di calcolo, forme indeterminate, gerarchia e limiti notevoli per le funzioni; il cambio di variabile.
\item[\bul] \textbf{Lezione 17 --- Continuità e teorema degli zeri}. Limiti destro e sinistro; non esistenza del limite. Asintoti verticali, orizzontali, obliqui. Funzioni continue; punti di discontinuità e loro classificazione. Teorema degli zeri con il metodo di bisezione; teorema dei valori intermedi; teorema di Weierstrass.
\item[\bul] \textbf{Lezione 18 --- La derivata}. Rapporto incrementale e derivata; retta tangente; significato fisico. Derivata destra e sinistra, punto angoloso; derivabilità e continuità. Derivate delle funzioni elementari con i limiti notevoli. Regole della somma, del prodotto e del quoziente; tabella delle derivate.
\end{itemize}
\end{vuoto}

Il livello di difficoltà è indicato con \liv{1} (esercizio di base), \liv{2} (esercizio standard), \liv{3} (esercizio più impegnativo).

\begin{vuoto}[Promemoria: limiti notevoli e derivate]{}
Per $x\to0$:
\begin{gather*}
\frac{\sin x}x\to1,\qquad \frac{\tan x}x\to1,\qquad \frac{1-\cos x}{x^2}\to\frac12,\\
\frac{\log(1+x)}x\to1,\qquad \frac{e^x-1}x\to1,\qquad \frac{(1+x)^\alpha-1}x\to\alpha,\qquad (1+x)^{1/x}\to e;
\end{gather*}
per $x\to\pm\infty$: $\bigl(1+\frac1x\bigr)^x\to e$. \emph{Cambio di variabile}: si pone $t=g(x)$, si dice dove va $t$, e si riscrive tutto con $t$.

Derivate: $(x^\alpha)'=\alpha x^{\alpha-1}$, $(e^x)'=e^x$, $(a^x)'=a^x\log a$, $(\log x)'=\frac1x$, $(\log_ax)'=\frac1{x\log a}$, $(\sin x)'=\cos x$, $(\cos x)'=-\sin x$, $(\tan x)'=\frac1{\cos^2x}$. Regole: $(f+g)'=f'+g'$, $(fg)'=f'g+fg'$, $\bigl(\frac fg\bigr)'=\frac{f'g-fg'}{g^2}$.
\end{vuoto}

%-------------------------------------------------------------------
\section{La definizione di limite}

\begin{esercizio}[\liv{1}]
Verificare con la definizione (con le successioni): \ (a) $\lim_{x\to2}(x^2+1)=5$; \ (b) $\lim_{x\to+\infty}\dfrac{3x-1}{x+2}=3$; \ (c) $\lim_{x\to1}\dfrac1{(x-1)^2}=+\infty$.
\end{esercizio}
\risp{(a) Se $x_n\to2$, allora $x_n^2+1\to5$ (algebra dei limiti). (b) Se $x_n\to+\infty$, si raccoglie: $\frac{3-1/x_n}{1+2/x_n}\to3$. (c) Se $x_n\to1$ con $x_n\ne1$, allora $(x_n-1)^2\to0$ ed è positivo, quindi il reciproco tende a $+\infty$.}

\begin{esercizio}[\liv{1}]
Per $\lim_{x\to1}(4x-3)=1$, trovare un $\delta$ che vada bene per $\varepsilon=0{,}01$ nella definizione con $\varepsilon$ e $\delta$.
\end{esercizio}
\risp{$|(4x-3)-1|=4|x-1|<0{,}01$ se $|x-1|<0{,}0025$: va bene $\delta=0{,}0025$.}

\begin{esercizio}[\liv{2}]
Dimostrare, trovando due successioni, che non esistono: \ (a) $\lim_{x\to+\infty}\sin x$; \ (b) $\lim_{x\to0}\cos\dfrac\pi x$.
\end{esercizio}
\risp{(a) Con $x_n=n\pi$ si ha $\sin x_n=0$; con $y_n=\frac\pi2+2n\pi$ si ha $\sin y_n=1$. (b) Con $x_n=\frac1{2n}$ si ha $\cos(2n\pi)=1$; con $y_n=\frac1{2n+1}$ si ha $\cos\bigl((2n+1)\pi\bigr)=-1$.}

%-------------------------------------------------------------------
\section{Forme indeterminate e gerarchia}

\begin{esercizio}[\liv{1}]
Calcolare:
\begin{gather*}
\text{(a) }\lim_{x\to+\infty}\frac{2x^3-x+1}{x^3+4x^2};\qquad \text{(b) }\lim_{x\to-\infty}\frac{x^2-3}{2x+1};\qquad \text{(c) }\lim_{x\to3}\frac{x^2-9}{x^2-5x+6};\\
\text{(d) }\lim_{x\to-1}\frac{x^3+1}{x+1};\qquad \text{(e) }\lim_{x\to0}\frac{\sqrt{4+x}-2}x;\qquad \text{(f) }\lim_{x\to+\infty}\bigl(\sqrt{x^2+3x}-x\bigr).
\end{gather*}
\end{esercizio}
\risp{(a) $2$; (b) $-\infty$ (domina $\frac{x^2}{2x}=\frac x2$); (c) $\frac{(x-3)(x+3)}{(x-3)(x-2)}=\frac{x+3}{x-2}\to6$; (d) $x^3+1=(x+1)(x^2-x+1)$: il limite è $3$; (e) razionalizzando, $\frac1{\sqrt{4+x}+2}\to\frac14$; (f) $\frac{3x}{\sqrt{x^2+3x}+x}=\frac3{\sqrt{1+3/x}+1}\to\frac32$.}

\begin{esercizio}[\liv{2}]
Calcolare con la gerarchia degli infiniti:
\begin{gather*}
\text{(a) }\lim_{x\to+\infty}\frac{e^x+x^{100}}{e^x-\log x};\qquad \text{(b) }\lim_{x\to+\infty}\frac{\log(x^3+1)}{\log x};\qquad \text{(c) }\lim_{x\to0^+}x^2\log x;\\
\text{(d) }\lim_{x\to+\infty}\frac{(\log x)^{10}}{x^{0{,}1}};\qquad \text{(e) }\lim_{x\to+\infty}x^{1/\log x}.
\end{gather*}
\end{esercizio}
\risp{(a) Si divide per $e^x$: $\frac{1+x^{100}/e^x}{1-\log x/e^x}\to1$; (b) $\log(x^3+1)=3\log x+\log\bigl(1+\frac1{x^3}\bigr)$, quindi il quoziente tende a $3$; (c) con $t=\frac1x\to+\infty$: $-\frac{\log t}{t^2}\to0$; (d) $0$: ogni potenza di $x$, anche piccola, batte ogni potenza del logaritmo; (e) $x^{1/\log x}=e^{\log x/\log x}=e$ per ogni $x>1$.}

%-------------------------------------------------------------------
\section{Limiti notevoli e cambio di variabile}

\begin{esercizio}[\liv{1}]
Calcolare con i limiti notevoli:
\begin{gather*}
\text{(a) }\lim_{x\to0}\frac{\sin4x}{3x};\qquad \text{(b) }\lim_{x\to0}\frac{1-\cos2x}{x^2};\qquad \text{(c) }\lim_{x\to0}\frac{\tan5x}x;\qquad \text{(d) }\lim_{x\to0}\frac{\log(1+4x)}x;\\
\text{(e) }\lim_{x\to0}\frac{e^{2x}-1}{3x};\qquad \text{(f) }\lim_{x\to0}\frac{\sqrt{1+x}-1}x;\\ \text{(g) }\lim_{x\to0}(1-3x)^{1/x};\qquad \text{(h) }\lim_{x\to+\infty}\Bigl(1+\frac1{2x}\Bigr)^x.
\end{gather*}
\end{esercizio}
\risp{(a) $\frac43\cdot\frac{\sin4x}{4x}\to\frac43$; (b) $4\cdot\frac{1-\cos2x}{(2x)^2}\to2$; (c) $5$; (d) $4$; (e) $\frac23$; (f) $\frac12$ ($\alpha=\frac12$; oppure razionalizzando); (g) $\bigl[(1-3x)^{\frac1{-3x}}\bigr]^{-3}\to e^{-3}$; (h) $\Bigl[\bigl(1+\frac1{2x}\bigr)^{2x}\Bigr]^{1/2}\to\sqrt e$.}

\begin{esercizio}[\liv{2}]
Calcolare con un cambio di variabile, dicendo ogni volta dove va la nuova variabile:
\begin{gather*}
\text{(a) }\lim_{x\to+\infty}x\tan\frac1x;\qquad \text{(b) }\lim_{x\to\pi}\frac{1+\cos x}{(x-\pi)^2};\qquad \text{(c) }\lim_{x\to1}\frac{x^3-1}{\log x};\\
\text{(d) }\lim_{x\to0^+}x\log^2x;\qquad \text{(e) }\lim_{x\to2}\frac{2^x-4}{x-2};\qquad \text{(f) }\lim_{x\to+\infty}x\bigl(2^{1/x}-1\bigr).
\end{gather*}
\end{esercizio}
\risp{(a) $t=\frac1x\to0^+$: $\frac{\tan t}t\to1$; (b) $t=x-\pi\to0$, $\cos x=-\cos t$: $\frac{1-\cos t}{t^2}\to\frac12$; (c) $t=x-1\to0$: $\frac{(1+t)^3-1}t\cdot\frac t{\log(1+t)}\to3\cdot1=3$; (d) $t=\frac1x\to+\infty$: $\frac{(\log t)^2}t\to0$; (e) $t=x-2\to0$: $\frac{4(2^t-1)}t=4\log2\cdot\frac{e^{t\log2}-1}{t\log2}\to4\log2$; (f) $t=\frac1x\to0^+$: $\frac{2^t-1}t\to\log2$.}

\begin{esercizio}[\liv{3}]
Calcolare:
\[
\text{(a) }\lim_{x\to0}\frac{e^{\sin3x}-1}{\log(1+x)};\qquad \text{(b) }\lim_{x\to0^+}(\sin x)^x;\qquad \text{(c) }\lim_{x\to0}(\cos x)^{1/x^2};\qquad \text{(d) }\lim_{x\to+\infty}\Bigl(\cos\frac1x\Bigr)^{x^2}.
\]
\end{esercizio}
\risp{(a) $\frac{e^{\sin3x}-1}{\sin3x}\cdot\frac{\sin3x}{3x}\cdot3\cdot\frac x{\log(1+x)}\to1\cdot1\cdot3\cdot1=3$ (con $t=\sin3x\to0$ nel primo fattore). (b) $(\sin x)^x=e^{x\log\sin x}$, e $x\log\sin x=x\log x+x\log\frac{\sin x}x\to0+0$: il limite è $1$. (c) Con $a=\cos x-1\to0$: $(\cos x)^{1/x^2}=\bigl[(1+a)^{1/a}\bigr]^{a/x^2}$, e $\frac a{x^2}\to-\frac12$: il limite è $e^{-1/2}$. (d) Con $t=\frac1x\to0^+$ diventa (c): $e^{-1/2}$.}

%-------------------------------------------------------------------
\section{Limiti destro e sinistro, asintoti}

\begin{esercizio}[\liv{1}]
Calcolare i limiti destro e sinistro e dire se il limite esiste: \ (a) $\dfrac1{x-2}$ per $x\to2$; \ (b) $e^{-1/x}$ per $x\to0$; \ (c) $\dfrac{|x|}{x^2+x}$ per $x\to0$; \ (d) $\arctan\dfrac1x$ per $x\to0$.
\end{esercizio}
\risp{(a) $+\infty$ e $-\infty$: no; (b) da destra $-\frac1x\to-\infty$, quindi $0$; da sinistra $-\frac1x\to+\infty$, quindi $+\infty$: no; (c) $\frac{|x|}{x(x+1)}$: da destra $\frac1{x+1}\to1$, da sinistra $-\frac1{x+1}\to-1$: no; (d) $\frac\pi2$ e $-\frac\pi2$: no.}

\begin{esercizio}[\liv{2}]
Determinare tutti gli asintoti di:
\[
\text{(a) }\frac{x^2+x}{x-1};\qquad \text{(b) }\frac{2x^2-1}{x^2+1};\qquad \text{(c) }\sqrt{x^2+1};
\]
\[
\text{(d) }x\,e^{1/x};\qquad \text{(e) }\log(x^2-4);\qquad \text{(f) }\frac{x^3}{x^2-4}.
\]
\end{esercizio}
\risp{(a) Verticale $x=1$; obliquo $y=x+2$ da entrambe le parti ($m=1$, $q=\lim\frac{2x}{x-1}=2$). (b) Orizzontale $y=2$; nessun verticale. (c) Obliqui: $y=x$ per $x\to+\infty$ e $y=-x$ per $x\to-\infty$ (per esempio $\sqrt{x^2+1}-x=\frac1{\sqrt{x^2+1}+x}\to0$). (d) Verticale $x=0$ da destra ($t=\frac1x\to+\infty$: $\frac{e^t}t\to+\infty$; da sinistra il limite è $0$); obliquo $y=x+1$ da entrambe le parti: $m=\lim e^{1/x}=1$, $q=\lim x\bigl(e^{1/x}-1\bigr)=1$ con $t=\frac1x\to0$. (e) Dominio $|x|>2$; verticali $x=\pm2$ (il limite è $-\infty$); nessun orizzontale né obliquo ($\frac{\log(x^2-4)}x\to0$ ma $\log(x^2-4)\to+\infty$). (f) Verticali $x=\pm2$; obliquo $y=x$ ($\frac{x^3}{x^2-4}-x=\frac{4x}{x^2-4}\to0$).}

%-------------------------------------------------------------------
\section{Continuità e punti di discontinuità}

\begin{esercizio}[\liv{1}]
Determinare il parametro in modo che la funzione sia continua nel suo insieme di definizione:
\begin{gather*}
\text{(a) }f(x)=\begin{cases}ax^2+1 & x\le1\\ 3x-a & x>1\end{cases};\qquad \text{(b) }f(x)=\begin{cases}\dfrac{\sin2x}x & x\ne0\\ b & x=0\end{cases};\\ \text{(c) }f(x)=\begin{cases}\dfrac{e^x-1}x & x<0\\ x+c & x\ge0\end{cases};\qquad \text{(d) }f(x)=\begin{cases}(1+x)^{1/x} & x>0\\ k & x=0\end{cases}.
\end{gather*}
\end{esercizio}
\risp{(a) In $1$: $a+1=3-a$, $a=1$. (b) $b=\lim_{x\to0}\frac{\sin2x}x=2$. (c) Il limite sinistro in $0$ è $1$, il valore è $c$: $c=1$. (d) $k=e$.}

\begin{esercizio}[\liv{2}]
Trovare e classificare i punti di discontinuità (eliminabile, di salto, di seconda specie):
\begin{gather*}
\text{(a) }\frac{x^2-4}{x-2};\qquad \text{(b) }\frac{\sin x}{|x|};\qquad \text{(c) }e^{1/x};\qquad \text{(d) }\frac{1-\cos x}{x^2};\\
\text{(e) }\frac x{x^2-1};\qquad \text{(f) }x\sin\frac1x;\qquad \text{(g) }\log|x|;\qquad \text{(h) }\lfloor x\rfloor+\lfloor-x\rfloor.
\end{gather*}
\end{esercizio}
\risp{(a) $2$, eliminabile (limite $4$). (b) $0$, di salto (limiti destro e sinistro $1$ e $-1$). (c) $0$, di seconda specie (limite destro $+\infty$). (d) $0$, eliminabile (limite $\frac12$). (e) $\pm1$, di seconda specie. (f) $0$, eliminabile: $|x\sin\frac1x|\le|x|\to0$. (g) $0$, di seconda specie (limite $-\infty$). (h) La funzione vale $0$ negli interi e $-1$ altrove: ogni intero è una discontinuità eliminabile (il limite è $-1$, il valore è $0$).}

%-------------------------------------------------------------------
\section{Zeri, valori intermedi, Weierstrass}

\begin{esercizio}[\liv{1}]
Dimostrare che ogni equazione ha almeno una soluzione nell'intervallo indicato: \ (a) $x^3+2x-5=0$ in $[1,2]$; \ (b) $\log x=2-x$ in $[1,2]$; \ (c) $2^x=3x$ in $[0,1]$.
\end{esercizio}
\risp{Si applica il teorema degli zeri a una funzione continua: (a) $f(1)=-2<0$, $f(2)=7>0$; (b) $g(x)=\log x+x-2$: $g(1)=-1$, $g(2)=\log2>0$; (c) $h(x)=2^x-3x$: $h(0)=1$, $h(1)=-1$.}

\begin{esercizio}[\liv{2}]
Per l'equazione (a) dell'esercizio precedente, eseguire quattro passi del metodo di bisezione. Quanto dista al più dallo zero il punto medio dell'ultimo intervallo? Quanti passi servono per un errore minore di $10^{-3}$?
\end{esercizio}
\risp{Punti medi e segni di $f$: $1{,}5$ ($+$), $1{,}25$ ($-$), $1{,}375$ ($+$), $1{,}3125$ ($-$): lo zero sta in $[1{,}3125;\,1{,}375]$, di lunghezza $\frac1{16}$; il punto medio $1{,}34375$ dista al più $\frac1{32}\approx0{,}031$ dallo zero (che è $1{,}3283\dots$). Dopo $n$ passi il punto medio dista al più $\frac1{2^{n+1}}$: serve $2^{n+1}>1000$, cioè $n\ge9$.}

\begin{esercizio}[\liv{2}]
(a) Dimostrare che $f(x)=x^5-3x+1$ assume il valore $10$ in qualche punto di $[1,2]$. (b) Quali delle seguenti funzioni hanno massimo e minimo? Il teorema di Weierstrass si applica? \ $x^2-2x$ su $[0,3]$; \ $\frac1x$ su $[1,2]$; \ $\frac1x$ su $]0,1]$; \ $e^{-x^2}$ su $\R$; \ $\arctan x$ su $\R$.
\end{esercizio}
\risp{(a) $f(1)=-1$, $f(2)=27$, e $f$ è continua: per i valori intermedi assume ogni valore fra $-1$ e $27$, quindi anche $10$. (b) $x^2-2x$ su $[0,3]$: minimo $-1$ in $1$, massimo $3$ in $3$ (Weierstrass si applica). $\frac1x$ su $[1,2]$: massimo $1$, minimo $\frac12$ (si applica). $\frac1x$ su $]0,1]$: minimo $1$, nessun massimo (intervallo non chiuso). $e^{-x^2}$ su $\R$: massimo $1$ in $0$, nessun minimo ($\inf=0$; $\R$ non è limitato). $\arctan x$ su $\R$: né massimo né minimo ($\sup=\frac\pi2$, $\inf=-\frac\pi2$).}

%-------------------------------------------------------------------
\section{La derivata}

\begin{esercizio}[\liv{1}]
Calcolare con la definizione: \ (a) $f'(1)$ per $f(x)=2x^2-3x$; \ (b) $f'(0)$ per $f(x)=\dfrac1{x+2}$; \ (c) $f'(9)$ per $f(x)=\sqrt x$.
\end{esercizio}
\risp{(a) $\frac{2(1+h)^2-3(1+h)+1}h=\frac{h+2h^2}h\to1$; (b) $\frac{\frac1{2+h}-\frac12}h=-\frac1{2(2+h)}\to-\frac14$; (c) razionalizzando, $\frac{\sqrt{9+h}-3}h=\frac1{\sqrt{9+h}+3}\to\frac16$.}

\begin{esercizio}[\liv{2}]
Dire se la funzione è derivabile nel punto indicato, calcolando derivata destra e sinistra:
\begin{gather*}
\text{(a) }|x^2-1|\text{ in }\pm1;\qquad \text{(b) }\begin{cases}x^2 & x\le0\\ x^3 & x>0\end{cases}\text{ in }0;\\ \text{(c) }\begin{cases}x+1 & x<0\\ e^x & x\ge0\end{cases}\text{ in }0;\qquad \text{(d) }\begin{cases}x^2 & x\le1\\ 2x & x>1\end{cases}\text{ in }1.
\end{gather*}
\end{esercizio}
\risp{(a) In $1$: derivata destra $2$ (lì $f=x^2-1$), sinistra $-2$ (lì $f=1-x^2$): punto angoloso; in $-1$ allo stesso modo (destra $2$, sinistra $-2$). (b) Continua; derivate destra e sinistra $0$: derivabile, $f'(0)=0$. (c) Continua ($f(0)=1$); derivate destra e sinistra $1$: derivabile, $f'(0)=1$. (d) $f(1)=1$, ma il limite destro è $2$: non è continua, quindi non è derivabile.}

\begin{esercizio}[\liv{1}]
Derivare con le regole:
\begin{gather*}
\text{(a) }5x^4-3x^2+2x;\qquad \text{(b) }x^3e^x;\qquad \text{(c) }x^2\log x;\qquad \text{(d) }\frac{x+1}{x^2+1};\\
\text{(e) }\sin x\cos x;\qquad \text{(f) }\frac{e^x}{1+e^x};\qquad \text{(g) }\sqrt x\,\log x;\qquad \text{(h) }x^{2/3};\\
\text{(i) }x\cdot3^x;\qquad \text{(j) }\tan x-x;\qquad \text{(k) }\frac{\log x}{x^2};\qquad \text{(l) }\frac x{\sin x}.
\end{gather*}
\end{esercizio}
\risp{(a) $20x^3-6x+2$; (b) $(x^3+3x^2)\,e^x$; (c) $2x\log x+x$; (d) $\frac{1-2x-x^2}{(x^2+1)^2}$; (e) $\cos^2x-\sin^2x$; (f) $\frac{e^x}{(1+e^x)^2}$; (g) $\frac{\log x}{2\sqrt x}+\frac{\sqrt x}x=\frac{\log x+2}{2\sqrt x}$; (h) $\frac23x^{-1/3}$; (i) $3^x(1+x\log3)$; (j) $\frac1{\cos^2x}-1=\tan^2x$; (k) $\frac{1-2\log x}{x^3}$; (l) $\frac{\sin x-x\cos x}{\sin^2x}$.}

\begin{esercizio}[\liv{2}]
(a) Scrivere l'equazione della retta tangente: a $y=x^3-2x$ in $x_0=1$; a $y=e^x$ in $x_0=1$; a $y=\log x$ in $x_0=e$; a $y=\sin x$ in $x_0=\pi$. (b) Trovare i punti in cui la tangente è orizzontale: $x^3-3x$; $x\,e^x$; $\dfrac{\log x}x$; $x^2e^x$.
\end{esercizio}
\risp{(a) $y=x-2$; \ $y=ex$; \ $y=\frac xe$; \ $y=-(x-\pi)$. (b) $3x^2-3=0$: $x=\pm1$; \ $(1+x)e^x=0$: $x=-1$; \ $\frac{1-\log x}{x^2}=0$: $x=e$; \ $(x^2+2x)e^x=0$: $x=0$ e $x=-2$.}

\begin{esercizio}[\liv{2}]
Un corpo si muove lungo una retta con legge oraria $s(t)=t^3-6t^2+9t$ (metri, secondi). Calcolare la velocità $v(t)=s'(t)$. In quali istanti il corpo è fermo? In quali intervalli di tempo torna indietro?
\end{esercizio}
\risp{$v(t)=3t^2-12t+9=3(t-1)(t-3)$ m/s: è fermo in $t=1$ e $t=3$; torna indietro ($v<0$) per $1<t<3$.}

\begin{esercizio}[\liv{3}]
Mostrare che la retta tangente alla parabola $y=x^2$ nel punto $(a,a^2)$ incontra l'asse $y$ nel punto $(0,-a^2)$. Usare questo fatto per disegnare, con riga e squadra, la tangente alla parabola in un punto dato.
\end{esercizio}
\risp{La tangente è $y=a^2+2a(x-a)=2ax-a^2$: per $x=0$ vale $-a^2$. Per disegnarla basta congiungere $(a,a^2)$ con il punto dell'asse $y$ simmetrico di $(0,a^2)$ rispetto all'origine.}

%-------------------------------------------------------------------
\section{Riepilogo}

\begin{esercizio}[\liv{2}]
Vero o falso? Motivare con una dimostrazione o con un controesempio. \ (a) Se $f$ è continua in $x_0$, allora è derivabile in $x_0$. \ (b) Se $f$ è derivabile in $x_0$, allora è continua in $x_0$. \ (c) Se $f(a)f(b)<0$, allora $f$ ha uno zero in $]a,b[$. \ (d) $\frac{\sin x}x\to1$ per $x\to+\infty$. \ (e) Una funzione continua in $]0,1[$ ha massimo. \ (f) La derivata di un prodotto è il prodotto delle derivate.
\end{esercizio}
\risp{(a) Falso: $|x|$ in $0$. (b) Vero (teorema di L18). (c) Falso senza la continuità: $\frac1x$ su $[-1,1]$. (d) Falso: per il confronto il limite è $0$; $\frac{\sin x}x\to1$ vale per $x\to0$. (e) Falso: $f(x)=x$ su $]0,1[$. (f) Falso: $(x\cdot x)'=2x$, mentre $1\cdot1=1$.}

%-------------------------------------------------------------------
\stampasoluzioni
   % Lezioni 16--18
% !TEX root =esercizi.tex
% ---------------------------------------------------------------------
% Scheda 7: Lezioni 19--21 (regole di derivazione; estremi, Fermat,
% ottimizzazione; L21 = esercitazione del 9/11, con gli esercizi del
% fascicolo esercitazione-9-novembre e le loro soluzioni).
% Le derivate delle inverse trigonometriche sono state svolte in L19 ma
% non sono materia della prima prova.
% ---------------------------------------------------------------------
\chapter{Lezioni 19--21 (4 -- 9 novembre 2026)}

\begin{vuoto}[Argomenti svolti]{}
\begin{itemize}
\item[\bul] \textbf{Lezione 19 --- Regole di derivazione}. Dimostrazione delle regole della somma, del prodotto e del quoziente; regola della funzione composta; derivata della funzione inversa (argomento geometrico) e delle inverse trigonometriche; approssimazione lineare (formula di Taylor al primo ordine), il simbolo $o$ piccolo e gli sviluppi al primo ordine delle funzioni elementari.
\item[\bul] \textbf{Lezione 20 --- Massimi, minimi e il teorema di Fermat}. Massimi e minimi assoluti e locali; teorema di Fermat; metodo dei candidati in un intervallo chiuso; problemi di ottimizzazione; il problema di Fermat dei tre punti; il rettangolo di perimetro assegnato. Lettura: il problema di Didone e la disuguaglianza isoperimetrica.
\item[\bul] \textbf{Lezione 21 --- Esercitazione in vista della prima prova}. Correzione del fascicolo distribuito prima della lezione: calcolo combinatorio, percentuali e progressioni, disequazioni, campi di esistenza, limiti di successioni e di funzioni, calcolo di derivate. Gli esercizi sono riportati qui sotto, con le soluzioni.
\end{itemize}
\end{vuoto}

Il livello di difficoltà è indicato con \liv{1} (esercizio di base), \liv{2} (esercizio standard), \liv{3} (esercizio più impegnativo).

\begin{vuoto}[Promemoria: regole di derivazione]{}
$(f+g)'=f'+g'$, \quad $(fg)'=f'g+fg'$, \quad $\Bigl(\dfrac fg\Bigr)'=\dfrac{f'g-fg'}{g^2}$, \quad $\bigl(g(f(x))\bigr)'=g'\bigl(f(x)\bigr)\,f'(x)$.\\[2pt]
Inversa: $(f^{-1})'(y_0)=\dfrac1{f'(x_0)}$ con $y_0=f(x_0)$. \quad $(\arctan x)'=\dfrac1{1+x^2}$, \ $(\arcsin x)'=\dfrac1{\sqrt{1-x^2}}$, \ $(\arccos x)'=-\dfrac1{\sqrt{1-x^2}}$.\\[2pt]
Approssimazione lineare: $f(x_0+h)\approx f(x_0)+f'(x_0)\,h$; \ per $x\to0$: $e^x=1+x+o(x)$, $\sin x=x+o(x)$, $\log(1+x)=x+o(x)$, $(1+x)^\alpha=1+\alpha x+o(x)$.
\end{vuoto}

%-------------------------------------------------------------------
\section{Regole di derivazione e funzione composta (L19)}

\begin{esercizio}[\liv{1}]
Derivare: \ (a) $x^3\sin x$; \ (b) $\dfrac{x^2+1}{x-1}$; \ (c) $e^{-x^2}$; \ (d) $\log(\cos x)$; \ (e) $(3x-2)^7$; \ (f) $\sqrt{x^2+9}$.
\end{esercizio}
\risp{(a) $3x^2\sin x+x^3\cos x$. (b) $\frac{2x(x-1)-(x^2+1)}{(x-1)^2}=\frac{x^2-2x-1}{(x-1)^2}$. (c) $-2x\,e^{-x^2}$. (d) $\frac{-\sin x}{\cos x}=-\tan x$, dove $\cos x>0$. (e) $21(3x-2)^6$. (f) $\frac x{\sqrt{x^2+9}}$.}

\begin{esercizio}[\liv{2}]
Derivare: \ (a) $e^x\sin2x$; \ (b) $\log\dfrac{x+1}{x^2+1}$; \ (c) $x^x$ per $x>0$ (si scriva $x^x=e^{x\log x}$); \ (d) $\sin^3(x^2)$.
\end{esercizio}
\risp{(a) $e^x(\sin2x+2\cos2x)$. (b) Con le proprietà del logaritmo, $\log(x+1)-\log(x^2+1)$: derivata $\frac1{x+1}-\frac{2x}{x^2+1}$, per $x>-1$. (c) $e^{x\log x}(\log x+1)=x^x(\log x+1)$. (d) Tre funzioni annidate: $3\sin^2(x^2)\cdot\cos(x^2)\cdot2x=6x\sin^2(x^2)\cos(x^2)$.}

\begin{esercizio}[\liv{2}]
(a) Con la regola della funzione inversa, ricavare $(\log y)'=\frac1y$ da $(e^x)'=e^x$, e $(\sqrt y)'=\frac1{2\sqrt y}$ da $(x^2)'=2x$ (per $x>0$). \ (b) Derivare $\arctan(e^x)$ e $\arcsin\sqrt x$ (per $0<x<1$).
\end{esercizio}
\risp{(a) Se $y=e^x$, allora $(\log y)'=\frac1{e^x}=\frac1y$; se $y=x^2$ con $x>0$, allora $(\sqrt y)'=\frac1{2x}=\frac1{2\sqrt y}$. (b) $\frac{e^x}{1+e^{2x}}$; \ $\frac1{\sqrt{1-x}}\cdot\frac1{2\sqrt x}=\frac1{2\sqrt{x(1-x)}}$.}

\begin{esercizio}[\liv{1}]
Scrivere l'equazione della retta tangente: \ (a) a $y=\log x$ in $x_0=e$; \ (b) a $y=x^3-2x$ in $x_0=-1$; \ (c) a $y=\tan x$ in $x_0=\frac\pi4$.
\end{esercizio}
\risp{(a) $f(e)=1$, $f'(e)=\frac1e$: $y=1+\frac{x-e}e=\frac xe$ (passa per l'origine). (b) $f(-1)=1$, $f'(-1)=3-2=1$: $y=x+2$. (c) $f\bigl(\frac\pi4\bigr)=1$, $f'\bigl(\frac\pi4\bigr)=\frac1{\cos^2(\pi/4)}=2$: $y=1+2\bigl(x-\frac\pi4\bigr)$.}

\begin{esercizio}[\liv{2}]
Con l'approssimazione lineare stimare $\sqrt{4{,}1}$, $\log(1{,}02)$ ed $e^{0{,}03}$, e confrontare con la calcolatrice.
\end{esercizio}
\risp{$\sqrt{4+h}\approx2+\frac h4$: $\sqrt{4{,}1}\approx2{,}025$ (calcolatrice $2{,}02485$). $\log(1+h)\approx h$: $\log1{,}02\approx0{,}02$ (calcolatrice $0{,}01980$). $e^h\approx1+h$: $e^{0{,}03}\approx1{,}03$ (calcolatrice $1{,}03045$). L'errore è dell'ordine di $h^2$.}

\begin{esercizio}[\liv{2}]
Il raggio di una sfera è misurato in $r=5$ cm, con un errore di $0{,}1$ cm. Stimare con l'approssimazione lineare l'errore sul volume $V=\frac43\pi r^3$, assoluto e relativo.
\end{esercizio}
\risp{$\Delta V\approx V'(r)\,\Delta r=4\pi r^2\cdot0{,}1=10\pi\approx31{,}4$ cm$^3$, su un volume di circa $523{,}6$ cm$^3$: errore relativo $\approx6\%$, cioè tre volte l'errore relativo sul raggio ($2\%$), perché $V$ è proporzionale a $r^3$.}

\begin{esercizio}[\liv{3}]
Con gli sviluppi al primo ordine calcolare: \ (a) $\displaystyle\lim_{x\to0}\frac{e^{2x}-1-\sin x}x$; \ (b) $\displaystyle\lim_{x\to0}\frac{\sqrt{1+2x}-\cos x}x$; \ (c) $\displaystyle\lim_{x\to0}\frac{(1+x)^5-1}{\sin3x}$.
\end{esercizio}
\risp{(a) $e^{2x}=1+2x+o(x)$, $\sin x=x+o(x)$: il numeratore è $x+o(x)$, limite $1$. (b) $\sqrt{1+2x}=1+x+o(x)$, $\cos x=1+o(x)$: limite $1$. (c) $\frac{5x+o(x)}{3x+o(x)}\to\frac53$.}

%-------------------------------------------------------------------
\section{Massimi, minimi e ottimizzazione (L20)}

\begin{esercizio}[\liv{1}]
Trovare massimo e minimo assoluti con il metodo dei candidati: \ (a) $x^3-3x^2+1$ su $[-1,3]$; \ (b) $x+\dfrac4x$ su $[1,4]$; \ (c) $|x^2-1|$ su $[-2,2]$.
\end{esercizio}
\risp{(a) $f'=3x^2-6x$, critici $0$ e $2$; $f(-1)=-3$, $f(0)=1$, $f(2)=-3$, $f(3)=1$: massimo $1$ in $0$ e in $3$, minimo $-3$ in $-1$ e in $2$. (b) $f'=1-\frac4{x^2}$, critico $2$; $f(1)=5$, $f(2)=4$, $f(4)=5$: massimo $5$ in $1$ e $4$, minimo $4$ in $2$. (c) Candidati: estremi $\pm2$ (valore $3$), punti di non derivabilità $\pm1$ (valore $0$), punto critico $0$ (valore $1$): massimo $3$, minimo $0$.}

\begin{esercizio}[\liv{2}]
Una funzione derivabile su $[0,4]$ ha il massimo in $x=4$, e la sua derivata sinistra in $4$ vale $3$. Questo contraddice il teorema di Fermat? E se il massimo fosse in $x=2$, con $f'(2)=3$?
\end{esercizio}
\risp{Nel primo caso no: Fermat vale per i punti \emph{interni}, e $4$ è un estremo dell'intervallo. Nel secondo caso sì: in un punto interno di massimo in cui $f$ è derivabile deve essere $f'=0$; quindi una tale funzione non esiste.}

\begin{esercizio}[\liv{2}]
Fra i rettangoli inscritti in un semicerchio di raggio $R$, con un lato sul diametro, trovare quello di area massima.
\end{esercizio}
\risp{Con il lato sul diametro lungo $2x$ e l'altezza $\sqrt{R^2-x^2}$: $A=2x\sqrt{R^2-x^2}$ per $0\le x\le R$. Conviene il quadrato $A^2=4x^2(R^2-x^2)$, con derivata $8xR^2-16x^3$, nulla in $x=\frac R{\sqrt2}$; agli estremi l'area è $0$. Il rettangolo ha lati $\sqrt2\,R$ e $\frac R{\sqrt2}$, e area $R^2$.}

\begin{esercizio}[\liv{2}]
Un recinto rettangolare di $200$ m$^2$ ha un lato lungo un muro, dove la rete non serve. Quali dimensioni rendono minima la rete? Quanta rete servirebbe per un recinto a semicerchio della stessa area, con il diametro sul muro?
\end{esercizio}
\risp{Con $x$ i lati perpendicolari al muro: rete $L(x)=2x+\frac{200}x$ per $x>0$, che tende a $+\infty$ agli estremi; $L'=2-\frac{200}{x^2}=0$ per $x=10$: recinto $10\times20$ m, $40$ m di rete. Semicerchio: $\frac12\pi r^2=200$, $r\approx11{,}28$ m, rete $\pi r\approx35{,}4$ m. È il problema di Didone (\S20.6).}

\begin{esercizio}[\liv{3}]
Tre villaggi sono ai vertici di un triangolo equilatero di lato $6$ km. Dove conviene mettere il punto di raccordo della rete di strade più corta? Quanto è lunga la rete? Confrontarla con la rete fatta da due lati del triangolo.
\end{esercizio}
\risp{Per la regola dei $120^\circ$ il raccordo è il centro del triangolo, che vede ogni lato sotto un angolo di $120^\circ$. Dista $\frac6{\sqrt3}=2\sqrt3$ km da ciascun vertice: la rete è lunga $6\sqrt3\approx10{,}39$ km, contro i $12$ km di due lati.}

\begin{esercizio}[\liv{1}]
Con una corda lunga $60$ m si recinta (a) un quadrato, (b) un esagono regolare, (c) un cerchio. Calcolare le tre aree.
\end{esercizio}
\risp{(a) $15^2=225$ m$^2$. (b) Sei triangoli equilateri di lato $10$: $6\cdot\frac{\sqrt3}4\cdot100\approx259{,}8$ m$^2$. (c) $\frac{60^2}{4\pi}\approx286{,}5$ m$^2$. Il cerchio vince, come dice la disuguaglianza isoperimetrica.}

%-------------------------------------------------------------------
\section{Lezione 21: esercitazione in vista della prima prova}

Sono gli esercizi del fascicolo dell'esercitazione del 9 novembre (capitolo 21 degli appunti). Ogni soluzione si chiude con un breve riepilogo della teoria che serve (\emph{In breve}).

\subsection*{Calcolo combinatorio}
\begin{esercizio}[\liv{1}]
(a) In quanti modi si possono disporre $5$ libri diversi su uno scaffale? (b) Quante ``parole'' di $4$ lettere (anche prive di senso) si possono scrivere con le $21$ lettere dell'alfabeto italiano? E quante senza lettere ripetute? (c) In un'aiuola ci sono $10$ piante, e se ne scelgono $3$ da campionare: in quanti modi?
\end{esercizio}
\risp{(a) Permutazioni: $5!=120$. (b) Disposizioni con ripetizione: $21^4=194\,481$; senza ripetizione: $21\cdot20\cdot19\cdot18=143\,640$. (c) L'ordine non conta: $\binom{10}3=\frac{10\cdot9\cdot8}{3!}=120$. \textit{In breve:} \emph{Permutazioni} di $n$ oggetti: $n!$. \emph{Disposizioni} di $k$ oggetti scelti fra $n$ (l'ordine conta): $n^k$ con ripetizione, $n(n-1)\cdots(n-k+1)$ senza. \emph{Combinazioni} (l'ordine non conta): $\binom nk=\frac{n!}{k!\,(n-k)!}$.}

\begin{esercizio}[\liv{2}]
(a) Un test ha $10$ domande vero/falso. Quanti sono i modi di rispondere a tutte? In quanti modi si possono dare esattamente $7$ risposte giuste? (b) Calcolare il coefficiente di $x^3$ nello sviluppo di $(2+x)^5$. (c) Da un gruppo di $6$ biologi e $4$ geologi si forma una commissione di $3$ persone. Quante commissioni contengono almeno un geologo?
\end{esercizio}
\risp{(a) $2^{10}=1024$; per scegliere quali $7$ domande sono giuste: $\binom{10}7=120$. (b) Binomio di Newton: il termine con $x^3$ è $\binom53\,2^2x^3$, e il coefficiente è $10\cdot4=40$. (c) Conviene contare il complementare: tutte le commissioni sono $\binom{10}3=120$, quelle di soli biologi $\binom63=20$; quindi $120-20=100$. \textit{In breve:} $n$ scelte fra due possibilità: $2^n$; scegliere \emph{quali} $k$ fra $n$: $\binom nk$. Binomio di Newton: $(a+b)^n=\sum_{k=0}^n\binom nk a^{n-k}b^k$. ``Almeno uno'': conviene contare il complementare (tutti meno ``nessuno'').}


\subsection*{Percentuali e progressioni}
\begin{esercizio}[\liv{1}]
Una popolazione di $2000$ individui aumenta del $10\%$ il primo anno, diminuisce del $10\%$ il secondo, aumenta del $20\%$ il terzo. (a) Quanti individui ci sono dopo tre anni? (b) Qual è la variazione percentuale complessiva? (c) Quale tasso annuo costante darebbe lo stesso risultato in tre anni?
\end{esercizio}
\risp{(a) $2000\cdot1{,}1\cdot0{,}9\cdot1{,}2=2376$. (b) $\frac{2376}{2000}=1{,}188$: $+18{,}8\%$ (non $+20\%$: le variazioni si moltiplicano, non si sommano). (c) $(1+p)^3=1{,}188$, quindi $1+p=\sqrt[3]{1{,}188}\approx1{,}059$: circa $5{,}9\%$ all'anno. \textit{In breve:} Aumentare del $p\%$ vuol dire moltiplicare per $1+\frac p{100}$, diminuire del $p\%$ per $1-\frac p{100}$: le variazioni successive si \emph{moltiplicano}, non si sommano. Un tasso costante $p$ per $n$ periodi dà il fattore $(1+p)^n$.}

\begin{esercizio}[\liv{1}]
Un prezzo viene scontato del $20\%$ e poi aumentato del $20\%$. (a) Il prezzo finale è maggiore, minore o uguale a quello iniziale, e di quanto per cento? (b) Di quanto per cento si sarebbe dovuto aumentare il prezzo scontato per tornare a quello iniziale?
\end{esercizio}
\risp{(a) $0{,}8\cdot1{,}2=0{,}96$: il prezzo è diminuito del $4\%$. (b) Serve $0{,}8\,(1+p)=1$, cioè $1+p=1{,}25$: un aumento del $25\%$. \textit{In breve:} $(1-p)(1+p)=1-p^2<1$: uno sconto e un aumento della stessa percentuale non si compensano. Per annullare il fattore $f$ serve il fattore $\frac1f$.}

\begin{esercizio}[\liv{2}]
(a) Calcolare $1+2+4+\dots+2^{10}$. (b) Un paziente assume ogni giorno, alla stessa ora, una dose di $100$ mg di un farmaco; in un giorno l'organismo ne elimina il $40\%$, cioè ne resta il $60\%$. Mostrare che subito dopo la $n$-esima dose la quantità nell'organismo è
\[
Q_n=100\,\bigl(1+0{,}6+0{,}6^2+\dots+0{,}6^{n-1}\bigr)\ \text{mg},
\]
e calcolarla con la formula della somma geometrica. Quanto vale dopo $5$ dosi? A che cosa tende per $n\to\infty$?
\end{esercizio}
\risp{(a) È una somma geometrica di ragione $2$ con $11$ termini: $\frac{2^{11}-1}{2-1}=2047$. (b) Dopo la prima dose: $100$; prima della seconda ne resta $60$, e con la seconda dose $100+60$; prima della terza ne resta $0{,}6(100+60)=60+36$, e con la terza $100+60+36=100(1+0{,}6+0{,}6^2)$; e così via. Per la formula, $Q_n=100\cdot\frac{1-0{,}6^n}{1-0{,}6}=250\,(1-0{,}6^n)$. Dopo $5$ dosi: $250\,(1-0{,}07776)\approx230{,}6$ mg. Poiché $0{,}6^n\to0$, $Q_n\to250$ mg: la quantità si stabilizza. \textit{In breve:} Somma geometrica: $1+q+q^2+\dots+q^{n-1}=\frac{1-q^n}{1-q}$ per $q\ne1$. Se $|q|<1$, allora $q^n\to0$ e la somma tende a $\frac1{1-q}$.}


\subsection*{Disequazioni}
\begin{esercizio}[\liv{1}]
Risolvere: \ (a) $\dfrac{x^2-5x+6}{x-1}\le0$; \ (b) $|2x-3|<x$; \ (c) $\sqrt{x+3}>x-3$.
\end{esercizio}
\risp{(a) Numeratore $(x-2)(x-3)$, denominatore $x-1$; tabella dei segni: $x<1$ oppure $2\le x\le3$. (b) Serve $x>0$; poi $-x<2x-3<x$, cioè $x>1$ e $x<3$: $1<x<3$. (c) Dominio $x\ge-3$. Se $x<3$ il secondo membro è negativo e la disequazione è vera: $-3\le x<3$. Se $x\ge3$ si eleva al quadrato: $x+3>x^2-6x+9$, cioè $x^2-7x+6<0$, $1<x<6$: quindi $3\le x<6$. In tutto: $-3\le x<6$. \textit{In breve:} \emph{Fratte}: tutto a un membro, segno di numeratore e denominatore, tabella dei segni; il denominatore non si annulla. \emph{Modulo}: $|A|<B\iff-B<A<B$. \emph{Irrazionali}: $\sqrt A>B$ se $B<0$ e $A\ge0$, oppure se $B\ge0$ e $A>B^2$.}

\begin{esercizio}[\liv{2}]
Risolvere: \ (a) $4^x-3\cdot2^x-4>0$; \ (b) $\log_2(x-1)+\log_2(x+1)\le3$; \ (c) $(\log x)^2-\log x-2\ge0$.
\end{esercizio}
\risp{(a) Con $t=2^x>0$: $t^2-3t-4=(t-4)(t+1)>0$; poiché $t+1>0$, serve $t>4$, cioè $x>2$. (b) Dominio $x>1$; $\log_2(x^2-1)\le3$, cioè $x^2-1\le8$, $|x|\le3$: con il dominio, $1<x\le3$. (c) Con $t=\log x$ ($x>0$): $(t-2)(t+1)\ge0$, cioè $t\le-1$ oppure $t\ge2$: $0<x\le\frac1e$ oppure $x\ge e^2$. \textit{In breve:} \emph{Esponenziali}: con $t=a^x>0$ si passa a una disequazione algebrica; poi si torna a $x$ ricordando che $a^x$ è crescente se $a>1$. \emph{Logaritmi}: prima il dominio (argomenti $>0$), poi $\log_aA+\log_aB=\log_a(AB)$, e $\log_a$ crescente se $a>1$.}

\begin{esercizio}[\liv{2}]
Risolvere in $[0,2\pi]$: \ (a) $2\sin x-1\ge0$; \ (b) $2\cos^2x-\cos x-1<0$. \ Risolvere in $\bigl]-\frac\pi2,\frac\pi2\bigr[$: \ (c) $\tan x>1$.
\end{esercizio}
\risp{(a) $\sin x\ge\frac12$: $\frac\pi6\le x\le\frac{5\pi}6$. (b) Con $t=\cos x$: $(2t+1)(t-1)<0$, cioè $-\frac12<\cos x<1$: $0<x<\frac{2\pi}3$ oppure $\frac{4\pi}3<x<2\pi$ (si escludono $0$ e $2\pi$, dove $\cos x=1$). (c) $\frac\pi4<x<\frac\pi2$. \textit{In breve:} Si risolve l'equazione associata nell'intervallo dato e si legge il segno sulla circonferenza goniometrica (o sul grafico). Con $t=\sin x$ o $t=\cos x$ si passa a una disequazione algebrica, ricordando che $-1\le t\le1$. $\tan x$ è crescente in $\bigl]-\frac\pi2,\frac\pi2\bigr[$.}


\subsection*{Campi di esistenza}
\begin{esercizio}[\liv{1}]
Determinare il dominio di: \ (a) $\sqrt{\dfrac{x-1}{x+2}}$; \ (b) $\log(4-x^2)$; \ (c) $\dfrac1{e^x-1}$; \ (d) $\arcsin(2x-1)$.
\end{esercizio}
\risp{(a) $\frac{x-1}{x+2}\ge0$: $x<-2$ oppure $x\ge1$. (b) $4-x^2>0$: $-2<x<2$. (c) $e^x\ne1$: $x\ne0$. (d) $-1\le2x-1\le1$: $0\le x\le1$. \textit{In breve:} Le condizioni di esistenza, da mettere a sistema: radice di indice pari, radicando $\ge0$; logaritmo, argomento $>0$; denominatore $\ne0$; $\arcsin$ e $\arccos$, argomento fra $-1$ e $1$.}

\begin{esercizio}[\liv{2}]
Determinare il dominio di: \ (a) $\sqrt{\log(x-2)}$; \ (b) $\log(\sin x)$ in $[0,2\pi]$; \ (c) $\sqrt{2-|x-1|}$; \ (d) $\dfrac1{\sqrt{x^2-3x}}+\log(5-x)$; \ (e) $\arccos(\log x)$.
\end{esercizio}
\risp{(a) $x>2$ e $\log(x-2)\ge0$, cioè $x-2\ge1$: $x\ge3$. (b) $\sin x>0$: $0<x<\pi$. (c) $|x-1|\le2$: $-1\le x\le3$. (d) $x^2-3x>0$ ($x<0$ oppure $x>3$) e $x<5$: $x<0$ oppure $3<x<5$. (e) $x>0$ e $-1\le\log x\le1$: $\frac1e\le x\le e$. \textit{In breve:} In una composizione ogni funzione impone le sue condizioni, anche su ciò che sta dentro: in $\sqrt{\log(x-2)}$ servono $x-2>0$ (logaritmo) e $\log(x-2)\ge0$ (radice). $\log t\ge0\iff t\ge1$; $\sin x>0$ in $]0,\pi[$.}


\subsection*{Limiti di successioni}
\begin{esercizio}[\liv{1}]
Calcolare: \ (a) $\dfrac{3n^2-n}{1-2n^2}$; \ (b) $\sqrt{n^2+n}-n$; \ (c) $\dfrac{2^n+n^3}{3^n}$; \ (d) $n\sin\dfrac2n$; \ (e) $\Bigl(1+\dfrac3n\Bigr)^n$.
\end{esercizio}
\risp{(a) $-\frac32$ (stesso grado: rapporto dei coefficienti direttivi). (b) Razionalizzando, $\frac n{\sqrt{n^2+n}+n}=\frac1{\sqrt{1+1/n}+1}\to\frac12$. (c) $\bigl(\frac23\bigr)^n+\frac{n^3}{3^n}\to0$ (gerarchia). (d) $2\cdot\frac{\sin(2/n)}{2/n}\to2$. (e) $\Bigl[\bigl(1+\frac1{n/3}\bigr)^{n/3}\Bigr]^3\to e^3$. \textit{In breve:} Quoziente di polinomi: si raccoglie la potenza più alta. $\infty-\infty$ con le radici: si razionalizza. Gerarchia: $n^\alpha\ll b^n$ ($b>1$). Limiti notevoli con $a_n\to0$: $\frac{\sin a_n}{a_n}\to1$; e $\bigl(1+\frac1{a_n}\bigr)^{a_n}\to e$ con $a_n\to\pm\infty$.}

\begin{esercizio}[\liv{2}]
Calcolare: \ (a) $\Bigl(\dfrac{n-1}{n+1}\Bigr)^n$; \ (b) $n\bigl(e^{1/n}-1\bigr)$; \ (c) $n^2\Bigl(1-\cos\dfrac1n\Bigr)$; \ (d) $\dfrac{(\log n)^2}{\sqrt n}$; \ (e) $\dfrac{n^2+2^n}{n!}$.
\end{esercizio}
\risp{(a) $\frac{n-1}{n+1}=1-\frac2{n+1}=1+\frac1{a_n}$ con $a_n=-\frac{n+1}2\to-\infty$; l'esponente è $n=a_n\cdot\frac{-2n}{n+1}$, con $\frac{-2n}{n+1}\to-2$: il limite è $e^{-2}$. (b) Con $a_n=\frac1n\to0$: $\frac{e^{a_n}-1}{a_n}\to1$. (c) $\frac{1-\cos a_n}{a_n^2}\to\frac12$, con $a_n=\frac1n$. (d) $0$ (gerarchia: ogni potenza di $n$ batte ogni potenza del logaritmo). (e) $0$: il fattoriale batte sia le potenze sia gli esponenziali. \textit{In breve:} La forma $1^\infty$ si riconduce a $\bigl(1+\frac1{a_n}\bigr)^{a_n}\to e$. Con $a_n\to0$: $\frac{e^{a_n}-1}{a_n}\to1$, $\frac{1-\cos a_n}{a_n^2}\to\frac12$. Gerarchia: $\log n\ll n^\alpha\ll b^n\ll n!$.}


\subsection*{Limiti di funzioni}
\begin{esercizio}[\liv{1}]
Calcolare: \ (a) $\lim\limits_{x\to2}\dfrac{x^2-4}{x^2-3x+2}$; \ (b) $\lim\limits_{x\to+\infty}\dfrac{2x^3+1}{x^3-x}$; \ (c) $\lim\limits_{x\to0}\dfrac{\sin3x}{2x}$; \ (d) $\lim\limits_{x\to0}\dfrac{e^{2x}-1}{\sin x}$; \ (e) $\lim\limits_{x\to0}\dfrac{\log(1+3x)}x$; \ (f) $\lim\limits_{x\to+\infty}x\,e^{-x}$.
\end{esercizio}
\risp{(a) $\frac{(x-2)(x+2)}{(x-2)(x-1)}=\frac{x+2}{x-1}\to4$. (b) $2$. (c) $\frac32\cdot\frac{\sin3x}{3x}\to\frac32$. (d) $\frac{e^{2x}-1}{2x}\cdot\frac x{\sin x}\cdot2\to2$. (e) $3\cdot\frac{\log(1+3x)}{3x}\to3$. (f) $\frac x{e^x}\to0$ (gerarchia). \textit{In breve:} $\frac00$ con polinomi: si scompone e si semplifica il fattore che si annulla. Limiti notevoli per $x\to0$: $\frac{\sin x}x$, $\frac{e^x-1}x$, $\frac{\log(1+x)}x\to1$; per $x\to+\infty$ l'esponenziale batte ogni potenza.}

\begin{esercizio}[\liv{2}]
Calcolare: \ (a) $\lim\limits_{x\to0}\dfrac{1-\cos2x}{x\tan x}$; \ (b) $\lim\limits_{x\to+\infty}\Bigl(1-\dfrac2x\Bigr)^x$; \ (c) $\lim\limits_{x\to0^+}x\log x$; \ (d) $\lim\limits_{x\to1}\dfrac{\log x}{x^2-1}$; \ (e) $\lim\limits_{x\to+\infty}x\bigl(\log(x+1)-\log x\bigr)$; \ (f) $\lim\limits_{x\to0}\dfrac{\sqrt{1+x}-1}{e^x-1}$.
\end{esercizio}
\risp{(a) $4\cdot\frac{1-\cos2x}{(2x)^2}\cdot\frac x{\tan x}\to4\cdot\frac12\cdot1=2$. (b) Con $t=-\frac x2\to-\infty$: $\bigl(1+\frac1t\bigr)^{-2t}\to e^{-2}$. (c) Con $t=\frac1x\to+\infty$: $-\frac{\log t}t\to0$. (d) Con $t=x-1\to0$: $\frac{\log(1+t)}t\cdot\frac1{x+1}\to1\cdot\frac12=\frac12$. (e) $x\log\bigl(1+\frac1x\bigr)$; con $t=\frac1x\to0^+$: $\frac{\log(1+t)}t\to1$. (f) $\frac{(1+x)^{1/2}-1}x\cdot\frac x{e^x-1}\to\frac12\cdot1=\frac12$. \textit{In breve:} \emph{Cambio di variabile}: si pone $t=g(x)$, si dice dove va $t$, e ci si riconduce a un limite notevole. Utili: $\bigl(1+\frac ax\bigr)^x\to e^a$ per $x\to\pm\infty$; $x\log x\to0$ per $x\to0^+$; $\frac{(1+x)^\alpha-1}x\to\alpha$ per $x\to0$.}

\begin{esercizio}[\liv{2}]
(a) Trovare tutti gli asintoti di $f(x)=\dfrac{x^2+1}{x-1}$. (b) Calcolare i limiti destro e sinistro di $e^{1/x}$ in $0$, e classificare il punto di discontinuità. (c) Per quali $a$ e $b$ la funzione che vale $ax+b$ per $x<0$ e $e^{2x}$ per $x\ge0$ è continua in $0$? Per quali è anche derivabile in $0$?
\end{esercizio}
\risp{(a) Verticale $x=1$ (limiti $\pm\infty$); obliquo: $m=\lim\frac{f(x)}x=1$, $q=\lim\bigl(f(x)-x\bigr)=\lim\frac{x+1}{x-1}=1$: $y=x+1$, da entrambe le parti. (b) Per $x\to0^+$, $\frac1x\to+\infty$ e $e^{1/x}\to+\infty$; per $x\to0^-$, $e^{1/x}\to0$: discontinuità di seconda specie. (c) Continuità: $b=e^0=1$. Derivabilità: la derivata sinistra è $a$, la destra è $2e^0=2$: serve $a=2$ (e $b=1$). \textit{In breve:} Asintoto verticale: un limite laterale infinito. Obliquo: $m=\lim\frac{f(x)}x$, $q=\lim\bigl(f(x)-mx\bigr)$, entrambi finiti. Seconda specie: un limite laterale infinito o inesistente. Continua in $x_0$: limiti laterali uguali a $f(x_0)$; derivabile: derivate laterali uguali (e continuità).}


\subsection*{Calcolo di derivate}
\begin{esercizio}[\liv{1}]
Derivare: \ (a) $x^3-4x+2$; \ (b) $x^2\sin x$; \ (c) $\dfrac{x-1}{x+1}$; \ (d) $e^{3x}$; \ (e) $\log(x^2+1)$; \ (f) $\sqrt{4-x^2}$; \ (g) $\sin^2x$; \ (h) $x\,e^{-x}$.
\end{esercizio}
\risp{(a) $3x^2-4$; (b) $2x\sin x+x^2\cos x$; (c) $\frac{(x+1)-(x-1)}{(x+1)^2}=\frac2{(x+1)^2}$; (d) $3e^{3x}$; (e) $\frac{2x}{x^2+1}$; (f) $-\frac x{\sqrt{4-x^2}}$, per $|x|<2$; (g) $2\sin x\cos x$; (h) $e^{-x}-x\,e^{-x}=(1-x)\,e^{-x}$. \textit{In breve:} $(fg)'=f'g+fg'$; \ $\bigl(\frac fg\bigr)'=\frac{f'g-fg'}{g^2}$; \ composta: $\bigl(g(f(x))\bigr)'=g'\bigl(f(x)\bigr)\,f'(x)$. Tabella: $(e^x)'=e^x$, $(\log x)'=\frac1x$, $(\sqrt x)'=\frac1{2\sqrt x}$, $(\sin x)'=\cos x$, $(\cos x)'=-\sin x$.}

\begin{esercizio}[\liv{2}]
Derivare: \ (a) $e^{x^2}\cos x$; \ (b) $\log(\sin x)$; \ (c) $(x^2+1)^{10}$; \ (d) $\tan(x^2)$; \ (e) $\sqrt x\,\log x$; \ (f) $\log\dfrac{1+x}{1-x}$.
\end{esercizio}
\risp{(a) $e^{x^2}(2x\cos x-\sin x)$; (b) $\frac{\cos x}{\sin x}$, dove $\sin x>0$; (c) $20x\,(x^2+1)^9$; (d) $\frac{2x}{\cos^2(x^2)}=2x\bigl(1+\tan^2(x^2)\bigr)$ (attenzione: non è $\tan^2x$); (e) $\frac{\log x}{2\sqrt x}+\frac{\sqrt x}x=\frac{\log x+2}{2\sqrt x}$, per $x>0$; (f) scrivendo $\log(1+x)-\log(1-x)$: $\frac1{1+x}+\frac1{1-x}=\frac2{1-x^2}$, per $|x|<1$. \textit{In breve:} Più funzioni annidate: si deriva dall'esterno verso l'interno, moltiplicando le derivate. Prima di derivare conviene semplificare con le proprietà del logaritmo. $(\tan u)'=\frac{u'}{\cos^2u}=u'(1+\tan^2u)$.}

\begin{esercizio}[\liv{2}]
Scrivere l'equazione della retta tangente: \ (a) a $y=\log(1+x)$ in $x_0=0$; \ (b) a $y=x^2e^x$ in $x_0=1$; \ (c) a $y=e^{-x^2}$ in $x_0=1$.
\end{esercizio}
\risp{(a) $f(0)=0$, $f'(0)=1$: $y=x$. (b) $f(1)=e$, $f'(x)=(x^2+2x)\,e^x$, $f'(1)=3e$: $y=e+3e(x-1)$. (c) $f(1)=\frac1e$, $f'(x)=-2x\,e^{-x^2}$, $f'(1)=-\frac2e$: $y=\frac1e-\frac2e\,(x-1)$, cioè $y=\frac{3-2x}e$. \textit{In breve:} Retta tangente al grafico di $f$ in $x_0$: $y=f(x_0)+f'(x_0)(x-x_0)$. Servono due numeri: il valore $f(x_0)$ e la pendenza $f'(x_0)$.}

\stampasoluzioni
   % Lezioni 19--21 (L21: esercitazione)

% --- Schede di recupero: sempre in fondo, dopo le schede settimanali ---
% !TEX root =esercizi.tex
% ---------------------------------------------------------------------
% Scheda di recupero R1: argomenti della prima prova (11/11), con gli
% errori più frequenti emersi nella correzione. Esercizi nuovi: non sono
% presi dalla prova né dal fascicolo dell'esercitazione del 9/11.
% Tutti i risultati sono stati verificati con sympy.
% Le schede di recupero stanno in fondo a esercizi.tex, dopo le schede
% settimanali; la numerazione R1, R2, ... non interferisce con le altre.
% ---------------------------------------------------------------------
\renewcommand{\thechapter}{R1}
\chapter{Recupero della prima prova}
\renewcommand{\thesection}{\arabic{section}}

\begin{vuoto}[Come usare questa scheda]{}
La scheda riprende gli argomenti della prima prova (11 novembre) seguendo l'ordine dei quesiti. Ogni sezione si apre con un promemoria e con \textbf{l'errore da evitare}, cioè l'errore che nella correzione si è visto più spesso. Gli esercizi sono nuovi e vanno in ordine di difficoltà: \liv{1} di base, \liv{2} standard, \liv{3} più impegnativo.

Il consiglio è di svolgere ogni esercizio per iscritto, \emph{giustificando i passaggi}, e di confrontarvi con le soluzioni solo alla fine. Nella prova un risultato giusto ottenuto con un ragionamento sbagliato vale come esercizio sbagliato.
\end{vuoto}

%-------------------------------------------------------------------
\section{Calcolo combinatorio: conta l'ordine?}

\begin{vuoto}[Promemoria]{}
Si scelgono $k$ oggetti distinti fra $n$. Se \textbf{conta l'ordine} (primo, secondo, terzo; parole; podi; cariche diverse) si contano le disposizioni $D_{n,k}=n(n-1)\cdots(n-k+1)$. Se \textbf{non conta l'ordine} (gruppi, commissioni, campioni) si contano le combinazioni $\binom nk=\dfrac{D_{n,k}}{k!}$.
\end{vuoto}
\boxr{\emergencystretch=1.5em \textbf{Errore da evitare.} Usare $\binom nk$ quando l'ordine conta. Prima di calcolare chiedetevi: \emph{se scambio due degli oggetti scelti, ottengo un caso diverso?} Se sì, l'ordine conta.}

\begin{esercizio}[\liv{1}]
Dite se conta l'ordine e calcolate. \ (a) In un consiglio di $9$ persone si eleggono presidente, vicepresidente e segretario. \ (b) Fra le stesse $9$ persone si sceglie una commissione di $3$ membri. \ (c) Quante parole di $4$ lettere tutte diverse si formano con le lettere A, B, C, D, E, F? \ (d) Fra $11$ campioni d'acqua se ne scelgono $4$ da mandare in laboratorio.
\end{esercizio}
\risp{(a) Conta l'ordine (le cariche sono diverse): $D_{9,3}=9\cdot8\cdot7=504$. (b) Non conta: $\binom93=\frac{504}{3!}=84$. (c) Conta: $D_{6,4}=6\cdot5\cdot4\cdot3=360$. (d) Non conta: $\binom{11}4=330$.}

\begin{esercizio}[\liv{2}]
In un'area protetta ci sono $12$ stazioni di monitoraggio, $4$ delle quali vicino al fiume. Se ne ispezionano $5$. In quanti modi si possono scegliere: \ (a) in tutto; \ (b) con esattamente $2$ stazioni vicino al fiume; \ (c) con almeno una stazione vicino al fiume?
\end{esercizio}
\risp{(a) $\binom{12}5=792$. (b) Si scelgono $2$ stazioni fra le $4$ vicino al fiume e $3$ fra le altre $8$: $\binom42\binom83=6\cdot56=336$. (c) Conviene il complementare: le scelte \emph{senza} stazioni vicino al fiume sono $\binom85=56$, quindi $792-56=736$.}

\begin{esercizio}[\liv{2}]
Un'escursione tocca $3$ rifugi scelti fra $8$. (a) Quanti percorsi diversi ci sono, se conta l'ordine in cui si visitano? \ (b) Quante scelte dei rifugi ci sono, se conta solo quali si visitano? \ (c) Spiegate perché il primo numero è $6$ volte il secondo.
\end{esercizio}
\risp{(a) $D_{8,3}=336$. (b) $\binom83=56$. (c) Ogni scelta di $3$ rifugi si può percorrere in $3!=6$ ordini diversi: $336=56\cdot6$.}

%-------------------------------------------------------------------
\section{Progressioni geometriche}

\begin{vuoto}[Promemoria]{}
Se ogni termine è il precedente moltiplicato per $q$: $a_n=a_1\,q^{\,n-1}$, \quad $a_1+a_2+\dots+a_n=a_1\,\dfrac{q^n-1}{q-1}$ \ (con $q\ne1$).
\end{vuoto}
\boxr{\emergencystretch=1.5em \textbf{Errore da evitare.} Sbagliare l'esponente di una unità: il primo termine è $a_1=a_1q^0$, quindi l'$n$-esimo è $a_1q^{\,n-1}$. Controllate sempre la formula su $n=1$ e $n=2$.}

\begin{esercizio}[\liv{1}]
Calcolate la somma $3+6+12+\dots+384$. Quanti sono gli addendi?
\end{esercizio}
\risp{$q=2$ e $384=3\cdot2^7$, quindi l'ultimo termine è $a_8$: gli addendi sono $8$. Somma: $3\cdot\frac{2^8-1}{2-1}=3\cdot255=765$.}

\begin{esercizio}[\liv{2}]
Una coltura di batteri triplica ogni ora; alla fine della prima ora ci sono $600$ batteri. \ (a) Quanti batteri ci sono alla fine della sesta ora? \ (b) Alla fine di quale ora si superano per la prima volta $10^6$ batteri?
\end{esercizio}
\risp{$a_n=600\cdot3^{\,n-1}$. (a) $a_6=600\cdot3^5=145\,800$. (b) Serve $3^{\,n-1}>\frac{10^6}{600}\approx1666{,}7$, cioè $n-1>\frac{\log1666{,}7}{\log3}\approx6{,}76$: la prima ora intera è $n=8$ (controllo: $a_7=437\,400$, $a_8=1\,312\,200$).}

%-------------------------------------------------------------------
\section{Disequazioni esponenziali e logaritmiche: il verso}

\begin{vuoto}[Promemoria]{}
Se la base è $a>1$, le funzioni $a^x$ e $\log_a x$ sono crescenti e \textbf{il verso si conserva}. Se la base è $0<a<1$, sono decrescenti e \textbf{il verso si inverte}. Per i logaritmi, prima di tutto le condizioni di esistenza: gli argomenti devono essere positivi.
\end{vuoto}
\boxr{\emergencystretch=1.5em \textbf{Errore da evitare.} Con base minore di $1$ si dimentica di invertire il verso. Lo stesso errore si fa dividendo per un logaritmo negativo, per esempio $\log0{,}9<0$.}

\begin{esercizio}[\liv{1}]
Risolvete: \ (a) $\Bigl(\dfrac13\Bigr)^{x-1}>9$; \qquad (b) $0{,}5^{\,2x}\le0{,}125$; \qquad (c) $2^{x^2}<16$.
\end{esercizio}
\risp{(a) $9=\bigl(\frac13\bigr)^{-2}$ e la base è minore di $1$: $x-1<-2$, cioè $x<-1$. (b) $0{,}125=0{,}5^3$, base minore di $1$: $2x\ge3$, cioè $x\ge\frac32$. (c) $16=2^4$, base maggiore di $1$: $x^2<4$, cioè $-2<x<2$.}

\begin{esercizio}[\liv{2}]
Risolvete: \ (a) $\log_{1/2}(x-3)\ge-2$; \qquad (b) $\log_3(2x+1)<2$;\\ (c) $\log x+\log(x-2)<\log3$.
\end{esercizio}
\risp{(a) Esistenza: $x>3$. Poiché $-2=\log_{1/2}4$ e la base è minore di $1$: $x-3\le4$. Soluzione: $3<x\le7$. (b) Esistenza: $x>-\frac12$; $2x+1<9$, cioè $x<4$. Soluzione: $-\frac12<x<4$. (c) Esistenza: $x>2$. La disequazione diventa $\log\bigl(x(x-2)\bigr)<\log3$, cioè $x^2-2x-3<0$, $-1<x<3$. Con l'esistenza: $2<x<3$.}

\begin{esercizio}[\liv{2}]
La massa di una sostanza radioattiva, in grammi, è $m(t)=80\cdot0{,}9^{\,t}$ dopo $t$ anni. Dopo quanti anni interi la massa scende per la prima volta sotto $20$ grammi? Attenzione al verso quando dividete per $\log0{,}9$.
\end{esercizio}
\risp{$0{,}9^{\,t}<\frac14$. Passando ai logaritmi: $t\log0{,}9<\log\frac14$. Poiché $\log0{,}9<0$, dividendo il verso si inverte: $t>\frac{\log(1/4)}{\log0{,}9}\approx13{,}16$. Quindi dopo $14$ anni.}

%-------------------------------------------------------------------
\section{Disequazioni fratte: niente denominatori «mandati via»}

\begin{vuoto}[Promemoria]{}
Si porta tutto a primo membro, si fa il denominatore comune e si studia il \textbf{segno} di $\dfrac{N(x)}{D(x)}$: segno di $N$, segno di $D$, tabella dei segni. Gli zeri di $D$ sono sempre esclusi; gli zeri di $N$ sono inclusi solo con $\le$ o $\ge$.
\end{vuoto}
\boxr{\emergencystretch=1.5em \textbf{Errore da evitare.} Moltiplicare entrambi i membri per il denominatore. Non se ne conosce il segno, quindi non si sa se il verso resta o si inverte. Altri due errori frequenti: studiare solo il segno dei denominatori, e includere un estremo in cui il denominatore si annulla.}

\begin{esercizio}[\liv{1}]
Risolvete: \ (a) $\dfrac{x+3}{x-2}\ge2$; \qquad (b) $\dfrac3x<\dfrac1{x+4}$.
\end{esercizio}
\risp{(a) $\frac{x+3-2(x-2)}{x-2}=\frac{7-x}{x-2}\ge0$: il numeratore è positivo o nullo per $x\le7$, il denominatore è positivo per $x>2$. Soluzione: $2<x\le7$. (b) $\frac{3(x+4)-x}{x(x+4)}=\frac{2(x+6)}{x(x+4)}<0$. Gli zeri sono $-6$, $-4$, $0$; dalla tabella dei segni: $x<-6$ oppure $-4<x<0$.}

\begin{esercizio}[\liv{2}]
Uno studente risolve $\dfrac1{x-1}\ge2$ così: «moltiplico per $x-1$: $1\ge2x-2$, quindi $x\le\frac32$». Mostrate con un valore di $x$ che la conclusione è sbagliata, trovate l'errore e risolvete correttamente.
\end{esercizio}
\risp{Per $x=0$ la conclusione dello studente è soddisfatta, ma $\frac1{0-1}=-1<2$. L'errore: per $x<1$ il fattore $x-1$ è negativo e il verso si inverte. Correttamente: $\frac{1-2(x-1)}{x-1}=\frac{3-2x}{x-1}\ge0$, quindi $1<x\le\frac32$.}

\begin{esercizio}[\liv{2}]
Risolvete $\dfrac{x^2-4}{x^2-6x+5}\le0$.
\end{esercizio}
\risp{$\frac{(x-2)(x+2)}{(x-1)(x-5)}\le0$. Zeri del numeratore: $\pm2$, inclusi; zeri del denominatore: $1$ e $5$, esclusi. Dalla tabella dei segni: $-2\le x<1$ oppure $2\le x<5$.}

%-------------------------------------------------------------------
\section{Disequazioni irrazionali: i due sistemi}

\begin{vuoto}[Promemoria]{}
\[
\sqrt{A(x)}<B(x)\iff\begin{cases}A(x)\ge0\\ B(x)>0\\ A(x)<B(x)^2\end{cases}
\]
\[
\sqrt{A(x)}>B(x)\iff\begin{cases}B(x)<0\\ A(x)\ge0\end{cases}\ \text{oppure}\ \begin{cases}B(x)\ge0\\ A(x)>B(x)^2\end{cases}
\]
\end{vuoto}
\boxr{\emergencystretch=1.5em \textbf{Errore da evitare.} Elevare al quadrato senza condizioni, oppure, nel caso $\sqrt A>B$, dimenticare il primo sistema: quando $B<0$ la disequazione è vera ovunque la radice esiste.}

\begin{esercizio}[\liv{1}]
Risolvete: \ (a) $\sqrt{x+5}<x-1$; \qquad (b) $\sqrt{2x+6}\ge x+1$.
\end{esercizio}
\risp{(a) $x\ge-5$, $x>1$, $x+5<x^2-2x+1$, cioè $x^2-3x-4>0$: $x<-1$ oppure $x>4$. Mettendo insieme: $x>4$. (b) Primo sistema: $x+1<0$ e $2x+6\ge0$, cioè $-3\le x<-1$. Secondo sistema: $x\ge-1$ e $2x+6\ge x^2+2x+1$, cioè $x^2\le5$: $-1\le x\le\sqrt5$. Unione: $-3\le x\le\sqrt5$.}

\begin{esercizio}[\liv{2}]
Risolvete $\sqrt{x^2-16}>x-2$.
\end{esercizio}
\risp{Primo sistema: $x-2<0$ e $x^2\ge16$, cioè $x\le-4$. Secondo sistema: $x\ge2$ e $x^2-16>x^2-4x+4$, cioè $x>5$. Soluzione: $x\le-4$ oppure $x>5$. Chi dimentica il primo sistema perde tutta la semiretta $x\le-4$.}

%-------------------------------------------------------------------
\section{Disequazioni trigonometriche e domini}

\begin{vuoto}[Promemoria]{}
$\sin^2x\ge a$ (con $0\le a\le1$) equivale a $|\sin x|\ge\sqrt a$, cioè $\sin x\ge\sqrt a$ \emph{oppure} $\sin x\le-\sqrt a$; analogamente per il coseno. Per i domini si scrivono \emph{tutte} le condizioni, una per ogni «pezzo» della funzione, e poi si mettono a sistema.
\end{vuoto}
\boxr{\emergencystretch=1.5em \textbf{Errori da evitare.} (1) Da $\sin^2x\ge a$ passare a $\sin x\ge\sqrt a$, perdendo i valori negativi. (2) Nei domini con $\arcsin$ o $\arccos$, imporre la condizione $-1\le\cdot\le1$ alla somma intera invece che al solo argomento.}

\begin{esercizio}[\liv{1}]
Risolvete $2\sin^2x-1\le0$ nell'intervallo $[0,2\pi]$.
\end{esercizio}
\risp{$\sin^2x\le\frac12$, cioè $|\sin x|\le\frac{\sqrt2}2$: $-\frac{\sqrt2}2\le\sin x\le\frac{\sqrt2}2$. In $[0,2\pi]$: $0\le x\le\frac\pi4$, \ $\frac{3\pi}4\le x\le\frac{5\pi}4$, \ $\frac{7\pi}4\le x\le2\pi$.}

\begin{esercizio}[\liv{2}]
Risolvete $\cos^2x-\cos x\ge0$ nell'intervallo $[0,2\pi]$. Attenzione ai punti isolati.
\end{esercizio}
\risp{$\cos x\,(\cos x-1)\ge0$. Poiché $\cos x-1\le0$ sempre, serve $\cos x\le0$ oppure $\cos x=1$. Soluzione: $\frac\pi2\le x\le\frac{3\pi}2$, più i punti isolati $x=0$ e $x=2\pi$.}

\begin{esercizio}[\liv{2}]
Determinate il dominio di:\\ (a) $f(x)=\log\dfrac{x-1}{4-x}+\sqrt x$; \qquad (b) $g(x)=\arcsin\dfrac{x+1}3+\dfrac1{x^2-4}$.
\end{esercizio}
\risp{(a) $\frac{x-1}{4-x}>0$, cioè $1<x<4$, e $x\ge0$: dominio $1<x<4$. (b) Per l'arcoseno, solo sul suo argomento: $-1\le\frac{x+1}3\le1$, cioè $-4\le x\le2$. Per la frazione: $x\ne\pm2$. Dominio: $-4\le x<-2$ oppure $-2<x<2$.}

%-------------------------------------------------------------------
\section{Limiti di successioni e la forma $1^\infty$}

\begin{vuoto}[Promemoria]{}
$1^\infty$ è una \textbf{forma indeterminata}. Se $a_n\to0$ e $b_n\to+\infty$, allora $(1+a_n)^{b_n}=\Bigl[(1+a_n)^{1/a_n}\Bigr]^{a_nb_n}$ e la quantità fra parentesi quadre tende a $e$: \textbf{decide il limite di $a_nb_n$}. Se $a_nb_n\to\ell$ il limite è $e^\ell$; se $a_nb_n\to+\infty$ è $+\infty$; se $a_nb_n\to0$ è $1$.
\end{vuoto}
\boxr{\emergencystretch=1.5em \textbf{Errore da evitare.} Scrivere «$1^\infty=1$» oppure «$1^\infty=\infty$». Anche $\infty\cdot0$ è una forma indeterminata, e non vale automaticamente $0$.}

\begin{esercizio}[\liv{1}]
Calcolate i limiti per $n\to+\infty$ e osservate che sono tutti della forma $1^\infty$: \ (a) $\Bigl(1-\dfrac1{2n}\Bigr)^n$; \ (b) $\Bigl(1+\dfrac1{n^3}\Bigr)^{n^2}$; \ (c) $\Bigl(1+\dfrac1{\sqrt n}\Bigr)^n$.
\end{esercizio}
\risp{(a) $a_nb_n=-\frac1{2n}\cdot n=-\frac12$: limite $e^{-1/2}=\frac1{\sqrt e}$. (b) $a_nb_n=\frac{n^2}{n^3}=\frac1n\to0$: limite $e^0=1$. (c) $a_nb_n=\frac n{\sqrt n}=\sqrt n\to+\infty$: limite $+\infty$. La stessa forma $1^\infty$ dà tre risultati diversi.}

\begin{esercizio}[\liv{2}]
Calcolate $\displaystyle\lim_{n\to+\infty}\Bigl(\frac{n+3}{n+1}\Bigr)^{2n}$.
\end{esercizio}
\risp{$\frac{n+3}{n+1}=1+\frac2{n+1}$, quindi $a_n=\frac2{n+1}$, $b_n=2n$ e $a_nb_n=\frac{4n}{n+1}\to4$. Limite: $e^4$.}

\begin{esercizio}[\liv{2}]
Calcolate: \ (a) $\displaystyle\lim_{n\to+\infty}n\log\Bigl(1+\frac2n\Bigr)$; \qquad (b) $\displaystyle\lim_{n\to+\infty}\sqrt n\,\bigl(e^{1/n}-1\bigr)$. Entrambi sono della forma $\infty\cdot0$: i risultati sono uguali?
\end{esercizio}
\risp{(a) $n\log\bigl(1+\frac2n\bigr)=2\cdot\frac{\log(1+2/n)}{2/n}\to2\cdot1=2$. (b) $\sqrt n\,(e^{1/n}-1)=\frac{e^{1/n}-1}{1/n}\cdot\frac1{\sqrt n}\to1\cdot0=0$. Sono diversi: $\infty\cdot0$ è indeterminata.}

%-------------------------------------------------------------------
\section{Limiti notevoli}

\begin{vuoto}[Promemoria]{}
Se $t(x)\to0$: \quad $\dfrac{\sin t}t\to1$, \quad $\dfrac{1-\cos t}{t^2}\to\dfrac12$, \quad $\dfrac{e^t-1}t\to1$, \quad $\dfrac{\log(1+t)}t\to1$,\\[2pt]
$\dfrac{(1+t)^\alpha-1}t\to\alpha$.\\[2pt]
L'argomento $t$ può essere \emph{qualsiasi} espressione che tende a $0$, per esempio $5x$, $x^2$ o $\frac1x$ per $x\to+\infty$. Al denominatore deve comparire la \emph{stessa} $t$: se manca, la si fa comparire moltiplicando e dividendo.
\end{vuoto}
\boxr{\emergencystretch=1.5em \textbf{Errore da evitare.} Sostituire direttamente il valore limite in una forma $\frac00$ o $\infty\cdot0$ e concludere «$=0$». Una forma indeterminata va trasformata, non «calcolata».}

\begin{esercizio}[\liv{1}]
Calcolate per $x\to0$: \ (a) $\dfrac{\sin5x}{2x}$; \qquad (b) $\dfrac{e^{3x}-1}{\sin x}$; \qquad (c) $\dfrac{\log(1+x^2)}{1-\cos x}$.
\end{esercizio}
\risp{(a) $\frac52\cdot\frac{\sin5x}{5x}\to\frac52$. (b) $3\cdot\frac{e^{3x}-1}{3x}\cdot\frac x{\sin x}\to3$. (c) $\frac{\log(1+x^2)}{x^2}\cdot\frac{x^2}{1-\cos x}\to1\cdot2=2$.}

\begin{esercizio}[\liv{2}]
Calcolate: \ (a) $\displaystyle\lim_{x\to0}\frac{\sqrt[3]{1+x}-1}{\sin2x}$; \quad (b) $\displaystyle\lim_{x\to+\infty}x\sin\frac2x$; \quad (c) $\displaystyle\lim_{x\to0}\frac{1-\cos x}{x\sin2x}$.
\end{esercizio}
\risp{(a) $\frac{(1+x)^{1/3}-1}x\cdot\frac{2x}{\sin2x}\cdot\frac12\to\frac13\cdot1\cdot\frac12=\frac16$. (b) Con $t=\frac2x\to0$: $x\sin\frac2x=2\cdot\frac{\sin t}t\to2$. (c) $\frac{1-\cos x}{x^2}\cdot\frac{2x}{\sin2x}\cdot\frac12\to\frac12\cdot1\cdot\frac12=\frac14$.}

\begin{esercizio}[\liv{2}]
Trovate l'errore e correggete. \ (a) «$\displaystyle\lim_{x\to0}\frac{e^{2x}-1}x=\frac{e^0-1}0=\frac00=0$.» \quad (b) «$\displaystyle\lim_{x\to+\infty}x\sin\frac1x=\infty\cdot\sin0=\infty\cdot0=0$.»
\end{esercizio}
\risp{In entrambi i casi si è «calcolata» una forma indeterminata. (a) $2\cdot\frac{e^{2x}-1}{2x}\to2$. (b) Con $t=\frac1x\to0^+$: $\frac{\sin t}t\to1$.}

%-------------------------------------------------------------------
\section{Retta tangente con una condizione}

\begin{vuoto}[Promemoria]{}
La retta tangente al grafico di $f$ nel punto di ascissa $x_0$ è $y=f(x_0)+f'(x_0)(x-x_0)$. Se il punto non è dato, ma è data una condizione sulla pendenza (per esempio «parallela a $y=mx+q$», oppure «orizzontale», cioè $m=0$), si procede così:
\begin{enumerate}
\item si calcola $f'(x)$;
\item si risolve $f'(x_0)=m$, trovando $x_0$;
\item si calcola $f(x_0)$ e si scrive la retta.
\end{enumerate}
\end{vuoto}
\boxr{\emergencystretch=1.5em \textbf{Errore da evitare.} Mettere a sistema la curva con la retta data. La retta data serve solo per la pendenza: la tangente cercata le è parallela e quindi non la incontra. Senza la derivata l'esercizio non si può impostare.}

\begin{esercizio}[\liv{1}]
Scrivete la retta tangente al grafico di $f(x)=x^2-3x$ nel punto di ascissa $x_0=2$.
\end{esercizio}
\risp{$f(2)=-2$, $f'(x)=2x-3$, $f'(2)=1$: $y=-2+(x-2)$, cioè $y=x-4$.}

\begin{esercizio}[\liv{2}]
Scrivete la retta tangente al grafico di $f(x)=e^{2x}$ parallela alla retta $y=4x+1$.
\end{esercizio}
\risp{$f'(x)=2e^{2x}=4$ dà $e^{2x_0}=2$, cioè $x_0=\frac{\log2}2$, e $f(x_0)=2$. Tangente: $y=2+4\bigl(x-\frac{\log2}2\bigr)$, cioè $y=4x+2-2\log2$.}

\begin{esercizio}[\liv{2}]
Scrivete la retta tangente al grafico di $f(x)=\sqrt{x+1}$ parallela alla retta $y=\dfrac x6$.
\end{esercizio}
\risp{$f'(x)=\frac1{2\sqrt{x+1}}=\frac16$ dà $\sqrt{x_0+1}=3$, cioè $x_0=8$, e $f(8)=3$. Tangente: $y=3+\frac16(x-8)$, cioè $y=\frac x6+\frac53$.}

\begin{esercizio}[\liv{2}]
In quale punto il grafico di $f(x)=x\,e^{-x}$ ha la retta tangente orizzontale? Scrivetene l'equazione.
\end{esercizio}
\risp{$f'(x)=e^{-x}-xe^{-x}=(1-x)e^{-x}=0$ per $x_0=1$; $f(1)=\frac1e$. Tangente orizzontale: $y=\frac1e$, nel punto $\bigl(1,\frac1e\bigr)$.}

%-------------------------------------------------------------------
\section{Asintoti}

\begin{vuoto}[Promemoria: la procedura completa]{}
\begin{enumerate}
\item Dominio.
\item Limiti ai bordi del dominio, \textbf{destro e sinistro} nei punti esclusi. Un limite infinito dà un asintoto verticale.
\item Limiti per $x\to+\infty$ e per $x\to-\infty$, \textbf{separatamente}. Un limite finito $\ell$ dà l'asintoto orizzontale $y=\ell$.
\item Se il limite è infinito, si cerca l'asintoto obliquo $y=mx+q$, con $m=\lim\frac{f(x)}x$ (finito e non nullo) e $q=\lim\bigl(f(x)-mx\bigr)$ (finito).
\end{enumerate}
\end{vuoto}
\boxr{\emergencystretch=1.5em \textbf{Errori da evitare.} Fermarsi a $x\to+\infty$: a $-\infty$ l'asintoto può essere diverso. Calcolare un solo limite laterale. Dichiarare un asintoto verticale in un punto dove il limite è finito.}

\begin{esercizio}[\liv{1}]
Trovate tutti gli asintoti di $f(x)=\dfrac{2x^2+1}{x-1}$.
\end{esercizio}
\risp{Dominio: $x\ne1$. $\lim_{x\to1^\pm}f(x)=\pm\infty$: asintoto verticale $x=1$. Per $x\to\pm\infty$ il limite è $\pm\infty$; $m=\lim\frac{2x^2+1}{x^2-x}=2$, $q=\lim\Bigl(\frac{2x^2+1}{x-1}-2x\Bigr)=\lim\frac{2x+1}{x-1}=2$. Asintoto obliquo $y=2x+2$ sia a $+\infty$ sia a $-\infty$.}

\begin{esercizio}[\liv{2}]
Sia $f(x)=\dfrac{\log x}{x-1}$. Trovate il dominio e gli asintoti, e dite che tipo di discontinuità c'è in $x=1$.
\end{esercizio}
\risp{Dominio: $x>0$, $x\ne1$. Per $x\to0^+$: $\log x\to-\infty$ e $x-1\to-1$, quindi $f\to+\infty$: asintoto verticale $x=0$ (da destra). Per $x\to1$: con $t=x-1\to0$, $\frac{\log(1+t)}t\to1$. Il limite è finito, quindi in $x=1$ c'è una discontinuità eliminabile e \emph{nessun} asintoto. Per $x\to+\infty$: $f\to0$, asintoto orizzontale $y=0$.}

\begin{esercizio}[\liv{3}]
Trovate gli asintoti di $f(x)=\sqrt{x^2+4x}$, distinguendo $x\to+\infty$ e $x\to-\infty$.
\end{esercizio}
\risp{Dominio: $x\le-4$ oppure $x\ge0$; non ci sono asintoti verticali. Ricordate che $\sqrt{x^2}=|x|$. Per $x\to+\infty$: $m=\lim\frac{\sqrt{x^2+4x}}x=1$, $q=\lim\bigl(\sqrt{x^2+4x}-x\bigr)=\lim\frac{4x}{\sqrt{x^2+4x}+x}=2$: asintoto $y=x+2$. Per $x\to-\infty$: $\sqrt{x^2+4x}=|x|\sqrt{1+4/x}=-x\sqrt{1+4/x}$, quindi $m=-1$, e $q=\lim\bigl(\sqrt{x^2+4x}+x\bigr)=\lim\frac{4x}{\sqrt{x^2+4x}-x}=-2$: asintoto $y=-x-2$.}

%-------------------------------------------------------------------
\section{Continuità e derivabilità con parametri}

\begin{vuoto}[Promemoria]{}
Sia $f(x)=g(x)$ per $x\le x_0$ e $f(x)=h(x)$ per $x>x_0$, con $g$ e $h$ derivabili.
\begin{itemize}
\item \textbf{Continuità} in $x_0$: $g(x_0)=\lim_{x\to x_0^+}h(x)$.
\item \textbf{Derivabilità} in $x_0$: \emph{serve la continuità}, e in più $g'(x_0)=\lim_{x\to x_0^+}h'(x)$, cioè derivata sinistra uguale a derivata destra.
\end{itemize}
Si ottengono due equazioni nei due parametri.
\end{vuoto}
\boxr{\emergencystretch=1.5em \textbf{Errore da evitare.} Imporre solo l'uguaglianza delle derivate, dimenticando la continuità: una funzione discontinua non è mai derivabile.}

\begin{esercizio}[\liv{1}]
Sia $f(x)=ax+b$ per $x<0$ e $f(x)=e^{2x}$ per $x\ge0$. Per quali $a$, $b$ la funzione $f$ è continua in $0$? E derivabile?
\end{esercizio}
\risp{Continuità: $\lim_{x\to0^-}(ax+b)=b$ deve valere $e^0=1$, quindi $b=1$ (con $a$ qualsiasi). Derivabilità: inoltre $a=(e^{2x})'\big|_{x=0}=2$. Quindi $a=2$, $b=1$.}

\begin{esercizio}[\liv{2}]
Sia $f(x)=x^2+a$ per $x\le1$ e $f(x)=b\log x+3$ per $x>1$. Trovate $a$ e $b$ in modo che $f$ sia derivabile in $x=1$.
\end{esercizio}
\risp{Continuità: $1+a=b\log1+3=3$, quindi $a=2$. Derivate: a sinistra $2x\big|_{x=1}=2$, a destra $\frac bx\big|_{x=1}=b$. Quindi $b=2$.}

\begin{esercizio}[\liv{2}]
Sia $f(x)=x^2$ per $x\le1$ e $f(x)=ax$ per $x>1$. Esiste un valore di $a$ per cui $f$ è derivabile in $x=1$?
\end{esercizio}
\risp{La continuità richiede $1=a$. Con $a=1$ la derivata sinistra è $2$ e quella destra è $1$: sono diverse. Quindi non esiste nessun valore di $a$: per $a=1$ la funzione è continua ma non derivabile, per $a\ne1$ non è nemmeno continua.}

\begin{esercizio}[\liv{3}]
Sia $f(x)=a\sin x+b$ per $x\le0$ e $f(x)=\sqrt{x+4}$ per $x>0$. (a) Per quali $a$, $b$ la funzione $f$ è continua in $0$? \ (b) Per quali è derivabile? \ (c) Per quali è continua ma non derivabile?
\end{esercizio}
\risp{(a) $f(0)=b$ e $\lim_{x\to0^+}\sqrt{x+4}=2$: $b=2$, con $a$ qualsiasi. (b) Derivata sinistra $a\cos0=a$, derivata destra $\frac1{2\sqrt{x+4}}\big|_{x=0}=\frac14$: $a=\frac14$, $b=2$. (c) $b=2$ e $a\ne\frac14$.}

\stampasoluzioni
   % R1: recupero della prima prova (11/11)

\end{document}
